%TODO: semua yg ada di judul wajib ada subbabnya di bab 2. termasiuk multimodal dan deteksi objek

\chapter{TINJAUAN PUSTAKA}
\section{Penyakit Tanaman}
Penyakit tanaman secara fundamental didefinisikan sebagai gangguan atau penyimpangan dari kondisi fisiologis normal yang disebabkan oleh interaksi berkelanjutan dengan faktor eksternal, yang bermanifestasi pada perubahan morfologi, fisiologi, maupun perilaku tanaman \parencite{vishnoiPlantDiseaseDetection2021}. \textcite{muhammadComprehensiveReviewCrop2025} memperluas definisi ini dengan terminologi cekaman (\textit{crop stress}), yakni segala kondisi eksternal yang memberikan dampak negatif terhadap metabolisme, pertumbuhan, dan produktivitas tanaman. Segala bentuk gangguan yang memengaruhi vitalitas tanaman tersebut menuntut adanya pemahaman ilmiah yang komprehensif.

Disiplin ilmu yang mempelajari fenomena ini adalah Fitopatologi (\textit{Phytopathology}). Berdasarkan etimologinya, istilah ini berasal dari bahasa Yunani \textit{phyton} (tanaman), \textit{pathos} (penyakit), dan \textit{logos} (pengetahuan) \parencite{bhargavaPlantLeafDisease2024}. Bhargava dkk. menjelaskan bahwa Fitopatologi bukan sekadar ilmu pengenalan penyakit, melainkan sebuah sistem terpadu yang mencakup empat tujuan utama, yaitu etiologi yang berfokus pada studi sumber atau penyebab penyakit, patogenesis yang mempelajari mekanisme perkembangan penyakit, epidemiologi yang mengkaji interaksi antara populasi inang, patogen, dan lingkungan, serta manajemen yang merancang strategi pengendalian untuk memitigasi kerugian. Fitopatologi mengintegrasikan berbagai subdisiplin ilmu seperti bakteriologi, virologi, mikologi, hingga biologi molekuler dan meteorologi untuk memahami dinamika kesehatan tanaman secara holistik.

Identifikasi penyebab cekaman pada tanaman secara mendasar dimulai dengan mengklasifikasikannya ke dalam dua kategori utama, yaitu faktor abiotik dan faktor biotik \parencite{bhargavaPlantLeafDisease2024, muhammadComprehensiveReviewCrop2025, georgePresentFutureDeep2025}. Visualisasi klasifikasi cekaman ini disajikan pada Gambar \ref{fig:taxonomy_plant_disease}. Cekaman abiotik dipicu oleh faktor lingkungan non-hidup seperti suhu ekstrem, salinitas tanah, dan ketidakseimbangan nutrisi yang bertindak sebagai stresor fisik pengganggu homeostasis tanaman \parencite{muhammadComprehensiveReviewCrop2025}. Secara karakteristik, cekaman abiotik bersifat tidak menular dengan sebaran di lahan yang cenderung seragam mengikuti kondisi fisik lingkungan, serta gejala yang muncul secara serentak pada banyak individu. Diagnosis untuk kategori ini umumnya dilakukan melalui analisis tanah, cuaca, dan jaringan tanaman. Sebaliknya, cekaman biotik disebabkan oleh organisme hidup parasit seperti jamur (\textit{fungi}), bakteri, virus, dan nematoda, maupun serangan hama yang dapat berperan sebagai vektor pembawa penyakit. Berbeda dengan faktor abiotik, cekaman biotik memiliki sifat menular yang menyebar secara bertahap seiring siklus hidup patogen, dengan pola sebaran di lahan yang cenderung berkelompok (\textit{patchy}) atau menyebar secara radial. Identifikasi cekaman biotik memerlukan pengamatan visual terhadap tanda, kultur, atau teknik molekuler.
\begin{figure}[ht]
\centering
\includegraphics[width=\textwidth]{bab2/abioticvsbiotic.jpg}
\caption[Klasifikasi cekaman biotik dan abiotik pada tanaman]{Klasifikasi cekaman biotik dan abiotik pada tanaman \parencite{georgePresentFutureDeep2025}}
\label{fig:taxonomy_plant_disease}
\end{figure}

% Terjadinya penyakit biotik yang sukses memerlukan interaksi kompleks yang dikenal sebagai Segitiga Penyakit (\textit{Disease Triangle}). Konsep ini menjelaskan bahwa infeksi hanya akan terjadi jika terdapat pertemuan antara inang yang rentan (\textit{susceptible host}), patogen yang virulen (\textit{virulent pathogen}), dan kondisi lingkungan yang mendukung (\textit{favorable environment}) (Bhargava dkk., 2024). Ketiadaan salah satu komponen tersebut akan menggagalkan proses patogenesis.

Dalam proses identifikasi, membedakan antara gejala dan tanda merupakan langkah fundamental bagi identifikasi yang akurat. Gejala merujuk pada respon internal atau perubahan fisiologis tanaman akibat gangguan, seperti klorosis (menguning), nekrosis (kematian jaringan), atau distorsi bentuk \parencite{muhammadComprehensiveReviewCrop2025}. Sementara itu, tanda (\textit{signs}) adalah bukti fisik keberadaan patogen, seperti kumpulan benang jamur (miselium), lendir bakteri (\textit{ooze}), atau keberadaan hama itu sendiri \parencite{bhargavaPlantLeafDisease2024}. Pengenalan visual terhadap indikator-indikator ini, baik yang bersifat makroskopis maupun mikroskopis, menjadi basis utama pengembangan sistem deteksi otomatis berbasis kecerdasan buatan. 

\section{\textit{Artificial Intelligence}}
\textit{Artificial Intelligence} (AI) atau kecerdasan buatan merupakan cabang ilmu komputer yang berfokus pada pengembangan sistem yang mampu melakukan tugas-tugas yang biasanya memerlukan kecerdasan manusia, seperti pengenalan pola, penalaran, pembelajaran, dan pengambilan keputusan \parencite{russell2020artificial}. Dalam perkembangannya, Russell dan Norvig mengategorikan AI ke dalam empat dimensi utama, yakni sistem yang berpikir seperti manusia, sistem yang bertindak seperti manusia, sistem yang berpikir secara rasional, dan sistem yang bertindak secara rasional. Dimensi terakhir ini yang menjadi landasan bagi pengembangan sistem cerdas yang mampu memecahkan masalah kompleks secara otonom.

Berdasarkan mekanisme operasionalnya, sistem yang menerapkan AI dapat diklasifikasikan menjadi dua paradigma utama, yakni \textit{rule-based system} (sistem berbasis aturan), dan \textit{machine learning system} (sistem pembelajaran mesin) \parencite{cholletDeepLearningPython2026}. \textit{Rule-based system} merepresentasikan paradigma pemrograman klasik, yang bekerja dengan mengikuti sekumpulan instruksi atau aturan yang didefinisikan secara manual oleh pemrogram untuk memproses data menjadi jawaban. Meskipun pendekatan ini sangat andal pada domain yang terstruktur, ia memiliki keterbatasan dalam menangani data yang kompleks dan tidak terduga. Sebaliknya, \textit{machine learning} memperkenalkan paradigma baru yang membalikkan proses tersebut. Sistem ini tidak diprogram secara eksplisit dengan aturan, melainkan dilatih untuk mempelajari aturan tersebut secara otomatis dari data. Perbandingan alur kerja antara kedua paradigma tersebut dapat digambarkan pada Gambar \ref{fig:comparison_rule_learning}. 

\begin{figure}[H] 
\centering
\begin{subfigure}[b]{0.8\textwidth}
    \centering
    \includegraphics[width=\textwidth]{bab2/rulebasedsystem.png}
    \caption{}
    \label{fig:rule_based}
\end{subfigure}
\vspace{0.5cm}
\begin{subfigure}[b]{0.8\textwidth}
    \centering
    \includegraphics[width=\textwidth]{bab2/learningbasedsystem.png}
    \caption{}
    \label{fig:learning_based}
\end{subfigure}
\caption[Perbandingan paradigma \textit{rule-based system} dan \textit{machine learning system}]{Perbandingan paradigma \textit{rule-based system} dan \textit{machine learning system}: (a) \textit{Rule-based system} memproses data berdasarkan aturan untuk menghasilkan jawaban, sedangkan (b) \textit{Machine learning system} mempelajari aturan dari data dan jawaban yang diberikan.}
\label{fig:comparison_rule_learning}
\end{figure}

Sistem \textit{machine learning} dilatih dengan dihadapkan pada banyak contoh yang relevan, di mana ia menemukan struktur statistik yang memungkinkannya membuat aturan untuk mengotomatisasi tugas \parencite{cholletDeepLearningPython2026}. Fleksibilitas ini menjadi kunci dalam aplikasi praktis seperti identifikasi penyakit tanaman, di mana karakteristik visual gejala seringkali terlalu kompleks untuk dipetakan secara manual ke dalam aturan logika yang kaku \parencite{pacalSystematicReviewDeep2024}. Untuk mengoptimalkan proses akuisisi pengetahuan tersebut, algoritma \textit{machine learning} dikembangkan ke dalam berbagai metode yang disesuaikan dengan karakteristik data yang tersedia.

Secara umum, algoritma \textit{machine learning} dikelompokkan ke dalam empat kategori utama berdasarkan jenis data dan metode pembelajarannya \parencite{sarkerMachineLearningAlgorithms2021}. Visualisasi dari kategori-kategori tersebut disajikan pada Gambar \ref{fig:taxonomy_ml}. Keempat kategori tersebut dijabarkan sebagai berikut:
\begin{enumerate}[nosep]
    \item \textit{Supervised Learning}: Pendekatan di mana model dilatih menggunakan \textit{dataset} yang memiliki label eksplisit, yang terdiri dari pasangan masukan dan luaran yang benar. Tujuannya adalah untuk memetakan fungsi dari masukan ke luaran sehingga model dapat melakukan prediksi atau klasifikasi pada data baru yang belum pernah dilihat sebelumnya.
    \item \textit{Unsupervised Learning}: Pendekatan ini bekerja pada data yang tidak memiliki label atau anotasi sebelumnya. Fokus utamanya adalah untuk menemukan struktur, pola, atau hubungan tersembunyi di dalam data tanpa adanya panduan luaran yang spesifik, seperti pada tugas klasterisasi atau asosiasi.
    \item \textit{Semi-supervised Learning}: Kategori ini merupakan hibrida yang memanfaatkan \textit{dataset} berlabel dalam jumlah kecil bersama dengan \textit{dataset} tidak berlabel dalam jumlah besar selama proses pelatihan. Pendekatan ini dapat digunakan dalam skenario di mana biaya anotasi data oleh pakar sangat mahal, namun data mentah tersedia secara melimpah.
    \item \textit{Reinforcement Learning}: Metode pembelajaran yang berbasis pada interaksi agen dengan lingkungan melalui mekanisme \textit{reward} (penghargaan) dan \textit{penalty} (hukuman). Agen belajar untuk mengambil serangkaian tindakan yang dapat memaksimalkan akumulasi \textit{reward} dalam jangka panjang melalui proses \textit{trial and error}.
\end{enumerate}


\begin{figure}[H]
\centering
\includegraphics[width=0.8\textwidth]{bab2/kategoriml.png}
\caption[Ragam teknik \textit{machine learning}]{Ragam teknik \textit{machine learning} \parencite{sarkerMachineLearningAlgorithms2021}}
\label{fig:taxonomy_ml}
\end{figure}

% [TODO: MENTION ABOUT SELF-SUPERVISED LEARNING AS A SUBSET OF SUPERVISED LEARNING]

\section{\textit{Deep Learning}}
\textit{Deep learning} merupakan subbidang dari \textit{machine learning} yang menggunakan arsitektur \textit{artificial neural networks} (jaringan saraf tiruan) dengan banyak lapisan tersembunyi untuk mempelajari representasi data secara hierarkis \parencite{cholletDeepLearningPython2026,goodfellowDeepLearning2016, russell2020artificial}. Secara taksonomi, \textit{deep learning} merupakan bagian integral dari \textit{machine learning}, yang mana \textit{machine learning} itu sendiri berada di bawah payung besar AI. Hubungan hierarkis ini dapat divisualisasikan pada Gambar \ref{fig:hierarchy_ai_ml_dl}.
%NOTE: COPYRIGHT RED
\begin{figure}[H] 
\centering
\includegraphics[width=0.45\textwidth]{bab2/hierarkiai.png}
\caption[Hierarki antara \textit{artificial intelligence}, \textit{machine learning}, \textit{representation learning}, dan \textit{deep learning}]{Hierarki antara \textit{artificial intelligence}, \textit{machine learning}, dan \textit{deep learning} \parencite{cholletDeepLearningPython2026}}
\label{fig:hierarchy_ai_ml_dl}
\end{figure}

\textit{Artificial neural networks} (ANN) merupakan fondasi utama dari \textit{deep learning} yang terinspirasi oleh struktur biologis otak manusia untuk memproses informasi melalui unit-unit komputasi yang saling terhubung \parencite{sarkerDeepLearningComprehensive2021}. Unit dasar dalam ANN adalah neuron yang menerima sinyal masukan, memberikan bobot dan bias tertentu, serta menerapkan fungsi aktivasi untuk menghasilkan luaran. Arsitektur dasar sebuah jaringan saraf terdiri dari tiga jenis lapisan utama, yaitu lapisan masukan, satu atau lebih lapisan tersembunyi, dan lapisan luaran. Secara spesifik, lapisan masukan berfungsi meneruskan data mentah ke dalam jaringan, sedangkan lapisan tersembunyi bertugas melakukan transformasi fitur melalui operasi bobot dan fungsi non-linear. Peningkatan jumlah lapisan tersembunyi akan memperdalam arsitektur jaringan, yang memungkinkan model untuk mempelajari hierarki representasi data yang lebih kompleks. Selanjutnya, hasil pemrosesan tersebut dikonversi oleh lapisan luaran menjadi nilai probabilitas atau prediksi akhir. Struktur jaringan saraf dan mekanisme operasional pada tingkat neuron diilustrasikan pada Gambar \ref{fig:nn_architecture}.

\begin{figure}[H] 
\centering
\begin{subfigure}[b]{0.8\textwidth}
    \centering
    \includegraphics[width=\textwidth]{bab2/ann.png}
    \caption{}
\end{subfigure}
\vspace{0.5cm}
\begin{subfigure}[b]{\textwidth}
    \centering
    \includegraphics[width=\textwidth]{bab2/neuron.png}
    \caption{}
\end{subfigure}
\caption[Arsitektur jaringan saraf tiruan dan mekanisme kerja neuron]{Arsitektur jaringan saraf tiruan: (a) Hierarki lapisan pada ANN dan (b) Mekanisme matematis pada tingkat neuron \parencite{guillodArtificialNeuralNetwork2020}}
\label{fig:nn_architecture}
\end{figure}

Mekanisme pembelajaran pada model \textit{deep learning} berfokus pada pencarian sekumpulan nilai \textit{weights} (bobot) yang mampu memetakan masukan ke target yang benar secara akurat \parencite{cholletDeepLearningPython2026}. Spesifikasi transformasi data pada setiap lapisan ditentukan oleh bobot-bobot ini. Selama proses pelatihan, kualitas prediksi model ($Y'$) diukur menggunakan \textit{loss function} (fungsi kerugian) yang membandingkannya dengan target sebenarnya ($Y$) untuk menghasilkan skor kerugian (\textit{loss score}). Skor ini bertindak sebagai sinyal umpan balik yang mengindikasikan seberapa jauh penyimpangan model dari kinerja ideal. Informasi tersebut kemudian digunakan oleh \textit{optimizer} untuk memodifikasi nilai-nilai bobot melalui algoritma \textit{Backpropagation}, dengan tujuan meminimalkan skor kerugian pada iterasi berikutnya. Siklus iteratif yang terdiri dari prediksi, pengukuran kesalahan, dan penyesuaian bobot ini dikenal sebagai \textit{training loop}, yang terus dijalankan hingga model mencapai performa yang diharapkan. Proses pembelajaran ini divisualisasikan pada Gambar \ref{fig:dl_learning_process}.
% #NOTE: COPYRIGHT RED
\begin{figure}[H]
\centering
\includegraphics[width=0.7\textwidth]{bab2/dllearning.png}
\caption[Mekanisme pembelajaran pada model \textit{deep learning}]{Mekanisme pembelajaran pada model \textit{deep learning} \parencite{cholletDeepLearningPython2026}}
\label{fig:dl_learning_process}
\end{figure} 

Dalam perkembangannya, teknik-teknik \textit{deep learning} dapat dikategorikan berdasarkan karakteristik pembelajaran dan aplikasinya. \textcite{sarkerDeepLearningComprehensive2021} mengajukan taksonomi yang membagi teknik \textit{deep learning} ke dalam tiga kategori utama, yaitu \textit{supervised} atau \textit{discriminative learning}, \textit{unsupervised} atau \textit{generative learning}, serta \textit{hybrid learning}. Pembelajaran diskriminatif berfokus pada penyediaan fungsi pembeda untuk tugas klasifikasi terawasi, dengan arsitektur populer seperti \textit{Multi-Layer Perceptron} (MLP), \textit{Convolutional Neural Networks} (CNN), dan \textit{Recurrent Neural Networks} (RNN). Di sisi lain, pembelajaran generatif digunakan untuk mengekstraksi fitur korelasi tingkat tinggi atau sintesis data tanpa label, yang mencakup metode seperti \textit{Generative Adversarial Networks} (GAN), \textit{Autoencoders} (AE), dan \textit{Deep Belief Networks} (DBN). Kategori terakhir, pembelajaran hibrida, mengintegrasikan kedua pendekatan tersebut untuk menggabungkan keunggulan model generatif dan diskriminatif dalam menyelesaikan masalah yang lebih kompleks. Visualisasi taksonomi teknik \textit{deep learning} ini disajikan pada Gambar \ref{fig:dl_technique}.

\begin{figure}[H]
\centering
\includegraphics[width=\textwidth]{bab2/dltechnique.png}
\caption[Taksonomi teknik \textit{deep learning}]{Taksonomi teknik \textit{deep learning} \parencite{sarkerDeepLearningComprehensive2021}}
\label{fig:dl_technique}
\end{figure}

Perbedaan utama antara \textit{machine learning} tradisional dan \textit{deep learning} terletak pada cara fitur data diperoleh \parencite{goodfellowDeepLearning2016}. Fitur data dapat didefinisikan sebagai atribut atau karakteristik spesifik yang diekstraksi dari data mentah untuk merepresentasikan informasi penting bagi model. Goodfellow dkk. menjelaskan bahwa \textit{machine learning} tradisional sangat bergantung pada proses \textit{feature extraction} (ekstraksi fitur) secara manual yang membutuhkan pengetahuan domain mendalam. Untuk mengatasi ketergantungan tersebut, dikembangkan pendekatan \textit{representation learning} (pembelajaran representasi) yang otomatis menemukan representasi data yang diperlukan untuk menyelesaikan tugas tertentu. \textit{Deep learning} merupakan bentuk spesifik dari \textit{representation learning} yang mampu mengekstraksi fitur secara otomatis langsung dari data mentah melalui struktur lapisan hierarkis yang merepresentasikan konsep secara bertingkat dari ANN. 

Secara konseptual, kedalaman dalam arsitektur ini memungkinkan komputer untuk mempelajari konsep yang rumit dengan membangunnya dari konsep-konsep yang lebih sederhana \parencite{goodfellowDeepLearning2016}. Sebagai contoh dalam pengolahan citra, lapisan awal jaringan dapat mempelajari deteksi tepian, lapisan menengah mempelajari kontur atau sudut, dan lapisan akhir mengenali bagian objek yang spesifik. Hierarki fitur ini terbentuk secara otomatis selama proses pelatihan tanpa intervensi manusia. Hal ini menjadi keunggulan utama dibandingkan metode tradisional, karena sistem dapat mengadaptasi ekstraksi fiturnya sendiri sesuai dengan karakteristik data yang spesifik. Perbandingan alur pengolahan data ini disajikan dalam Gambar \ref{fig:workflow_ml_dl}.

\begin{figure}[H] 
\centering
\includegraphics[max width=\textwidth, max height=0.6\textheight]{bab2/traditionalvsdeeplearning.png}
\caption[Perbandingan alur kerja antara \textit{machine learning} tradisional, \textit{representation learning}, dan \textit{deep learning}]{Perbandingan alur kerja antara \textit{machine learning} tradisional, \textit{representation learning}, dan \textit{deep learning} \parencite{goodfellowDeepLearning2016}}
\label{fig:workflow_ml_dl}
\end{figure}

Dalam implementasi \textit{deep learning}, terdapat konsep yang sering dikenal dengan \textit{pre-training}, di mana model dilatih terlebih dahulu pada \textit{dataset} berskala besar untuk mempelajari fitur umum sebelum diterapkan pada tugas yang lebih spesifik melalui \textit{transfer learning} \parencite{cholletDeepLearningPython2026}. Chollet menjelaskan bahwa fitur pada lapisan awal cenderung bersifat umum (\textit{general}), sementara fitur pada lapisan akhir lebih spesifik terhadap \textit{dataset} pelatihan asli. Oleh karena itu, terdapat dua strategi \textit{transfer learning} yang umum digunakan dalam pengembangan model \textit{deep learning}, yaitu \textit{feature extraction} dan \textit{fine-tuning}. Perbedaan mekanisme antara kedua strategi tersebut diilustrasikan pada Gambar \ref{fig:transfer_learning_mechanisms}.

\begin{figure}[H]
\centering
\includegraphics[width=\textwidth]{bab2/transferlearningmechanisms.png}
\caption[Contoh mekanisme \textit{feature extraction} dan \textit{fine-tuning} pada model CNN]{Contoh mekanisme \textit{feature extraction} dan \textit{fine-tuning} pada model CNN \parencite{asifDeepLearningbasedFramework2022}}
\label{fig:transfer_learning_mechanisms}
\end{figure}

Strategi \textit{feature extraction} bekerja dengan menggunakan bobot model yang sudah dilatih sebagai ekstraktor fitur tetap tanpa melakukan pembaruan bobot pada lapisan dasar model, sehingga hanya lapisan klasifikasi terakhir yang dilatih ulang \parencite{cholletDeepLearningPython2026}. Strategi ini ideal digunakan ketika \textit{dataset} target berukuran kecil atau memiliki kemiripan visual yang sangat tinggi dengan \textit{dataset} asal. Sebaliknya, strategi \textit{fine-tuning} melibatkan pelatihan kembali sebagian atau seluruh lapisan model menggunakan \textit{dataset} target dengan laju pembelajaran (\textit{learning rate}) yang lebih rendah guna mengadaptasi fitur yang sudah dipelajari ke domain data yang baru. Strategi ini lebih cocok diterapkan pada \textit{dataset} target yang berukuran besar atau memiliki karakteristik data yang berbeda secara signifikan dari \textit{dataset} asal.

Sebagian besar model \textit{deep learning} secara tradisional dilatih menggunakan pendekatan \textit{supervised learning} \parencite{goodfellowDeepLearning2016}. Momentum kebangkitan \textit{deep learning} modern pun dipicu oleh keberhasilan arsitektur AlexNet dari \textcite{krizhevskyImageNetClassificationDeep2017} yang dilatih secara supervisi untuk mengklasifikasikan citra pada \textit{dataset} ImageNet. Kendati demikian, \textcite{huyenAIEngineeringBuilding2025} memaparkan bahwa ketergantungan pada supervisi menghadirkan hambatan skalabilitas yang signifikan karena proses pelabelan data secara manual membutuhkan waktu dan biaya yang tinggi. Beban biaya ini akan meningkat seiring dengan bertambahnya jumlah kategori objek yang ingin dicakup oleh model. Selain itu, terdapat variabilitas dalam tingkat kesulitan pelabelan, yaitu meskipun pelabelan objek umum relatif murah, pelabelan pada domain spesifik yang menuntut kepakaran khusus memerlukan sumber daya yang jauh lebih besar.

Untuk mengatasi hambatan pelabelan data tersebut, dikembangkan paradigma \textit{self-supervised learning}. Dalam pendekatan ini, model tidak memerlukan label eksplisit, melainkan menginferensi label secara mandiri langsung dari informasi yang terkandung dalam data masukan. Kemampuan tersebut menjadi fondasi krusial bagi pengembangan \textit{foundation models} yang mampu mengintegrasikan berbagai modalitas data. Perkembangan ini yang kemudian melahirkan \textit{Vision-Language Models}, yang menggabungkan kemampuan pemahaman visual dan tekstual secara sinergis untuk menyelesaikan tugas-tugas penalaran yang kompleks \parencite{bordesIntroductionVisionLanguageModeling2024}. Pendekatan ini menjadi sangat relevan dalam mengatasi tantangan kelangkaan data berlabel pada domain-domain spesifik, termasuk dalam identifikasi penyakit tanaman yang memerlukan anotasi dari pakar.

% Note: Ini versi lama. Disimpan untuk referensi saja.
% Walaupun \textit{deep learning} menawarkan keunggulan dalam akurasi, model ini secara tradisional sangat bergantung pada ketersediaan data berlabel dalam jumlah besar terutama untuk pelatihan berbasis \textit{supervised learning} \parencite{goodfellowDeepLearning2016}. Keterbatasan ini mendorong perkembangan paradigma \textit{self-supervised learning}, di mana model dirancang untuk mempelajari representasi data dari data yang tidak berlabel dengan cara menciptakan tugas prediksi internal dari struktur data itu sendiri \parencite{minaeeLargeLanguageModels2025}. Paradigma ini memungkinkan pemanfaatan data dalam skala masif untuk membangun representasi fitur yang kaya tanpa bergantung pada anotasi manusia yang mahal. Kemampuan tersebut menjadi fondasi krusial bagi pengembangan \textit{foundation models} yang mampu mengintegrasikan berbagai modalitas data. Perkembangan ini yang kemudian melahirkan \textit{Vision-Language Models}, yang menggabungkan kemampuan pemahaman visual dan tekstual secara sinergis untuk menyelesaikan tugas-tugas penalaran yang kompleks \parencite{bordesIntroductionVisionLanguageModeling2024}. Pendekatan ini menjadi sangat relevan dalam mengatasi tantangan kelangkaan data berlabel pada domain-domain spesifik, termasuk dalam identifikasi penyakit tanaman yang memerlukan anotasi dari pakar.

\section{\textit{Large-Language Models}}
\textit{Large Language Models} (LLM) merupakan salah satu terobosan paling signifikan dalam bidang \textit{Natural Language Processing} (NLP) \parencite{cholletDeepLearningPython2026}. NLP didefinisikan sebagai subbidang AI yang berfokus pada pengembangan algoritma untuk memproses, memahami, serta menghasilkan bahasa alami manusia, baik dalam bentuk teks maupun suara \parencite{jm3}. Dalam hierarki teknologi, LLM menempati posisi sebagai model bahasa tingkat lanjut yang memanfaatkan teknik \textit{deep learning} untuk memproses informasi tekstual dalam skala yang sangat besar.

\begin{figure}[H]
\centering
\includegraphics[width=0.6\textwidth]{bab2/llmnlp.png}
\caption[Hubungan antara AI, NLP, dan LLM]{Hubungan antara AI, NLP, dan LLM \parencite{Bottino2024}}
\label{fig:ai_ml_dl_llm}
\end{figure}

Terminologi LLM terdiri dari dua komponen utama, yakni model bahasa (\textit{language model}) dan skala besar (\textit{large}). Model bahasa berfungsi untuk mengodekan informasi statistik mengenai satu atau lebih bahasa \parencite{huyenAIEngineeringBuilding2025}. Secara intuitif, informasi ini mengindikasikan seberapa besar kemungkinan suatu kata muncul dalam konteks tertentu. Dari sini, model bahasa dapat memprediksi kata yang paling mungkin muncul berikutnya dalam suatu \textit{sequence} (urutan) berdasarkan konteks yang diberikan. Mekanisme prediksi kata berikutnya ini memanfaatkan sifat sekuensial alami bahasa untuk melatih model agar mampu memahami konteks, struktur, serta hubungan antar elemen di dalam teks \parencite{raschkaBuildLargeLanguage2025}. Karena kemampuannya dalam membangkitkan teks tersebut, LLM juga sering diklasifikasikan sebagai bentuk kecerdasan buatan generatif (\textit{generative AI} atau \textit{GenAI}). Terminologi ini merujuk pada kemampuan model untuk menghasilkan luaran yang bersifat terbuka (\textit{open-ended}), di mana kosakata yang terbatas digunakan untuk mengonstruksi kemungkinan luaran yang tak terbatas \parencite{huyenAIEngineeringBuilding2025}.

% NOTE: Ini versi lama. Disimpan untuk referensi saja.
% Model bahasa secara fundamental adalah model probabilitas berbasis jaringan saraf yang menerima masukan berupa konteks atau prefiks teks, dan menghasilkan distribusi probabilitas terhadap kata atau token berikutnya yang paling mungkin muncul. Ilustrasi mekanisme dasar LLM sebagai model prediksi urutan kata disajikan pada Gambar \ref{fig:llm_illustration}. 
% \begin{figure}[H]
% \centering
% \includegraphics[width=0.8\textwidth]{bab2/llmilust.png}
% \caption[Ilustrasi mekanisme prediksi kata pada LLM]{Ilustrasi mekanisme LLM yang menerima konteks dan menghasilkan distribusi kata berikutnya \parencite{jm3}}
% \label{fig:llm_illustration}
% \end{figure}
Atribut "besar" pada LLM merujuk pada dua aspek utama, yaitu ukuran model dalam hal parameter serta \textit{dataset} masif yang digunakan sebagai basis pelatihannya \parencite{raschkaBuildLargeLanguage2025}. Model jenis ini umumnya memiliki puluhan hingga ratusan miliar parameter, yang merupakan bobot yang dapat disesuaikan dalam jaringan dan dioptimalkan selama pelatihan untuk memprediksi kata berikutnya dalam suatu urutan. Skala parameter ini berbanding lurus dengan kapasitas komputasi yang dibutuhkan, yang umumnya diukur dalam satuan FLOPS (\textit{Floating Point Operations per Second}) untuk merepresentasikan jumlah operasi aritmatika yang dilakukan per detik selama proses pelatihan maupun inferensi \parencite{amidiSuperStudyGuide2024}. Visualisasi mengenai korelasi antara jumlah parameter dan arsitektur model disajikan pada Gambar \ref{fig:model_parameters}.
\begin{figure}[H]
\centering
\includegraphics[width=0.8\textwidth]{bab2/modelparameters.png}
\caption[Visualisasi parameter pada model bahasa]{Visualisasi parameter pada model bahasa \parencite{amidiSuperStudyGuide2024}}
\label{fig:model_parameters}
\end{figure}
Perkembangan model bahasa modern saat ini hampir seluruhnya didasarkan pada arsitektur \textit{Transformer} yang diperkenalkan pertama kali oleh \textcite{vaswaniAttentionAllYou2023} melalui penelitian "Attention Is All You Need" \parencite{jm3,raschkaBuildLargeLanguage2025}. Arsitektur ini memperkenalkan mekanisme \textit{self-attention} yang memungkinkan model untuk memproses seluruh urutan kata secara paralel dan memahami ketergantungan antar kata terlepas dari jarak posisinya dalam kalimat. Inovasi ini menjadi fondasi utama bagi hampir seluruh LLM populer yang ada saat ini.

Dalam proses kerjanya, LLM mengandalkan mekanisme tokenisasi untuk memecah teks masukan menjadi unit-unit kecil yang disebut sebagai token \parencite{jm3}. Proses ini merupakan tahap fundamental pertama yang memungkinkan model untuk menginterpretasikan bahasa manusia ke dalam format yang dapat diproses secara komputasional. Terdapat berbagai strategi tokenisasi, mulai dari pemisahan berbasis kata, karakter, hingga sub-kata (\textit{sub-word}) seperti \textit{Byte-Pair Encoding} (BPE) yang mampu menangani kata-kata langka secara lebih efisien. Setelah melalui proses tokenisasi (\textit{encode}), urutan token tersebut dapat dikembalikan ke bentuk teks asli melalui proses sebaliknya (\textit{decode}), seperti yang diilustrasikan pada Gambar \ref{fig:encode_decode}.

\begin{figure}[H]
\centering
\includegraphics[width=1\textwidth]{bab2/encodedecode.png}
\caption[Mekanisme tokenisasi]{Mekanisme tokenisasi: (a) Proses \textit{encode} dari teks ke token dan (b) Proses \textit{decode} dari token ke teks \parencite{amidiSuperStudyGuide2024}}
\label{fig:encode_decode}
\end{figure}

Setiap token kemudian dikonversi menjadi representasi numerik dalam bentuk vektor dimensi tinggi yang dikenal sebagai \textit{embedding} \parencite{jm3}. Vektor-vektor ini ditempatkan dalam ruang fitur di mana kata-kata dengan makna serupa akan berada pada posisi yang berdekatan. Kedekatan atau kemiripan antar vektor tersebut umumnya diukur menggunakan teknik kesamaan vektor, seperti \textit{cosine similarity}, yang memungkinkan model untuk menangkap hubungan semantik antar token dalam dimensi ruang tertentu. Representasi ini memungkinkan model melakukan operasi aritmatika vektor untuk menangkap relasi analogi antar kata, seperti diilustrasikan pada Gambar \ref{fig:vector_representation}.

\begin{figure}[H]
\centering
\includegraphics[width=0.7\textwidth]{bab2/kingmanwomanqueen.png}
\caption[Visualisasi representasi vektor pada ruang \textit{embedding}]{Visualisasi representasi vektor pada ruang \textit{embedding} dengan contoh analogi $king - man + woman \approx queen$ oleh \textcite{mikolovEfficientEstimationWord2013} \parencite{sutorMetaconceptsIsolatingContext2019a}}
\label{fig:vector_representation}
\end{figure}

Kemampuan utama LLM dalam menghasilkan teks dilakukan melalui mekanisme \textit{autoregressive}, di mana model memprediksi token berikutnya berdasarkan urutan token sebelumnya secara berulang \parencite{jm3}. Visualisasi dari proses pembangkitan teks secara \textit{autoregressive} ini disajikan pada Gambar \ref{fig:llm_generation}. Proses pembangkitan ini dipandu oleh instruksi atau masukan teks awal yang disebut sebagai \textit{prompt}. Melalui konfigurasi \textit{prompt} yang tepat, model dapat diarahkan untuk menjalankan tugas-tugas tertentu tanpa memerlukan pembaruan parameter model secara permanen, sebuah fenomena yang dikenal sebagai \textit{in-context learning}.

\begin{figure}[H]
\centering
\includegraphics[width=0.8\textwidth]{bab2/llmgeneration.png}
\caption[Mekanisme pembangkitan teks autoregresif pada LLM]{Mekanisme pembangkitan teks autoregresif pada LLM \parencite{jm3}}
\label{fig:llm_generation}
\end{figure}

LLM dilatih melalui pendekatan \textit{self-supervised learning} \parencite{huyenAIEngineeringBuilding2025}. Pendekatan ini memungkinkan model untuk belajar dari urutan teks tanpa memerlukan pelabelan eksplisit, karena label diinferensi secara mandiri dari data masukan. Dalam konteks pemodelan bahasa, setiap urutan masukan menyediakan label berupa token yang akan diprediksi sekaligus konteks yang digunakan untuk melakukan prediksi tersebut. Mengingat ketersediaan urutan teks yang sangat melimpah, metode ini memungkinkan konstruksi data pelatihan dalam jumlah masif yang memfasilitasi peningkatan skala model bahasa menjadi LLM.

Meskipun memiliki kemampuan bahasa yang menyerupai manusia, LLM tetap memiliki batasan intrinsik berupa fenomena halusinasi \parencite{minaeeLargeLanguageModels2025}. Halusinasi merujuk pada kondisi di mana model menghasilkan informasi yang tampak masuk akal secara bahasa namun tidak memiliki dasar faktual atau tidak sesuai dengan data pelatihan yang ada. Hal ini karena LLM pada dasarnya bekerja sebagai model probabilitas yang memprediksi urutan kata yang paling mungkin, bukan sebagai sistem yang menyimpan pengetahuan faktual secara eksplisit. 

\textcite{minaeeLargeLanguageModels2025} mengidentifikasi berbagai metode augmentasi untuk mengurangi halusinasi pada LLM. Beberapa di antaranya adalah \textit{fine-tuning}, \textit{Retrieval-Augmented Generation} (RAG), penggunaan alat eksternal, serta pengembangan agen LLM. \textit{Fine-tuning} memungkinkan model beradaptasi dengan tugas tertentu melalui pelatihan ulang sebagian parameternya menggunakan data yang lebih kecil dan terkurasi. Akan tetapi, pendekatan ini menghadapi sejumlah kendala, seperti kebutuhan data pelatihan yang cukup besar, risiko terjadinya \textit{catastrophic forgetting} (hilangnya pengetahuan umum), serta tingginya biaya komputasi dalam pelatihan. Sebagai alternatif, RAG menawarkan solusi yang lebih fleksibel dengan menggabungkan LLM dan mekanisme pengambilan informasi dari sumber eksternal yang sudah disiapkan. Dengan cara ini, model dapat menghasilkan jawaban yang lebih relevan, akurat, dan dapat dipercaya.

Mekanisme pengambilan informasi eksternal ini, menurut \textcite{minaeeLargeLanguageModels2025}, merupakan salah satu contoh bagaimana LLM dapat menggunakan \textit{tools} (alat) eksternal. Alat eksternal ini pada dasarnya dapat berupa fungsi, API (\textit{Application Programming Interface}), atau layanan apa pun yang memberikan akses ke data dinamis, kalkulasi, atau kapabilitas lainnya. Lebih lanjut, kemampuan ini dapat ditingkatkan melalui konsep agen LLM yang dirancang untuk memungkinkan LLM menggunakan berbagai alat yang disediakan secara adaptif. 


\subsection{Arsitektur \textit{Transformer}}
Arsitektur \textit{Transformer} adalah arsitektur \textit{deep learning} yang pertama kali diperkenalkan oleh \textcite{vaswaniAttentionAllYou2023} untuk keperluan penerjemahan mesin (\textit{machine translation}).
Arsitektur ini awalnya digunakan untuk menerjemahkan teks bahasa Inggris ke bahasa Jerman dan Perancis \parencite{raschkaBuildLargeLanguage2025}. Kehadiran arsitektur ini kemudian membawa revolusi besar dalam bidang NLP serta berbagai tugas lain yang melibatkan pemetaan urutan-ke-urutan (\textit{sequence-to-sequence}) \parencite{altoBuildingLLMPowerd2024}. 

Secara fundamental, arsitektur \textit{Transformer} memanfaatkan mekanisme \textit{attention} untuk melakukan proses \textit{encode} dan \textit{decode} terhadap suatu urutan data \parencite{altoBuildingLLMPowerd2024}. Mekanisme ini hadir sebagai penyempurnaan dari metode sebelumnya yang diterapkan pada \textit{Recurrent Neural Networks} (RNN). Pada arsitektur RNN, pemrosesan data dilakukan secara sekuensial, di mana informasi diproses langkah demi langkah dengan meneruskan status tersembunyi (\textit{hidden state}) dari satu langkah ke langkah berikutnya. Pendekatan ini memiliki keterbatasan dalam menangkap hubungan ketergantungan jangka panjang. Sebaliknya, mekanisme \textit{attention} memungkinkan model untuk memusatkan perhatian pada bagian-bagian yang paling relevan dari urutan masukan saat menyusun luaran. 

Komponen utama yang menerapkan mekanisme ini dalam arsitektur \textit{Transformer} adalah lapisan \textit{Multi-Head Attention} \parencite{jm3}. Mekanisme ini sering juga disebut sebagai \textit{self-attention} karena diterapkan pada urutan yang sama dengan yang sedang dikodekan \parencite{altoBuildingLLMPowerd2024}. Secara sederhana, lapisan \textit{self-attention} bertugas untuk menentukan tingkat kepentingan dari setiap token masukan dalam hubungannya dengan token lain saat menghasilkan luaran. Dalam implementasinya, mekanisme \textit{self-attention} melibatkan tiga matriks bobot yang dipelajari selama pelatihan, yaitu \textit{Query} ($Q$), \textit{Key} ($K$), dan \textit{Value} ($V$). Di sini, $Q$ merepresentasikan token yang sedang diproses atau apa yang dicari, $K$ bertindak sebagai label identitas untuk setiap token dalam urutan atau apa yang tersedia, dan $V$ menyimpan konten informasi aktual dari token tersebut. Proses komputasi dimulai dengan menghitung skor atensi melalui perkalian $Q$ dan $K$ untuk menentukan seberapa relevan suatu token terhadap token lainnya. Skor ini kemudian dinormalisasi menggunakan fungsi \textit{Softmax} dan dikalikan dengan $V$ untuk menghasilkan vektor konteks yang telah diperkaya dengan informasi dari bagian input yang relevan. Mekanisme dekomposisi matriks input dan representasi operasi perkalian matriks ini diilustrasikan pada Gambar \ref{fig:attention_matrices}.

\begin{figure}[H]
\centering
\begin{subfigure}[b]{0.45\textwidth}
    \centering
    \includegraphics[width=\textwidth]{bab2/dekomposisimatriks.png}
    \caption{}
    \label{fig:dekomposisi_matriks}
\end{subfigure}
\hfill
\begin{subfigure}[b]{0.45\textwidth}
    \centering
    \includegraphics[width=\textwidth]{bab2/representasimatriks.png}
    \caption{}
    \label{fig:representasi_matriks}
\end{subfigure}
\caption[Dekomposisi dan representasi operasi matriks pada mekanisme \textit{self-attention}]{Dekomposisi dan representasi operasi matriks pada mekanisme \textit{self-attention}: (a) Dekomposisi matriks input menjadi vektor Q, K, dan V, serta (b) Representasi perkalian matriks Q, K, dan V untuk mendapatkan vektor konteks \parencite{altoBuildingLLMPowerd2024}}
\label{fig:attention_matrices}
\end{figure} 

Secara struktural, arsitektur \textit{Transformer} terdiri dari dua submodul utama, yakni pengode (\textit{encoder}) dan pendekode (\textit{decoder}) \parencite{raschkaBuildLargeLanguage2025}. Modul \textit{encoder} berfungsi untuk memproses teks masukan dan mengodekannya menjadi serangkaian representasi numerik atau vektor yang menangkap informasi kontekstual dari masukan tersebut. Selanjutnya, modul \textit{decoder} memanfaatkan vektor yang telah dikodekan ini untuk membangkitkan teks luaran. Dalam tugas penerjemahan bahasa, \textit{encoder} akan mengubah teks dari bahasa sumber menjadi representasi vektor, sedangkan \textit{decoder} akan menguraikan kode vektor tersebut untuk menghasilkan teks dalam bahasa target. Kedua modul ini tersusun atas lapisan-lapisan yang saling terhubung melalui mekanisme \textit{self-attention}. Keseluruhan arsitektur ini disajikan pada Gambar \ref{fig:transformer_architecture}.

\begin{figure}[H]
\centering
\includegraphics[max width=\linewidth, max totalheight=0.37\textheight, keepaspectratio]{bab1/arsitekturtransformer.png}
\caption[Arsitektur \textit{Transformer}]{Arsitektur \textit{Transformer} \parencite{vaswaniAttentionAllYou2023}}
\label{fig:transformer_architecture}
\end{figure}

Dalam pengembangannya, arsitektur \textit{encoder-decoder} dari \textcite{vaswaniAttentionAllYou2023} telah diadaptasi menjadi berbagai varian yang dioptimalkan untuk tugas-tugas spesifik. Adaptasi ini umumnya diklasifikasikan ke dalam tiga kategori utama berdasarkan bagian arsitektur yang digunakan, yaitu \textit{encoder-decoder}, \textit{encoder-only}, dan \textit{decoder-only} \parencite{amidiSuperStudyGuide2024}. Perbedaan struktur antara ketiga kategori arsitektur tersebut diilustrasikan pada Gambar \ref{fig:three_transformer_architectures}.

\begin{figure}[H]
\centering
\includegraphics[width=\textwidth]{bab2/tigaarsitekturtransformer.png}
\caption[Tiga varian arsitektur \textit{Transformer}]{Tiga varian arsitektur \textit{Transformer}: \textit{encoder-only}, \textit{encoder-decoder}, dan \textit{decoder-only} \parencite{jm3}}
\label{fig:three_transformer_architectures}
\end{figure}

BERT (\textit{Bidirectional Encoder Representations from Transformers}) dan GPT (\textit{Generative Pretrained Transformers}) merupakan contoh varian arsitektur \textit{Transformer} \parencite{raschkaBuildLargeLanguage2025}. BERT dibangun di atas submodul \textit{encoder} dari arsitektur asli \textit{Transformer}. Pendekatan pelatihannya berfokus pada prediksi kata yang disembunyikan atau \textit{masked word prediction}, di mana model memprediksi kata-kata yang ditutupi dalam suatu kalimat. Strategi ini memberikan BERT keunggulan dalam tugas-tugas klasifikasi teks, seperti prediksi sentimen dan kategorisasi dokumen. Di sisi lain, GPT berfokus pada bagian \textit{decoder} dari arsitektur \textit{Transformer} dan dirancang khusus untuk tugas-tugas generatif. Model ini menghasilkan teks kata demi kata dan mampu menangani berbagai tugas seperti penerjemahan mesin, peringkasan teks, hingga penulisan fiksi. Selain itu, model GPT juga menunjukkan kemampuan yang menonjol dalam melakukan \textit{zero-shot} dan \textit{few-shot learning}, yang memungkinkannya untuk menggeneralisasi tugas baru tanpa atau dengan sedikit contoh. Representasi visual dari perbedaan arsitektur \textit{encoder} pada BERT dan \textit{decoder} pada GPT dapat dilihat pada Gambar \ref{fig:bert_vs_gpt}.

\begin{figure}[H]
\centering
\includegraphics[max width=\linewidth, max totalheight=0.37\textheight, keepaspectratio]{bab2/encoderbertdecodergpt.png}
\caption[Representasi visual arsitektur \textit{encoder} pada BERT dan \textit{decoder} pada GPT]{Representasi visual arsitektur \textit{encoder} pada BERT dan \textit{decoder} pada GPT \parencite{raschkaBuildLargeLanguage2025}}
\label{fig:bert_vs_gpt}
\end{figure}

%TODO: Ini subsection ViT masih banyak klaim yang ga bersumber
\subsection{Arsitektur \textit{Vision Transformer}}
\textit{Vision Transformer} (ViT) diperkenalkan oleh \textcite{dosovitskiyImageWorth16x162020a} sebagai upaya untuk menerapkan arsitektur \textit{Transformer} standar secara langsung pada citra dengan modifikasi seminimal mungkin. Sebelum kemunculan ViT, mekanisme \textit{attention} dalam visi komputer umumnya diterapkan bersamaan dengan \textit{Convolutional Neural Networks} (CNN), atau digunakan untuk menggantikan komponen tertentu dari CNN sambil tetap mempertahankan struktur keseluruhannya. \textcite{dosovitskiyImageWorth16x162020a} menunjukkan bahwa ketergantungan pada CNN tidak sepenuhnya diperlukan dan \textit{Transformer} murni yang diterapkan pada urutan \textit{patch} citra dapat memberikan kinerja yang sangat baik pada tugas klasifikasi citra, terutama ketika dilatih pada \textit{dataset} berskala besar.

Mekanisme inti dari ViT adalah mengubah citra 2D menjadi urutan \textit{token} 1D agar dapat diproses oleh \textit{Transformer} standar yang dirancang untuk data sekuensial \textit{Vision Transformer} (ViT) diperkenalkan oleh \parencite{dosovitskiyImageWorth16x162020a} . Sebuah citra $x \in \mathbb{R}^{H \times W \times C}$ dipecah menjadi serangkaian \textit{patch} berukuran tetap $x_p \in \mathbb{R}^{N \times (P^2 \cdot C)}$, di mana $(H, W)$ adalah resolusi citra asli, $C$ adalah jumlah saluran, $(P, P)$ adalah resolusi setiap \textit{patch}, dan $N = HW/P^2$ adalah jumlah total \textit{patch} yang juga menjadi panjang urutan masukan bagi \textit{Transformer}. Setiap \textit{patch} kemudian diratakan (\textit{flattened}) dan dipetakan ke dalam dimensi model $D$ menggunakan proyeksi linier yang dapat dilatih, menghasilkan apa yang disebut sebagai \textit{patch embeddings}.

Untuk mempertahankan informasi posisi spasial yang hilang akibat perataan, \textit{position embeddings} ditambahkan ke setiap \textit{patch embeddings}. \textcite{dosovitskiyImageWorth16x162020a} menggunakan \textit{learnable position embeddings} 1D standar. Selain itu, mirip dengan \textit{token} [CLS] pada model BERT, sebuah \textit{embedding} khusus yang dapat dilatih ditambahkan di awal urutan \textit{patch} ($z_0^0 = x_{class}$). Keadaan keluaran dari \textit{token} ini pada lapisan terakhir \textit{encoder} ($z_L^0$) berfungsi sebagai representasi citra keseluruhan (\textit{image representation}) yang kemudian diteruskan ke \textit{classification head} (biasanya berupa MLP) untuk memprediksi kelas citra.

Arsitektur \textit{Transformer Encoder} pada ViT terdiri dari lapisan-lapisan yang bergantian antara \textit{Multi-Head Self-Attention} (MSA) dan blok \textit{Multi-Layer Perceptron} (MLP). Normalisasi lapisan (\textit{Layer Normalization} atau LN) diterapkan sebelum setiap blok, dan koneksi residual (\textit{residual connections}) diterapkan setelah setiap blok. Struktur keseluruhan dari \textit{Vision Transformer} diilustrasikan pada Gambar \ref{fig:vit_architecture}.

\begin{figure}[H]
\centering
\includegraphics[width=\textwidth]{bab2/arsitekturvisiontransformer.png} 
\caption[Arsitektur \textit{Vision Transformer}]{Arsitektur \textit{Vision Transformer} \parencite{dosovitskiyImageWorth16x162020a}}
\label{fig:vit_architecture}
\end{figure}

Berbeda dengan CNN yang memiliki bias induktif yang kuat terhadap invarian translasi dan lokalitas, ViT memiliki bias induktif yang jauh lebih sedikit \parencite{dosovitskiyImageWorth16x162020a}. Pada ViT, hubungan antar piksel atau \textit{patch} harus dipelajari dari awal melalui mekanisme \textit{self-attention}, yang memungkinkan model untuk menangkap konteks global (\textit{global context}) dari citra sejak lapisan awal. Hal ini memberikan keunggulan pada ViT dalam menangkap pola jarak jauh yang kompleks, namun juga mensyaratkan jumlah data pelatihan yang sangat besar agar model dapat menggeneralisasi dengan baik.

\subsection{\textit{Mixture-of-Experts}}
\textit{Mixture-of-Experts} (MoE) merupakan kerangka kerja jaringan saraf yang beroperasi dengan prinsip bahwa berbagai bagian dari sebuah model, yang dikenal sebagai ahli (\textit{experts}), memiliki spesialisasi pada aspek data yang berbeda \parencite{caiSurveyMixtureExperts2025}. Dengan paradigma ini, hanya ahli yang relevan yang diaktifkan untuk memproses suatu masukan, sehingga kapasitas pengetahuan model dapat diperluas tanpa menyebabkan lonjakan beban komputasi. Konsep dasar dari arsitektur ini pertama kali diperkenalkan oleh \textcite{jacobsAdaptiveMixturesLocal1991} melalui pendekatan (\textit{supervised learning}) yang membagi sebuah tugas kompleks menjadi beberapa subtugas yang masing-masing ditangani oleh jaringan ahli (\textit{expert network}) terpisah. Pendekatan tersebut kemudian dikembangkan lebih lanjut oleh \textcite{jordanHierarchicalMixturesExperts1993} dengan mengusulkan arsitektur terstruktur pohon (\textit{tree-structured architecture}) yang dikenal sebagai \textit{hierarchical mixtures of experts}. 

Integrasi konsep MoE ke dalam ranah \textit{deep learning} berskala besar kemudian berhasil diwujudkan oleh \textcite{shazeerOutrageouslyLargeNeural2017} melalui pengenalan lapisan \textit{Sparsely-Gated Mixture-of-Experts}. Penelitian tersebut memecahkan tantangan praktis dari komputasi bersyarat (\textit{conditional computation}) dengan menggunakan sebuah jaringan gerbang (\textit{gating network}). Jaringan gerbang tersebut dilatih untuk memilih kombinasi kecil dari ribuan jaringan ahli secara dinamis pada setiap masukan. Mekanisme ini memungkinkan kapasitas model ditingkatkan secara drastis hingga ratusan miliar parameter dengan tetap mempertahankan efisiensi komputasi.

Berbekal efisiensi komputasi tersebut, penerapan arsitektur MoE secara khusus pada LLM dipelopori oleh \textcite{lepikhinGShardScalingGiant2020} melalui model GShard. Penelitian ini mendemonstrasikan kelayakan pelatihan model multibahasa dengan 600 miliar parameter menggunakan teknik pemaralelan otomatis (\textit{automatic sharding}). GShard menggunakan mekanisme \textit{top-2 routing}, di mana setiap \textit{token} diproses oleh dua ahli yang paling sesuai, yang terbukti meningkatkan kualitas penerjemahan secara signifikan dibandingkan model \textit{dense}. Perbandingan arsitektur blok \textit{Transformer} standar dengan blok MoE pada GShard ditunjukkan pada Gambar \ref{fig:gshard_architecture}.

\begin{figure}[H]
\centering
\includegraphics[width=\textwidth]{bab2/arsitekturmoe.png}
\caption[Perbandingan arsitektur \textit{Transformer} standar dan MoE pada GShard]{Perbandingan arsitektur \textit{Transformer encoder} standar, MoE \textit{Transformer encoder}, dan MoE \textit{Transformer encoder} dengan \textit{device placement} \parencite{lepikhinGShardScalingGiant2020}}
\label{fig:gshard_architecture}
\end{figure}

Konsep ini kemudian dikembangkan lebih lanjut oleh \textcite{fedusSwitchTransformersScaling2022} melalui arsitektur \textit{Switch Transformer} yang menyederhanakan mekanisme pemilihan ahli. \textcite{fedusSwitchTransformersScaling2022} memperkenalkan strategi \textit{top-1 routing} di mana setiap \textit{token} hanya diarahkan ke satu ahli tunggal. Mekanisme tersebut diilustrasikan pada Gambar \ref{fig:moe_architecture}. Pendekatan ini mengurangi biaya komunikasi antar perangkat dan meningkatkan stabilitas pelatihan model berukuran triliunan parameter. Selain itu, untuk mengatasi masalah ketidakseimbangan beban di mana satu ahli mungkin menerima terlalu banyak \textit{token}, \textit{Switch Transformer} menerapkan mekanisme \textit{load balancing loss} sebagai penyeimbang. Arsitektur ini membuktikan bahwa skala parameter model dapat dipisahkan dari biaya komputasi (FLOPS), memungkinkan peningkatan performa model yang berkelanjutan. 

\begin{figure}[H]
\centering
\includegraphics[width=0.8\textwidth]{bab2/arsitekturswitchtransformer.png}
\caption[Mekanisme \textit{routing} pada \textit{Switch Transformer}]{Mekanisme \textit{routing} pada \textit{Switch Transformer} dengan strategi \textit{top-1 routing} \parencite{fedusSwitchTransformersScaling2022}}
\label{fig:moe_architecture}
\end{figure}

Saat ini, arsitektur MoE telah menjadi standar dalam pengembangan \textit{Large Language Models} modern guna mencapai performa tinggi dengan efisiensi inferensi yang optimal \parencite{caiSurveyMixtureExperts2025}. Dominasi arsitektur ini didorong oleh kemampuannya dalam menyeimbangkan skala parameter yang masif dengan biaya operasional yang tetap rendah. Beberapa model yang mengadopsi arsitektur ini antara lain Mixtral 8x7B \parencite{jiangMixtralExperts2024}, DeepSeek-V3 \parencite{deepseek-aiDeepSeekV3TechnicalReport2024}, hingga keluarga model Gemini \parencite{comaniciGemini25Pushing2025}.

\section{Model \textit{Multimodal}}
Model \textit{multimodal} merupakan paradigma pengembangan AI yang dirancang untuk dapat bekerja dengan lebih dari satu \textit{modality} atau jenis data \parencite{huyenAIEngineeringBuilding2025}. Mengacu pada \textcite{liangFoundationsTrendsMultimodal2023}, \textit{modality} didefinisikan sebagai cara suatu fenomena alam dipersepsikan atau diekspresikan, yang merentang dari data mentah (\textit{raw}) seperti sinyal suara atau citra, hingga data abstrak (\textit{abstract}) seperti bahasa atau kategori objek. Secara lebih umum, \textit{multimodal} merujuk pada pengekspresian atau persepsi hal-hal kompleks melalui berbagai \textit{modality} \parencite{10386743}. \textit{Multimodal} dapat diklasifikasikan menjadi \textit{modality} homogen, seperti citra yang diambil dari dua kamera berbeda, dan \textit{modality} heterogen, seperti hubungan antara citra dan bahasa tekstual. Visualisasi definisi ini diilustrasikan pada Gambar \ref{fig:multimodal_definition}.

\begin{figure}[H]
\centering
\includegraphics[width=0.8\textwidth]{bab2/multimodal.png}
\caption[Visualisasi definisi \textit{multimodal}]{Visualisasi definisi \textit{multimodal} \parencite{10386743}}
\label{fig:multimodal_definition}
\end{figure} 

Menurut \textcite{wuMultimodalLargeLanguage2023}, pengembangan teknis model \textit{multimodal} mencakup lima aspek fundamental, yaitu representasi pengetahuan, pemilihan tujuan pembelajaran, konstruksi struktur model, fusi informasi, dan penggunaan \textit{prompts}. Kelima aspek yang diilustrasikan pada Gambar \ref{fig:multimodal_model_points} ini dijabarkan sebagai berikut. Representasi pengetahuan dimulai dengan proses tokenisasi dan \textit{embedding} untuk mengubah data teks dan citra menjadi format vektor yang dapat diproses. Untuk teks, metode seperti \textit{byte pair encoding} umum digunakan, sedangkan citra diproses melalui pendekatan berbasis \textit{region}, \textit{grid}, atau \textit{patch} yang membagi gambar menjadi bagian-bagian kecil. Setelah data terwakili, pemilihan tujuan pembelajaran menjadi krusial untuk melatih model memahami hubungan antar \textit{modality}. Teknik seperti \textit{Image-Text Contrast} (ITC) digunakan untuk menyelaraskan representasi citra dan teks, sementara \textit{Masked Language Modeling} (MLM) dan \textit{Masked Visual Modeling} (MVM) melatih model merekonstruksi informasi yang hilang untuk memahami konteks secara mendalam.

\begin{figure}[H]
\centering
\includegraphics[width=0.8\textwidth]{bab2/multimodalmodelpoints.png}
\caption[Aspek utama dalam model \textit{multimodal}]{Aspek utama dalam model \textit{multimodal} \parencite{wuMultimodalLargeLanguage2023}}
\label{fig:multimodal_model_points}
\end{figure}

Konstruksi struktur model \textit{multimodal} umumnya dikategorikan menjadi arsitektur \textit{encoder-only} dan \textit{encoder-decoder}. Arsitektur \textit{encoder-only} memproses masukan secara langsung melalui \textit{encoder} dan lebih cocok untuk tugas pencarian (\textit{retrieval}), sedangkan arsitektur \textit{encoder-decoder} memanfaatkan kemampuan \textit{decoder} untuk membangkitkan luaran secara autoregresif sehingga ideal untuk tugas generatif. Fusi informasi (\textit{information fusion}) berfokus pada mekanisme penggabungan fitur multimodal, baik melalui fusi awal (\textit{early fusion}) maupun fusi akhir (\textit{late fusion}), untuk menghasilkan pemahaman yang komprehensif. Dalam konteks ini, terdapat pendekatan \textit{fusion encoder} dan \textit{dual encoder}. \textit{Fusion encoder} memungkinkan interaksi mendalam antar \textit{modality} melalui mekanisme \textit{attention} yang memberikan hasil inferensi yang kuat namun dengan biaya komputasi yang lebih tinggi. Sebaliknya, \textit{dual encoder} memproses setiap \textit{modality} secara terpisah sebelum menghitung kesamaan menggunakan operasi sederhana, menjadikannya lebih efisien untuk tugas pencarian data.

Penggunaan \textit{prompt} memainkan peran penting dalam menjembatani kesenjangan antara fase pra-pelatihan dan aplikasi pada tugas hilir. Dengan memodifikasi format tugas melalui instruksi tekstual, \textit{prompt} membantu model beradaptasi dengan tugas baru tanpa memerlukan pembaruan parameter yang ekstensif, sehingga mengurangi biaya \textit{fine-tuning}. Pendekatan ini memungkinkan model menangani skenario dengan sedikit atau tanpa contoh data. Dalam praktiknya, penggunaan \textit{prompt} yang informatif dapat meningkatkan pemahaman model terhadap konteks visual dan tekstual, serta memfasilitasi interaksi yang lebih alami antara pengguna dan sistem \textit{multimodal}. 

\textcite{huyenAIEngineeringBuilding2025} mengemukakan bahwa perkembangan pesat model \textit{multimodal} saat ini berjalan seiring dengan kesuksesan arsitektur \textit{Large Language Models} (LLM) berbasis \textit{Transformer} yang telah berevolusi menjadi \textit{foundation models}. Model-model generatif ini menjalani pelatihan ekstensif menggunakan himpunan data yang luas dan beragam, yang memungkinkannya untuk menangkap pola umum dan hubungan dalam data yang tidak hanya terbatas pada teks, tetapi juga mencakup format data lain seperti citra, audio, dan video \parencite{altoBuildingLLMPowerd2024}. \textcite{huyenAIEngineeringBuilding2025} merujuk pada model \textit{multimodal} generatif ini sebagai \textit{Large Multimodal Models} (LMM), sementara \textcite{altoBuildingLLMPowerd2024} menggunakan istilah \textit{Large Foundation Models} (LFM), dan \textcite{yinSurveyMultimodalLarge2024} menyebutnya sebagai \textit{Multimodal Large Language Models} (MLLM). Walaupun menggunakan terminologi yang berbeda, ketiga istilah tersebut pada dasarnya mengacu pada konsep yang sama. Kemampuan model-model ini dalam mengintegrasikan berbagai \textit{modality} data menjadi fondasi penting bagi pengembangan sistem kecerdasan buatan yang lebih adaptif dan mampu memahami kompleksitas dunia nyata secara holistik.

\subsection{\textit{Multimodal Large Language Models} (MLLM)}
\textcite{yinSurveyMultimodalLarge2024} mendefinisikan MLLM secara formal sebagai model berbasis LLM yang memiliki kemampuan untuk menerima, melakukan penalaran, dan menghasilkan luaran dengan informasi \textit{multimodal}. Arsitektur tipikal MLLM terdiri dari \textit{encoder}, konektor, dan LLM. Sebuah generator opsional dapat ditambahkan pada LLM untuk menghasilkan \textit{modality} lain selain teks. \textit{Encoder} berfungsi menerima masukan berupa citra, audio, atau video dan menghasilkan fitur, yang kemudian diproses oleh konektor agar dapat dipahami oleh LLM. Ilustrasi arsitektur ini disajikan pada Gambar \ref{fig:mllm_architecture}.

\begin{figure}[H]
\centering
\includegraphics[width=\textwidth]{bab2/mllm.png}
\caption[Ilustrasi arsitektur tipikal MLLM]{Ilustrasi arsitektur tipikal MLLM yang mencakup \textit{encoder}, konektor, dan LLM \parencite{yinSurveyMultimodalLarge2024}}
\label{fig:mllm_architecture}
\end{figure}

Secara umum, pengembangan MLLM dapat diabstraksikan menjadi tiga modul utama, yaitu \textit{modality encoder} yang telah dilatih sebelumnya (\textit{pre-trained}), LLM yang telah dilatih sebelumnya, dan antarmuka \textit{modality} (\textit{modality interface}) untuk menghubungkan keduanya. \textcite{yinSurveyMultimodalLarge2024} menganalogikan struktur ini dengan manusia, di mana \textit{modality encoder} (seperti \textit{encoder} citra atau audio) berperan sebagai mata atau telinga yang menerima dan memproses sinyal optik atau akustik, sedangkan LLM bertindak layaknya otak yang memahami dan melakukan penalaran terhadap sinyal yang telah diproses tersebut.

Modul pertama, \textit{modality encoder}, bertugas mengompresi informasi mentah seperti citra atau audio menjadi representasi yang lebih ringkas \parencite{yinSurveyMultimodalLarge2024}. Penggunaan \textit{encoder} yang telah dilatih sebelumnya dan diselaraskan dengan \textit{modality} lain lebih umum dilakukan daripada melatih model dari awal. Sebagai contoh, CLIP \parencite{radfordLearningTransferableVisual2021} menggunakan \textit{visual encoder} yang secara semantik telah diselaraskan dengan teks melalui pra-pelatihan skala besar. Penggunaan \textit{encoder} semacam ini memudahkan proses penyelarasan dengan LLM. Modul kedua adalah LLM pra-latih. Memulai pengembangan dengan LLM yang telah dilatih sebelumnya dinilai lebih efisien dan praktis. Melalui pelatihan ekstensif pada korpus teks yang masif, LLM telah dilengkapi dengan pengetahuan dunia yang kaya serta kemampuan generalisasi dan penalaran yang kuat.

Modul ketiga adalah antarmuka \textit{modality} \parencite{yinSurveyMultimodalLarge2024}. Mengingat LLM hanya memproses teks dan tingginya biaya komputasi pelatihan \textit{end-to-end}, diperlukan jembatan untuk menghubungkan kesenjangan antara bahasa alami dan \textit{modality} lainnya. Oleh karena itu, pendekatan yang lebih praktis adalah dengan memperkenalkan konektor yang dapat dilatih (\textit{learnable connector}) di antara \textit{visual encoder} dan LLM. Dalam pendekatan menggunakan konektor, komponen ini berfungsi untuk menyelaraskan representasi dari \textit{modality} yang berbeda agar dapat diproses oleh LLM. Terdapat tiga jenis konektor utama, yaitu berbasis proyeksi, berbasis \textit{query}, dan berbasis fusi. Dua jenis pertama mengadopsi fusi tingkat token yang memproses fitur menjadi token untuk dikirim bersama token teks, sedangkan jenis terakhir memungkinkan fusi tingkat fitur di dalam LLM. Sebagai alternatif, pendekatan lain memanfaatkan model ahli (\textit{expert models}) seperti \textit{image captioning} untuk mengonversi masukan \textit{multimodal} menjadi teks, sehingga dapat diproses langsung oleh LLM tanpa pelatihan tambahan. Meskipun pendekatan ini lebih sederhana, penggunaannya tidak sefleksibel antarmuka yang dapat dilatih karena konversi modalitas asing ke dalam bentuk teks berpotensi menyebabkan hilangnya informasi.

\subsection{\textit{Vision-Language Models} (VLM)}
\textit{Vision-Language Models} (VLM) adalah kelas model kecerdasan buatan \textit{multimodal} yang dirancang untuk memproses dan menghubungkan informasi dari \textit{modality} visual (citra atau video) dan tekstual (bahasa alami) secara bersamaan. Tujuan utama dari VLM adalah untuk menjembatani kesenjangan semantik antara representasi visual dan linguistik, sehingga memungkinkan mesin untuk "melihat" dan "memahami" dunia visual serta mengomunikasikannya melalui bahasa manusia \parencite{bordesIntroductionVisionLanguageModeling2024}. Kemampuan ini menjadi fondasi bagi berbagai aplikasi cerdas, mulai dari pencarian gambar berbasis teks, deskripsi gambar otomatis, hingga sistem tanya jawab visual yang kompleks.

Berdasarkan strategi pelatihan dan arsitekturnya, \textcite{bordesIntroductionVisionLanguageModeling2024} mengklasifikasikan VLM ke dalam empat keluarga utama yang menggambarkan evolusi integrasi \textit{modality} visual dan tekstual, sebagaimana disajikan pada Gambar \ref{fig:vlm_families}. Kategori pertama adalah \textit{contrastive learning} yang berfokus pada penyelarasan ruang representasi antara citra dan teks menggunakan fungsi kerugian kontras untuk mendekatkan representasi pasangan data yang cocok. Kategori kedua, \textit{masking}, menerapkan mekanisme rekonstruksi bagian input yang ditutupi untuk mempelajari hubungan kontekstual antara elemen visual dan tekstual. Selanjutnya, kategori \textit{generative} mencakup model yang dilatih dengan tujuan utama membangkitkan data secara autoregresif, baik berupa teks dari citra maupun sebaliknya. Terakhir, kategori \textit{pretrained backbones} memanfaatkan model bahasa dan \textit{vision encoder} yang telah dilatih sebelumnya sebagai komponen dasar yang dibekukan, di mana pelatihan hanya difokuskan pada modul penghubung untuk memetakan fitur visual ke dalam ruang input model bahasa.

\begin{figure}[H]
\centering
\includegraphics[width=0.7\textwidth]{bab2/vlmfamilies.png}
\caption[Empat Kategori Utama \textit{Vision-Language Models}]{Empat Kategori Utama \textit{Vision-Language Models}: (1) \textit{Contrastive}, (2) \textit{Masking}, (3) \textit{Generative}, dan (4) \textit{Pretrained Backbones} \parencite{bordesIntroductionVisionLanguageModeling2024}}
\label{fig:vlm_families}
\end{figure}

Evaluasi kinerja VLM dilakukan menggunakan berbagai metrik dan tugas standar untuk mengukur kemampuan visio-linguistiknya secara komprehensif. Taksonomi evaluasi VLM sebagaimana dijelaskan oleh \textcite{bordesIntroductionVisionLanguageModeling2024} divisualisasikan pada Gambar \ref{fig:vlm_evaluation}. Salah satu tugas evaluasi yang paling fundamental adalah \textit{Visual Question Answering} (VQA). Dalam VQA, model diberikan sebuah citra dan pertanyaan bahasa alami terkait citra tersebut, kemudian diminta untuk menghasilkan jawaban yang akurat. Tugas ini menuntut kemampuan penalaran tingkat tinggi karena model tidak hanya harus mengenali objek dalam citra, tetapi juga memahami hubungan spasial, aktivitas, dan konteks semantik yang ditanyakan \parencite{bordesIntroductionVisionLanguageModeling2024}. VQA dianggap sebagai proksi yang kuat untuk mengukur pemahaman \textit{multimodal} yang mendalam. Selain VQA, evaluasi juga mencakup kemampuan \textit{zero-shot classification}, penalaran komposisional (\textit{compositional reasoning}) untuk mendeteksi nuansa linguistik yang halus, serta pengujian terhadap halusinasi objek untuk memastikan reliabilitas model dalam aplikasi dunia nyata.

\begin{figure}[H]
\centering
\includegraphics[width=0.8\textwidth]{bab2/vlmevaluation.png}
\caption[Taksonomi Evaluasi \textit{Vision-Language Models}]{Taksonomi Evaluasi \textit{Vision-Language Models} \parencite{bordesIntroductionVisionLanguageModeling2024}}
\label{fig:vlm_evaluation}
\end{figure}

\subsection{\textit{Contrastive Language-Image Pre-training} (CLIP)}
\textit{Contrastive Language-Image Pre-training} (CLIP) merupakan model \textit{multimodal} yang diperkenalkan oleh \textcite{radfordLearningTransferableVisual2021} untuk menjembatani kesenjangan antara \textit{computer vision} dan pemahaman bahasa alami. Berbeda dengan pendekatan visi komputer konvensional yang bergantung pada \textit{dataset} dengan label kelas yang kaku dan terbatas, CLIP dilatih menggunakan pengawasan alami (\textit{natural language supervision}) dari pasangan citra dan teks yang tersedia secara melimpah di internet. Pendekatan ini memungkinkan model untuk mempelajari representasi visual yang kaya dan fleksibel, serta memiliki kemampuan \textit{zero-shot transfer} yang memungkinkannya mengenali kategori objek baru tanpa perlu dilatih ulang secara spesifik.

\textcite{radfordLearningTransferableVisual2021} merancang arsitektur CLIP dengan dua komponen utama, yaitu \textit{Image Encoder} dan \textit{Text Encoder}. \textit{Image Encoder} bertugas untuk memproses data visual menjadi vektor fitur, yang dapat menggunakan arsitektur ResNet atau \textit{Vision Transformer} (ViT). Sementara itu, \textit{Text Encoder} menggunakan arsitektur \textit{Transformer} untuk mengubah deskripsi teks menjadi vektor fitur dalam ruang dimensi yang sama. Mekanisme pelatihan CLIP didasarkan pada pembelajaran kontras (\textit{contrastive learning}), di mana model dilatih untuk memprediksi pasangan citra dan teks yang benar dari sekumpulan kemungkinan pasangan acak dalam sebuah \textit{batch}. Secara matematis, CLIP memaksimalkan kesamaan kosinus (\textit{cosine similarity}) antara \textit{embedding} citra dan teks yang berpasangan ($N$ pasangan asli) sambil meminimalkan kesamaan antara pasangan yang tidak cocok ($N^2 - N$ pasangan salah). Arsitektur dan mekanisme pelatihan CLIP diilustrasikan pada Gambar \ref{fig:clip_architecture}.

\begin{figure}[H]
\centering
\includegraphics[width=\textwidth]{bab2/arsitekturclip.png}
\caption[Arsitektur dan mekanisme pelatihan CLIP]{Arsitektur dan mekanisme pelatihan CLIP \parencite{radfordLearningTransferableVisual2021}}
\label{fig:clip_architecture}
\end{figure}

Keunggulan CLIP dalam mempelajari representasi visual yang kuat telah menginspirasi pengembangan model turunan untuk domain spesifik, salah satunya adalah SCOLD (\textit{Soft-target COntrastive learning for Leaf Disease identification}) yang diusulkan oleh \textcite{quocVisionLanguageFoundationModel2025}. SCOLD merupakan \textit{Vision-Language Foundation Model} yang secara khusus dikembangkan untuk pengenalan penyakit tanaman. Model ini mengadaptasi arsitektur CLIP dan menyempurnakannya melalui teknik \textit{fine-tuning} serta penerapan \textit{Deep Metric Learning}.

Dalam pengembangannya, \textcite{quocVisionLanguageFoundationModel2025} merancang SCOLD untuk mengatasi tantangan variabilitas visual yang tinggi pada gejala penyakit tanaman yang sering kali sulit dibedakan oleh model umum. Arsitektur SCOLD mempertahankan struktur \textit{encoder} ganda dari CLIP, namun dioptimalkan untuk memetakan fitur penyakit tanaman ke dalam ruang representasi yang lebih diskriminatif. Hal ini dicapai dengan melatih model agar jarak antara representasi citra penyakit dan deskripsi teksnya menjadi lebih dekat, sementara jarak dengan kelas penyakit lain diperlebar. Visualisasi arsitektur SCOLD dapat dilihat pada Gambar \ref{fig:scold_architecture}.

\begin{figure}[H]
\centering
\includegraphics[width=\textwidth]{bab2/arsitekturscold.png}
\caption[Arsitektur SCOLD]{Arsitektur SCOLD untuk pengenalan penyakit tanaman \parencite{quocVisionLanguageFoundationModel2025}}
\label{fig:scold_architecture}
\end{figure}

Efektivitas SCOLD dibandingkan dengan model CLIP standar dievaluasi oleh \textcite{quocVisionLanguageFoundationModel2025} menggunakan \textit{dataset} LeafNet, yang dibuat secara khusus untuk mengembangkan SCOLD, melalui tugas pengambilan lintas \textit{modal} (\textit{cross-modal retrieval}), yaitu \textit{Image-to-Text Retrieval} dan \textit{Text-to-Image Retrieval}. Hasil evaluasi menunjukkan bahwa SCOLD secara konsisten mengungguli CLIP pada metrik Recall@K (R@1, R@5, dan R@10) untuk kedua tugas tersebut. Peningkatan performa ini mengindikasikan bahwa proses \textit{fine-tuning} dan adaptasi domain yang dilakukan pada SCOLD berhasil meningkatkan keselarasan semantik antara fitur visual gejala penyakit dan deskripsi tekstualnya. Perbandingan kinerja pengambilan informasi antara CLIP dan SCOLD disajikan pada Tabel \ref{tab:scold_performance}.

\begin{table}[H]
    \centering
    \caption[Perbandingan performa retrieval SCOLD dan CLIP]
    {Perbandingan \textit{retrieval} SCOLD dan CLIP pada \textit{dataset} LeafNet \parencite{quocVisionLanguageFoundationModel2025}}
    \label{tab:scold_performance}
    \setstretch{1.0}
    \begin{tblr}{
        width = \linewidth,
        colspec = {X[l, m] *{6}{X[c, m]}},
        hlines, vlines,
        rows = {m},
        rowsep = 4pt,
        row{1,2} = {font=\bfseries, halign=c, valign=m},
        cell{1}{1} = {r=2}{},
        cell{1}{2} = {c=3}{},
        cell{1}{5} = {c=3}{},
    }
    Metode & \textit{Image-to-Text} & & & \textit{Text-to-Image} & & \\
    & R@1 (\%) & R@5 (\%) & R@10 (\%) & R@1 (\%) & R@5 (\%) & R@10 (\%) \\
    CLIP & 95.01 & 97.89 & 99.85 & 95.31 & 98.18 & 99.89 \\
    SCOLD
        & \textbf{95.49} {\footnotesize(+0.48)}
        & \textbf{98.85} {\footnotesize(+0.96)}
        & \textbf{99.93} {\footnotesize(+0.08)}
        & \textbf{95.46} {\footnotesize(+0.15)}
        & \textbf{99.70} {\footnotesize(+1.52)}
        & \textbf{99.99} {\footnotesize(+0.10)} \\
    \end{tblr}
\end{table}


% Efektivitas SCOLD dibandingkan dengan model CLIP standar dievaluasi oleh \textcite{quocVisionLanguageFoundationModel2025} menggunakan \textit{dataset} LeafNet, yang dibuat secara khusus untuk mengembangkan SCOLD, melalui tugas pengambilan lintas \textit{modal} (\textit{cross-modal retrieval}), yaitu \textit{Image-to-Text Retrieval} dan \textit{Text-to-Image Retrieval}. Hasil evaluasi menunjukkan bahwa SCOLD secara konsisten mengungguli CLIP pada metrik \textit{Recall@K} (R@1, R@5, dan R@10). Pada tugas \textit{Image-to-Text Retrieval}, SCOLD mencatatkan akurasi R@1 sebesar 95,49\%, R@5 sebesar 98,85\%, dan R@10 sebesar 99,93\%, unggul dibandingkan CLIP yang masing-masing memperoleh 95,01\%, 97,89\%, dan 99,85\%. Tren serupa terlihat pada \textit{Text-to-Image Retrieval}, di mana SCOLD mencapai R@1 sebesar 95,46\%, R@5 sebesar 99,7\%, dan R@10 sebesar 99,99\%, melampaui CLIP yang mencatatkan 95,31\%, 98,18\%, dan 99,89\%. Peningkatan performa ini mengindikasikan bahwa proses \textit{fine-tuning} dan adaptasi domain yang dilakukan pada SCOLD berhasil meningkatkan keselarasan semantik antara fitur visual gejala penyakit dan deskripsi tekstualnya \parencite{quocVisionLanguageFoundationModel2025}.

\subsection{Keluarga Model Gemini}
Gemini merupakan keluarga model \textit{multimodal} yang dikembangkan oleh Google DeepMind \parencite{teamGeminiFamilyHighly2025}. Berbeda dengan pendekatan konvensional yang melatih komponen secara terpisah untuk \textit{modality} yang berbeda, Gemini dibangun secara natif (\textit{natively multimodal}). Model ini dilatih sejak awal pada sekumpulan data \textit{multimodal} yang mencakup teks, kode, gambar, audio, dan video secara bersamaan. Pendekatan ini memungkinkan model untuk memiliki pemahaman dan penalaran lintas \textit{modality} yang kuat serta kemampuan untuk menerima masukan dan menghasilkan luaran dalam berbagai format secara mulus \parencite{teamGeminiFamilyHighly2025}. Gambaran umum mengenai kemampuan \textit{multimodal} model ini diilustrasikan pada Gambar \ref{fig:gemini_overview}.

\begin{figure}[H]
\centering
\includegraphics[width=0.8\textwidth]{bab2/geminimodeloverview.png}
\caption[Gambaran umum kemampuan \textit{multimodal} Gemini]{Gambaran umum kemampuan \textit{multimodal} Gemini \parencite{teamGeminiFamilyHighly2025}}
\label{fig:gemini_overview}
\end{figure}

Evolusi keluarga model ini dimulai dengan Gemini 1.0 yang memperkenalkan tiga varian utama, yaitu Ultra, Pro, dan Nano, masing-masing dioptimalkan untuk tingkat kompleksitas tugas dan efisiensi sumber daya yang berbeda. Pengembangan model ini dilanjutkan dengan diperkenalkannya Gemini 1.5 \parencite{teamGemini15Unlocking2024}. Inovasi utama pada Gemini 1.5 adalah perluasan jendela konteks hingga jutaan token, yang memfasilitasi pemrosesan dokumen panjang atau basis kode yang besar dalam satu kali interaksi.

Evolusi selanjutnya hadir melalui keluarga Gemini 2.5, yang difokuskan pada peningkatan kemampuan penalaran lanjut dan menyeimbangkan performa tinggi dengan efisiensi operasional \parencite{comaniciGemini25Pushing2025}. Gemini 2.5 dibangun menggunakan arsitektur \textit{Mixture-of-Experts}, di mana model hanya mengaktifkan subset parameter yang relevan untuk setiap input, sehingga meningkatkan efisiensi inferensi tanpa mengorbankan kualitas luaran. Model ini juga dilatih menggunakan teknik \textit{Chain-of-Thought}, yang memberikannya kemampuan penalaran dalam menyelesaikan masalah logika kompleks, matematika, dan pemrograman. Gemini 2.5 dirancang untuk menangani instruksi yang kompleks serta mempertahankan koherensi dalam tugas-tugas yang membutuhkan perencanaan bertahap. Seri model ini mencakup varian Gemini 2.5 Pro yang dioptimalkan untuk performa penalaran tertinggi, serta Gemini 2.5 Flash yang dirancang untuk efisiensi dan latensi rendah. Kemampuan varian Pro dalam berbagai \textit{benchmark} diperlihatkan pada Tabel \ref{tab:gemini25_pro_benchmark}.

{
\setstretch{1.0}
\begin{longtblr}[
    caption = {Hasil \textit{benchmark} Gemini 2.5 Pro pada berbagai tugas \parencite{comaniciGemini25Pushing2025}},
    entry = {Hasil \textit{benchmark} Gemini 2.5 Pro pada berbagai tugas},
    label = {tab:gemini25_pro_benchmark},
]{
    width = \linewidth,
    colspec = {X[1.4,l] X[1.7,l] *{4}{X[c]}},
    rows = {m},
    rowhead = 1,
    row{1} = {font=\bfseries, halign=c, valign=m},
    row{16} = {font=\bfseries, halign=c, valign=m},
    hlines, vlines,
    rowsep = 3pt,
}
    % ---- BLOK 1: Header ----
    \textit{Capability} & \textit{Benchmark} & {Gemini 2.5\\Pro} & {o3\\high} & {o4-mini\\high} & {Claude 4\\Sonnet} \\
    % ---- Code ----
    \textit{Code} & LiveCodeBench & 74.2\% & 72.0\% & \textbf{75.8\%} & 48.9\% \\
    \textit{Code} & Aider Polyglot & \textbf{82.2\%} & 79.6\% & 72.0\% & 61.3\% \\
    \textit{Code} & SWE-bench Verified \textit{(single)} & 59.6\% & 69.1\% & 68.1\% & \textbf{72.7\%} \\
    \textit{Code} & SWE-bench Verified \textit{(multiple)} & 67.2\% & -- & -- & \textbf{80.2\%} \\
    % ---- Reasoning ----
    \textit{Reasoning} & GPQA (diamond) \textit{(single)} & \textbf{86.4\%} & 83.3\% & 81.4\% & 75.4\% \\
    \textit{Reasoning} & Humanity's Last Exam \textit{(no tools)} & \textbf{21.6\%} & 20.3\% & 18.1\% & 7.8\% \\
    % ---- Factuality ----
    \textit{Factuality} & SimpleQA & \textbf{54.0\%} & 48.6\% & 19.3\% & -- \\
    \textit{Factuality} & FACTS Grounding & \textbf{87.8\%} & 69.9\% & 62.1\% & 79.1\% \\
    % ---- Math ----
    \textit{Math} & AIME 2025 \textit{(single)} & 88.0\% & 88.9\% & \textbf{92.7\%} & 70.5\% \\
    % ---- Long-context ----
    \textit{Long-context} & LOFT (hard retrieval) $\le$128K & \textbf{87.0\%} & 77.0\% & 60.5\% & 81.6\% \\
    \textit{Long-context} & LOFT (hard retrieval) 1M & \textbf{69.8\%} & -- & -- & -- \\
    \textit{Long-context} & MRCR-V2 (8-needle) $\le$128K & \textbf{58.0\%} & 57.1\% & 36.3\% & 39.1\% \\
    \textit{Long-context} & MRCR-V2 (8-needle) 1M & \textbf{16.4\%} & -- & -- & -- \\
    % ---- Image ----
    \textit{Image Und.} & MMMU \textit{(single)} & 82.0\% & \textbf{82.9\%} & 81.6\% & 74.4\% \\
    % ---- BLOK 2: Header ulang ----
    \textit{Capability} & \textit{Benchmark} & {Claude 4\\Opus} & {Grok 3 Beta\\Ext. Thinking} & \SetCell[c=2]{c} {DeepSeek R1\\0528} \\
    % ---- Code ----
    \textit{Code} & LiveCodeBench & 51.1\% & -- & \SetCell[c=2]{c} 70.5\% \\
    \textit{Code} & Aider Polyglot & 72.0\% & 53.3\% & \SetCell[c=2]{c} 71.6\% \\
    \textit{Code} & SWE-bench Verified \textit{(single)} & 72.5\% & -- & \SetCell[c=2]{c} -- \\
    \textit{Code} & SWE-bench Verified \textit{(multiple)} & 79.4\% & -- & \SetCell[c=2]{c} 57.6\% \\
    % ---- Reasoning ----
    \textit{Reasoning} & GPQA (diamond) \textit{(single)} & 79.6\% & 80.2\% & \SetCell[c=2]{c} 81.0\% \\
    \textit{Reasoning} & Humanity's Last Exam \textit{(no tools)} & 10.7\% & -- & \SetCell[c=2]{c} 14.0\% \\
    % ---- Factuality ----
    \textit{Factuality} & SimpleQA & -- & 43.6\% & \SetCell[c=2]{c} 27.8\% \\
    \textit{Factuality} & FACTS Grounding & 77.7\% & 74.8\% & \SetCell[c=2]{c} 82.4\% \\
    % ---- Math ----
    \textit{Math} & AIME 2025 \textit{(single)} & 75.5\% & 77.3\% & \SetCell[c=2]{c} 87.5\% \\
    % ---- Long-context ----
    \textit{Long-context} & LOFT (hard retrieval) $\le$128K & -- & 73.1\% & \SetCell[c=2]{c} -- \\
    \textit{Long-context} & LOFT (hard retrieval) 1M & -- & -- & \SetCell[c=2]{c} -- \\
    \textit{Long-context} & MRCR-V2 (8-needle) $\le$128K & 16.1\% & 34.0\% & \SetCell[c=2]{c} -- \\
    \textit{Long-context} & MRCR-V2 (8-needle) 1M & -- & -- & \SetCell[c=2]{c} -- \\
    % ---- Image ----
    \textit{Image Und.} & MMMU \textit{(single)} & 76.5\% & 76.0\% & \SetCell[c=2]{c} \textit{No MM} \\
\end{longtblr}
}

Selain varian Pro, \textcite{comaniciGemini25Pushing2025} menyoroti kemampuan Gemini 2.5 Flash yang menawarkan keunggulan dalam kecepatan pemrosesan. Karakteristik ini menjadikannya solusi yang sesuai untuk aplikasi yang membutuhkan respons instan atau pemrosesan volume data yang besar. Perbandingan kecepatan inferensi yang signifikan pada varian Flash ditunjukkan pada Gambar \ref{fig:gemini25_flash_speed}.

\begin{figure}[H]
\centering
\includegraphics[width=0.8\textwidth]{bab2/gemini2.5flashisfast.png}
\caption[Perbandingan kecepatan Gemini 2.5 Flash]{Perbandingan kecepatan Gemini 2.5 Flash \parencite{comaniciGemini25Pushing2025}}
\label{fig:gemini25_flash_speed}
\end{figure}

Selain keunggulan dalam kecepatan, \textcite{comaniciGemini25Pushing2025} juga mencatat bahwa Gemini 2.5 menunjukkan peningkatan performa yang signifikan. Peningkatan ini terlihat jelas jika dibandingkan dengan model-model pendahulunya. Perbandingan kemampuan model-model dalam keluarga Gemini hingga versi 2.5 diperlihatkan pada Tabel \ref{tab:gemini25_comparisons}.

{
\setstretch{1.0}
\begin{longtblr}[
    caption = {Perbandingan kemampuan keluarga model Gemini hingga versi 2.5 \parencite{comaniciGemini25Pushing2025}},
    entry = {Perbandingan kemampuan keluarga model Gemini},
    label = {tab:gemini25_comparisons},
]{
        width = \linewidth,
        colspec = {X[1.1,l] *{6}{X[c]}},
        hlines, vlines,
        rows = {m},
        rowsep = 4pt,
        row{1} = {font=\bfseries, halign=c, valign=m},
        column{1} = {font=\bfseries},
        rowhead = 1,
    }
    Fitur & {Gemini 1.5\\Flash} & {Gemini 1.5\\Pro} & {Gemini 2.0\\Flash-Lite} & {Gemini 2.0\\Flash} & {Gemini 2.5\\Flash} & {Gemini 2.5\\Pro} \\
    \textit{Modality} masukan & Teks, Citra, Video, Audio & Teks, Citra, Video, Audio & Teks, Citra, Video, Audio & Teks, Citra, Video, Audio & Teks, Citra, Video, Audio & Teks, Citra, Video, Audio \\
    Panjang masukan & 1M & 2M & 1M & 1M & 1M & 1M \\
    \textit{Modality} luaran & Teks & Teks & Teks & Teks, Citra* & Teks, Audio* & Teks, Audio* \\
    Panjang luaran & 8K & 8K & 8K & 8K & 64K & 64K \\
    \textit{Thinking} & Tidak & Tidak & Tidak & Ya* & Dinamis & Dinamis \\
    {Dukung\newline an} \textit{tool} & Tidak & Tidak & Tidak & Ya & Ya & Ya \\
    {Batas pengetahu\newline an} & November 2023 & November 2023 & Juni 2024 & Juni 2024 & Januari 2025 & Januari 2025 \\
\end{longtblr}
}

Iterasi terbaru yang saat ini diperkenalkan adalah seri Gemini 3, yang kembali menghadirkan varian Pro dan Flash dengan peningkatan pada arsitektur dasarnya. Gemini 3 Pro dirancang sebagai model dengan kemampuan penalaran terbaik di kelasnya, mampu menangani nuansa bahasa dan logika yang paling kompleks \parencite{google2025}. Peningkatan ini memungkinkan model untuk memproses instruksi yang rumit dengan akurasi yang lebih tinggi, khususnya pada domain yang memerlukan pemahaman konteks yang mendalam. Di sisi lain, Gemini 3 Flash menargetkan keseimbangan optimal antara kecerdasan, kecepatan, dan biaya operasional, serta kemampuan agentik yang lebih andal untuk eksekusi tugas otomatis \parencite{googleGemini3Flash2025}. Dengan latensi yang lebih rendah, varian Flash ini dioptimalkan untuk skenario penggunaan skala besar yang menuntut efisiensi waktu dan sumber daya komputasi. Performa kedua model ini dalam berbagai pengujian standar dapat dilihat pada Tabel \ref{tab:gemini3_pro_benchmark}.

{
\setstretch{1.0}
\begin{longtblr}[
    caption = {Hasil \textit{benchmark} Gemini 3 Pro pada berbagai tugas \parencite{google2025}},
    label = {tab:gemini3_pro_benchmark},
    entry = {Hasil \textit{benchmark} Gemini 3 Pro},
]{
    width = \linewidth,
    colspec = {X[2,l] X[1.5,l] *{4}{X[1,c]}},
    rows = {m},
    rowhead = 1,
    row{1} = {font=\bfseries, halign=c, valign=m},
    hlines, vlines,
    rowsep = 3pt,
}
    \textit{Benchmark} & Deskripsi & {Gemini 3\\Pro} & {Gemini 2.5\\Pro} & {Claude 4.5\\Sonnet} & {GPT-5.1} \\
    Humanity's Last Exam & \textit{Academic reasoning} & {\textbf{37.5\%} \scriptsize (no tools) \\ \textbf{45.8\%} \scriptsize (code)} & 21.6\% & 13.7\% & 26.5\% \\
    ARC-AGI-2 & \textit{Visual reasoning puzzles} & \textbf{31.1\%} & 4.9\% & 13.6\% & 17.6\% \\
    GPQA Diamond & \textit{Scientific knowledge} & \textbf{91.9\%} & 86.4\% & 83.4\% & 88.1\% \\
    AIME 2025 & \textit{Mathematics} & {95.0\% \scriptsize (no tools) \\ \textbf{100\%} \scriptsize (code)} & 88.0\% & {87.0\% \\ \textbf{100\%}} & 94.0\% \\
    MathArena Apex & \textit{Math Contest problems} & \textbf{23.4\%} & 0.5\% & 1.6\% & 1.0\% \\
    MMMU-Pro & \textit{Multimodal reasoning} & \textbf{81.0\%} & 68.0\% & 68.0\% & 76.0\% \\
    ScreenSpot-Pro & \textit{Screen understanding} & \textbf{72.7\%} & 11.4\% & 36.2\% & 3.5\% \\
    CharXiv Reasoning & \textit{Chart information synthesis} & \textbf{81.4\%} & 69.6\% & 68.5\% & 69.5\% \\
    OmniDocBench 1.5 & \textit{OCR (lower is better)} & \textbf{0.115} & 0.145 & 0.145 & 0.147 \\
    Video-MMMU & \textit{Video knowledge} & \textbf{87.6\%} & 83.6\% & 77.8\% & 80.4\% \\
    LiveCodeBench Pro & \textit{Competitive coding} & \textbf{2,439} & 1,775 & 1,418 & 2,243 \\
    Terminal-Bench 2.0 & \textit{Agentic terminal coding} & \textbf{54.2\%} & 32.6\% & 42.8\% & 47.6\% \\
    SWE-Bench Verified & \textit{Agentic coding} & 76.2\% & 59.6\% & \textbf{77.2\%} & 76.3\% \\
    $\tau$2-bench & \textit{Agentic tool use} & \textbf{85.4\%} & 54.9\% & 84.7\% & 80.2\% \\
    Vending-Bench 2 & \textit{Long-horizon agentic tasks} & \textbf{\$5,478} & \$574 & \$3,839 & \$1,473 \\
    FACTS \textit{Benchmark} & \textit{Internal grounding} & \textbf{70.5\%} & 63.4\% & 50.4\% & 50.8\% \\
    SimpleQA Verified & \textit{Parametric knowledge} & \textbf{72.1\%} & 54.5\% & 29.3\% & 34.9\% \\
    MMMLU & \textit{Multilingual Q\&A} & \textbf{91.8\%} & 89.5\% & 89.1\% & 91.0\% \\
    Global PIQA & \textit{Commonsense reasoning} & \textbf{93.4\%} & 91.5\% & 90.1\% & 90.9\% \\
    MRCR v2 (8-needle) & \textit{Long context} & {\textbf{77.0\%} \scriptsize (128k) \\ \textbf{26.3\%} \scriptsize (1M)} & {58.0\% \\ 16.4\%} & {47.1\% \\ --} & {61.6\% \\ --} \\
\end{longtblr}
}

Melengkapi keluarga model generatif tersebut, \textcite{leeGeminiEmbeddingGeneralizable2025} juga menyediakan ekosistem Gemini dengan model \textit{embedding} yang dioptimalkan untuk tugas pencarian informasi semantik dan klasterisasi data. Model ini berfungsi memetakan teks ke dalam vektor representasi berdimensi tinggi yang menangkap makna semantik, yang merupakan komponen krusial dalam implementasi sistem \textit{Retrieval-Augmented Generation} (RAG). Efektivitas model \textit{embedding} ini dalam menangkap konteks semantik telah dibuktikan melalui berbagai pengujian \textit{benchmark}, sebagaimana ditampilkan pada Tabel \ref{tab:gemini_embedding_benchmark}.

{
\setstretch{1.0}
\begin{longtblr}[
    caption = {Hasil \textit{benchmark} model \textit{embedding} Gemini \parencite{leeGeminiEmbeddingGeneralizable2025}},
    entry = {Hasil \textit{benchmark} model \textit{embedding} Gemini},
    label = {tab:gemini_embedding_benchmark},
]{
    width = \linewidth,
    colspec = {X[1.6,l] *{6}{X[1,c]}},
    rows = {m},
    rowhead = 1,
    row{1} = {font=\bfseries, halign=c, valign=m},
    cell{1}{2-7} = {font=\scriptsize\bfseries},
    hlines, vlines,
    rowsep = 3pt,
}
    Task & {Gemini\\Embedding} & {Gecko\\Embedding$^\dagger$} & {gte-Qwen2\\7B-instruct} & {multilingual-e5\\large-instruct} & {Cohere-embed\\multilingual-v3.0} & {text-embedding\\3-large} \\
    \SetCell[c=7]{l} \textbf{MTEB(Multilingual)} & & & & & & \\
    \hspace{3mm} Mean (Task) & \textbf{68.32} & 62.13 & 62.51 & 63.23 & 61.10 & 58.92 \\
    \hspace{3mm} Mean (Type) & \textbf{59.64} & 54.32 & 56.00 & 55.17 & 53.31 & 51.48 \\
    \hspace{3mm} - Bitext Mining & 79.32 & 70.73 & 73.92 & \textbf{80.13} & 70.50 & 62.17 \\
    \hspace{3mm} - Classification & \textbf{71.84} & 64.64 & 61.55 & 64.94 & 62.95 & 60.27 \\
    \hspace{3mm} - Clustering & \textbf{54.99} & 48.47 & 53.36 & 51.54 & 47.61 & 47.49 \\
    \hspace{3mm} - Inst. Retrieval & \textbf{5.18} & 4.08 & 4.94 & -0.40 & -1.89 & -2.68 \\
    \hspace{3mm} - Multilabel Class. & \textbf{29.16} & 22.80 & 25.48 & 22.91 & 22.74 & 22.03 \\
    \hspace{3mm} - Pair Class. & 83.64 & 81.14 & \textbf{85.13} & 80.86 & 79.88 & 79.17 \\
    \hspace{3mm} - Reranking & \textbf{65.72} & 61.22 & 65.55 & 62.61 & 64.07 & 63.89 \\
    \hspace{3mm} - Retrieval & \textbf{67.71} & 59.68 & 60.08 & 57.12 & 59.16 & 59.27 \\
    \hspace{3mm} - STS & \textbf{79.40} & 76.11 & 73.98 & 76.81 & 74.80 & 71.68 \\
    \SetCell[c=7]{l} \textbf{MTEB(Eng, v2)} & & & & & & \\
    \hspace{3mm} Mean (Task) & \textbf{73.30} & 69.53 & 70.72 & 65.53 & 66.01 & 66.43 \\
    \hspace{3mm} Mean (Type) & \textbf{67.67} & 64.82 & 65.77 & 61.21 & 61.43 & 62.15 \\
    {\textbf{MTEB} \\ \textbf{(Code)$^*$} \\} & \textbf{74.66} & 65.40 & 56.41 & 57.94 & 51.94 & 58.95 \\
    {\textbf{XOR-Retrieve} \\} & \textbf{90.42} & 65.67 & N/A & N/A & N/A & 68.76 \\
    {\textbf{XTREME-UP} \\} & \textbf{64.33} & 34.97 & 17.39 & 18.68 & N/A & 18.80 \\
    \textbf{Commercial Use} & \checkmark & \checkmark & & & \checkmark & \checkmark \\
\end{longtblr}
} 

\section{\textit{Prompt Engineering}}

LLM yang telah melalui tahap \textit{fine-tuning} memiliki kemampuan yang sangat baik dalam mengikuti instruksi, menjawab pertanyaan, dan melakukan percakapan melalui mekanisme pemberian \textit{prompt} \parencite{jm3}. \textit{Prompt} didefinisikan sebagai string teks yang diberikan oleh pengguna kepada model bahasa untuk mendapatkan luaran yang bermanfaat. Dalam prosesnya, string \textit{prompt} tersebut diteruskan ke model, yang kemudian secara iteratif membangkitkan token berdasarkan kondisi yang ditetapkan dalam \textit{prompt}. Upaya sistematis untuk menemukan formulasi \textit{prompt} yang paling efektif dalam menyelesaikan suatu tugas tertentu dikenal sebagai \textit{prompt engineering}. 

\textit{Prompt engineering} berakar pada kemampuan \textit{in-context learning} LLM, yang memungkinkan model beradaptasi terhadap tugas baru melalui instruksi atau contoh dalam konteks inferensi tanpa pembaruan parameter \parencite{brownLanguageModelsAre2020}. \textcite{huyenAIEngineeringBuilding2025} mendeskripsikan mekanisme ini sebagai bentuk pembelajaran berkelanjutan (\textit{continual learning}) yang memungkinkan model untuk mengasimilasi informasi baru secara dinamis, sehingga tetap relevan tanpa perlu dilatih ulang. Dalam konteks ini, setiap contoh yang disertakan disebut sebagai \textit{shot}. Pendekatan dengan memberikan sejumlah contoh dikenal sebagai \textit{few-shot learning}, sedangkan tanpa contoh disebut \textit{zero-shot learning}. Secara spesifik dalam konteks \textit{prompting}, tugas yang melibatkan pemberian contoh ini sering disebut sebagai \textit{few-shot prompting}, berbeda dengan \textit{zero-shot prompting} yang merupakan instruksi tanpa disertai contoh berlabel \parencite{jm3}. Meskipun penambahan contoh (\textit{shots}) umumnya meningkatkan performa, jumlahnya dibatasi oleh kapasitas jendela konteks (\textit{context window})  model dan perlu dipertimbangkan terhadap peningkatan biaya inferensi yang ditimbulkannya \parencite{huyenAIEngineeringBuilding2025}.

Selain instruksi yang diberikan dalam \textit{prompt}, LLM umumnya memiliki komponen yang disebut \textit{system prompt}, yaitu instruksi teks tunggal yang diberikan di awal interaksi untuk menetapkan peran, nada, dan konteks model \parencite{jm3}. \textit{System prompt} ini secara otomatis disisipkan (\textit{silently prepended}) sebelum teks masukan pengguna, sehingga menjadi acuan utama bagi model dalam memproses keseluruhan konteks percakapan. Keberadaan jendela konteks yang sangat besar pada model bahasa modern memungkinkan penggunaan \textit{system prompt} yang panjang, yang memberikan fleksibilitas lebih besar dalam mengarahkan perilaku model untuk tugas-tugas spesifik. Contoh ilustrasi mengenai struktur dan penerapan \textit{system prompt} dalam interaksi model bahasa ditampilkan pada Gambar \ref{fig:system_prompt}.
\begin{figure}[H]
\centering
\includegraphics[width=0.7\textwidth]{bab2/exsystemprompt.png}
\caption[Contoh \textit{System Prompt}]{Contoh penggunaan \textit{System Prompt} dalam interaksi dengan model bahasa \parencite{jm3}}
\label{fig:system_prompt}
\end{figure}

\subsection{\textit{Chain-of-Thought}}
\textit{Chain-of-Thought} (CoT) \textit{prompting} merupakan teknik yang diperkenalkan oleh \textcite{weiChainofThoughtPromptingElicits2023} untuk meningkatkan kemampuan penalaran kompleks pada \textit{Large Language Models} (LLM). Metode ini mengarahkan model untuk menghasilkan serangkaian langkah penalaran perantara (\textit{intermediate reasoning steps}) sebelum menyajikan jawaban akhir, menyerupai proses kognitif manusia dalam memecahkan masalah yang rumit secara bertahap.

Dalam paradigma \textit{prompting} standar, model umumnya dilatih untuk memprediksi jawaban secara langsung dari masukan yang diberikan \parencite{weiChainofThoughtPromptingElicits2023}. Pendekatan tersebut sering kali menghasilkan kesalahan pada tugas-tugas yang memerlukan logika bertingkat, seperti soal cerita matematika atau penalaran simbolis. Sebaliknya, CoT memungkinkan model untuk menguraikan masalah menjadi komponen-komponen yang lebih kecil, sehingga mengurangi kemungkinan kesalahan logika pada hasil akhir. Perbedaan fundamental antara pendekatan standar dan CoT divisualisasikan pada Gambar \ref{fig:cot_comparison}.

\begin{figure}[H]
\centering
\includegraphics[width=0.6\textwidth]{bab2/cotcomparison.png}
\caption[Perbandingan \textit{standard prompting} dan \textit{chain-of-thought prompting}]{Perbandingan \textit{standard prompting} dan \textit{chain-of-thought prompting} dalam menyelesaikan soal cerita matematika \parencite{weiChainofThoughtPromptingElicits2023}}
\label{fig:cot_comparison}
\end{figure}

Penerapan CoT dapat dilakukan melalui dua strategi utama. Strategi pertama adalah \textit{Few-Shot CoT}, di mana \textit{prompt} diperkaya dengan sejumlah contoh demonstrasi (\textit{exemplars}) yang secara eksplisit menampilkan alur penalaran beserta jawabannya \parencite{weiChainofThoughtPromptingElicits2023}. Strategi kedua adalah \textit{Zero-Shot CoT} yang diusulkan oleh \textcite{kojimaLargeLanguageModels2023}. Pada strategi ini, kemampuan penalaran model diaktifkan hanya dengan menambahkan instruksi pemicu seperti "Mari berpikir langkah demi langkah" (\textit{"Let's think step by step"}) tanpa memerlukan contoh demonstrasi spesifik. Temuan ini mengindikasikan bahwa kapasitas penalaran merupakan kemampuan inheren yang dimiliki oleh LLM berskala besar, yang dapat dimunculkan melalui instruksi yang tepat.

\subsection{\textit{Reasoning and Acting}}
\textit{Reasoning and Acting} (ReAct) merupakan pola yang mengatur cara kerja \textit{Large Language Model} (LLM). Pola ini pertama kali diperkenalkan oleh \textcite{yaoReActSynergizingReasoning2023} sebagai sebuah teknik \textit{prompting} dalam \textit{prompt engineering}. ReAct memberikan \textit{prompt} yang mengarahkan LLM untuk menghasilkan jejak pemikiran sekaligus melaksanakan tindakan yang spesifik terhadap suatu tugas secara bersamaan. Melalui mekanisme tersebut, sistem mampu melakukan penalaran dinamis untuk menyusun, menjaga, dan menyesuaikan rencana tindakan, sekaligus berinteraksi dengan lingkungan eksternal, misalnya dengan melakukan pencarian tambahan melalui API atau basis data sebelum menggunakan hasilnya untuk mengambil keputusan yang lebih tepat. 

Dalam implementasinya, \textcite{yaoReActSynergizingReasoning2023} menggunakan \textit{prompt} untuk menginstruksikan LLM menghasilkan tindakan dalam skema atau struktur yang telah ditentukan, yang salah satunya dapat berupa pemanggilan (API). Dari struktur ini, lingkungan eksternal bertugas mengekstraksi pemanggilan API tersebut untuk dieksekusi pada kode alat yang sebenarnya. Gambar \ref{fig:react_illustration} memperlihatkan ilustrasi mekanisme ReAct dalam menjawab pertanyaan pada dataset HotpotQA dari \textcite{yangHotpotQADatasetDiverse2018}. Dalam ilustrasi tersebut, model berusaha mengidentifikasi perangkat lain yang dapat digunakan untuk mengontrol program menggunakan Apple Remote.

\begin{figure}[H]
\centering
\includegraphics[width=0.7\textwidth]{bab2/ilustrasireact.png}
\caption[Ilustrasi metode ReAct]{Ilustrasi metode ReAct dalam menyelesaikan tugas \parencite{yaoReActSynergizingReasoning2023}}
\label{fig:react_illustration}
\end{figure}

ReAct merupakan pengembangan dari teknik \textit{prompting} \textit{Chain-of-Thought} (CoT) dari \textcite{weiChainofThoughtPromptingElicits2023}, yang memungkinkan model menghasilkan jejak penalaran (\textit{reasoning traces}) untuk menjawab pertanyaan yang melibatkan aritmatika, penalaran berbasis logika, maupun \textit{commonsense reasoning}. Akan tetapi, CoT memiliki keterbatasan karena tidak dapat mengakses informasi eksternal atau memperbarui pengetahuannya, sehingga sering menimbulkan masalah seperti halusinasi fakta dan propagasi kesalahan. Pada penelitian \textcite{yaoReActSynergizingReasoning2023}, ReAct dan CoT dievaluasi pada dataset HotpotQA \parencite{yangHotpotQADatasetDiverse2018} untuk tugas tanya jawab dan FEVER \parencite{thorneFEVERLargescaleDataset2018} untuk verifikasi fakta, dengan menggunakan PaLM-540B sebagai model dasar. Hasil evaluasi menunjukkan bahwa ReAct secara umum memberikan kinerja lebih baik dibandingkan Act (yang hanya melibatkan tindakan) pada kedua tugas, dengan skor 27,4 pada HotpotQA dan 58,9 pada FEVER untuk ReAct, serta 25,7 dan 58,9 untuk Act. ReAct juga melampaui CoT pada FEVER dengan skor 60,9, meskipun masih tertinggal pada HotpotQA dengan skor 27,4 dibandingkan 33,4 untuk CoT.

Analisis lebih lanjut oleh \textcite{yaoReActSynergizingReasoning2023} menunjukkan bahwa CoT rentan terhadap halusinasi fakta, sementara ReAct memiliki keterbatasan fleksibilitas karena adanya struktur langkah yang kaku dalam merumuskan penalaran. Selain itu, ReAct sangat bergantung pada kualitas informasi yang diambil. Jika hasil pencarian tidak relevan, proses penalaran model dapat terganggu dan sulit untuk dipulihkan. Menariknya, ketika CoT dan ReAct dikombinasikan melalui mekanisme yang memungkinkan model beralih di antara keduanya, kinerjanya melampaui semua metode \textit{prompting} lain, dengan skor 35,1 pada HotpotQA dan 64,6 pada FEVER. Hasil ini memperlihatkan bahwa kekurangan ReAct dapat ditutupi oleh keunggulan CoT, sehingga kombinasi keduanya menghasilkan sistem yang lebih akurat.

\section{\textit{Object Detection}}
Deteksi objek adalah salah satu tugas fundamental dalam \textit{computer vision} yang bertujuan untuk mengidentifikasi keberadaan objek tertentu dalam citra digital dan menentukan lokasi tepatnya \parencite{raniObjectDetectionRecognition2022}. Berbeda dengan pengenalan gambar (\textit{image recognition}) atau klasifikasi citra yang hanya memberikan label kategori tunggal untuk keseluruhan gambar, deteksi objek memberikan informasi yang lebih granular berupa label kelas dan koordinat kotak pembatas untuk setiap instansi objek yang ditemukan. Perbedaan visual antara pengenalan gambar dan deteksi objek dapat dilihat pada Gambar \ref{fig:obj_vs_recog}.

\begin{figure}[H]
\centering
\includegraphics[width=0.4\textwidth]{bab2/Image-recognition-vs-object-detection.png}
\caption[Perbedaan Antara \textit{Image Recognition} dan \textit{Object Detection}]{Perbedaan Antara \textit{Image Recognition} dan \textit{Object Detection} \parencite{raniObjectDetectionRecognition2022}}
\label{fig:obj_vs_recog}
\end{figure}

Berdasarkan cakupan kategori objek yang dapat dikenali, sistem deteksi objek secara umum diklasifikasikan menjadi dua paradigma, yaitu \textit{closed-set object detection} (CSOD) dan \textit{open-vocabulary object detection} (OVD) \parencite{hosoyaOpenvocabularyVsClosedset2024}. CSOD merupakan pendekatan konvensional di mana model dilatih menggunakan dataset dengan daftar kategori yang telah ditentukan sebelumnya . Model jenis ini hanya mampu mendeteksi objek yang termasuk dalam kategori pelatihan tersebut dan akan memperlakukan objek lain sebagai latar belakang. Sebaliknya, OVD bertujuan untuk mengatasi keterbatasan ini dengan memanfaatkan kemampuan generalisasi dari model bahasa-visi. Pendekatan ini memungkinkan model untuk mendeteksi objek berdasarkan deskripsi teks bebas, termasuk kategori objek yang tidak pernah dilihat selama proses pelatihan. Perbedaan mekanisme antara kedua pendekatan ini diilustrasikan pada Gambar \ref{fig:closed_vs_open}.

\begin{figure}[H]
\centering
\includegraphics[width=\textwidth]{bab2/closedsetvsopenvocab.png}
\caption[Perbedaan \textit{closed-set} dan \textit{open-vocabulary object detection}]{Perbedaan \textit{closed-set} dan \textit{open-vocabulary object detection} \parencite{hosoyaOpenvocabularyVsClosedset2024}}
\label{fig:closed_vs_open}
\end{figure}

\subsection{\textit{You Only Look Once}}
\textit{You Only Look Once} (YOLO) merupakan algoritma deteksi objek yang pertama kali diperkenalkan oleh \textcite{redmonYouOnlyLook2016a} melalui pendekatan yang menyatukan proses prediksi lokasi objek dan klasifikasi dalam satu jaringan saraf tunggal. Dalam bidang \textit{computer vision}, deteksi objek merupakan tugas krusial yang bertujuan untuk mengidentifikasi jenis objek sekaligus menentukan posisinya dalam citra melalui koordinat \textit{bounding box}. Sebelum kemunculan YOLO, deteksi objek didominasi oleh pendekatan \textit{two-stage detector} seperti Faster R-CNN yang dikembangkan oleh \textcite{renFasterRCNNRealTime2017}. Visualisasi arsitektur dasar YOLO diilustrasikan pada Gambar \ref{fig:yolo_architecture}.

\begin{figure}[H]
\centering
\includegraphics[width=0.8\textwidth]{bab2/arsitekturyolo.png}
\caption[Arsitektur YOLO]{Arsitektur YOLO \parencite{redmonYouOnlyLook2016a}}
\label{fig:yolo_architecture}
\end{figure}

Perbedaan utama antara kedua pendekatan tersebut terletak pada alur pemrosesannya. \textit{Two-stage detector} bekerja dengan menghasilkan sejumlah usulan wilayah (\textit{region proposals}) terlebih dahulu, baru kemudian melakukan klasifikasi pada setiap wilayah tersebut secara terpisah \Parencite{renFasterRCNNRealTime2017}. Meskipun akurat, proses ini memerlukan waktu komputasi yang besar sehingga sulit digunakan untuk aplikasi yang membutuhkan respons cepat. Sebaliknya, YOLO sebagai \textit{one-stage detector} memperlakukan deteksi objek sebagai masalah regresi tunggal, yang memetakan piksel gambar langsung ke koordinat \textit{bounding box} dan probabilitas kelas dalam satu kali jalan (\textit{one-pass}) \parencite{redmonYouOnlyLook2016a}. Karakteristik ini memungkinkan YOLO untuk bekerja secara \textit{real-time} dengan kecepatan yang jauh mengungguli pendahulunya.

Perkembangan arsitektur YOLO terus mengalami peningkatan signifikan sejak pertama kali dirilis. Setelah versi awal oleh \textcite{redmonYouOnlyLook2016a}, berbagai iterasi seperti YOLOv2 \parencite{redmonYOLO9000BetterFaster2016}, YOLOv3 \parencite{redmonYOLOv3IncrementalImprovement2018}, hingga YOLOv7 \parencite{wangYOLOv7TrainableBagofFreebies2023} membawa peningkatan pada akurasi dan kecepatan melalui teknik-teknik seperti \textit{anchor boxes}, \textit{feature pyramid networks}, dan optimalisasi \textit{backbone}. Saat ini, pengembangan YOLO dikelola secara luas oleh Ultralytics \parencite{jocherUltralyticsYOLO2023}, yang telah merilis versi populer seperti YOLOv8 dan YOLOv9 \parencite{wangYOLOv9LearningWhat2025}. Iterasi terbaru ini mengintegrasikan kemajuan terbaru dalam \textit{deep learning}, seperti penggunaan \textit{backbone} yang lebih efisien serta modul \textit{neck} yang mampu menangkap fitur pada berbagai skala dengan lebih baik.

Versi terbaru yang diperkenalkan oleh Ultralytics adalah YOLOv11, yang membawa sejumlah pembaruan arsitektur untuk meningkatkan efisiensi ekstraksi fitur dan performa deteksi \parencite{khanamYOLOv11OverviewKey2024}. Arsitektur YOLOv11 memperkenalkan blok C3k2 yang merupakan evolusi dari modul C2f, serta penggunaan \textit{Spatial Pyramid Pooling - Fast} (SPPF) yang dioptimalkan untuk mempercepat pemrosesan tanpa mengurangi kualitas representasi spasial. Selain itu, integrasi modul \textit{Cross-Stage Partial with Shared Attention} (C2PSA) pada bagian \textit{neck} memungkinkan model untuk lebih fokus pada fitur-fitur penting dalam citra, sehingga meningkatkan akurasi pada objek dengan skala yang bervariasi. Visualisasi arsitektur YOLOv11 secara mendetail dapat dilihat pada Gambar \ref{fig:yolov11_architecture}.

\begin{figure}[H]
\centering
\includegraphics[width=0.9\textwidth]{bab2/yolov11architecture.png}
\caption[Arsitektur YOLOv11]{Arsitektur YOLOv11 \parencite{sapkotaYOLOAdvancesIts2025}}
\label{fig:yolov11_architecture}
\end{figure}

Dari sisi performa, YOLOv11 menunjukkan peningkatan efisiensi komputasi yang signifikan dibandingkan dengan versi sebelumnya. Berdasarkan hasil pengujian pada dataset COCO, YOLOv11 mampu mencapai nilai \textit{mean Average Precision} (mAP) yang lebih tinggi dengan jumlah parameter yang lebih sedikit jika dibandingkan dengan YOLOv8 dan YOLOv10 \parencite{khanamYOLOv11OverviewKey2024}. Hal ini membuat YOLOv11 sangat kompetitif untuk aplikasi \textit{real-time} yang membutuhkan keseimbangan antara akurasi tinggi dan latensi rendah. Perbandingan performa YOLOv11 terhadap iterasi YOLO sebelumnya diilustrasikan pada Gambar \ref{fig:yolov11_benchmark}.

\begin{figure}[H]
\centering
\includegraphics[width=0.8\textwidth]{bab2/yolov11benchmark.png}
\caption[Perbandingan performa YOLOv11]{Perbandingan performa YOLOv11 terhadap model YOLO lainnya \parencite{khanamYOLOv11OverviewKey2024}}
\label{fig:yolov11_benchmark}
\end{figure}

Saat ini, Ultralytics memposisikan YOLO11 sebagai versi yang paling stabil dan efisien untuk berbagai skenario penggunaan \parencite{jocherUltralyticsYOLO2023}, meskipun versi YOLOv12 telah mulai diperkenalkan oleh \textcite{tianYOLOv12AttentionCentricRealTime2025}. Peningkatan arsitektur pada YOLO11 difokuskan pada pengoptimalan aliran informasi dalam jaringan guna mengurangi beban komputasi tanpa mengorbankan akurasi. Hal ini dapat dibuktikkan berdasarkan hasil evaluasi kinerja yang dilakukan oleh \textcite{sapkotaComprehensivePerformanceEvaluation2026} pada deteksi objek di lingkungan pertanian yang kompleks, YOLO11 menunjukkan konsistensi performa yang lebih baik dibandingkan dengan varian YOLOv8 hingga YOLOv12. Dalam pengujian tersebut, varian terkecil yakni YOLO11n (\textit{nano}) tercatat memiliki kecepatan pemrosesan paling tinggi dibandingkan model lainnya, yang menjadikannya sangat ideal untuk implementasi pada perangkat dengan keterbatasan sumber daya. Hasil perbandingan kinerja dari berbagai varian YOLO tersebut disajikan pada Tabel \ref{tab:yolo_benchmark}.

{
\setstretch{1.0}
\begin{longtblr}[
    caption = {Hasil \textit{benchmark} berbagai varian YOLO pada tugas deteksi objek di lingkungan pertanian oleh \textcite{sapkotaComprehensivePerformanceEvaluation2026}},
    entry = {Hasil \textit{benchmark} berbagai varian YOLO},
    label = {tab:yolo_benchmark},
]{
    width = \linewidth,
    colspec = {X[1.5,l] *{3}{X[1,c]}},
    hlines, vlines,
    rows = {m},
    rowsep = 4pt,
    row{1} = {font=\bfseries, halign=c, valign=m},
    rowhead = 1,
}
    Model & Pre (ms) & Inf (ms) & Post (ms) \\
    \SetCell[c=4]{l} \textit{YOLOv8} & & & \\
    YOLOv8n & 1.3 & 4.1 & 2.3 \\
    YOLOv8s & 1.3 & 6.4 & 2.3 \\
    YOLOv8m & 1.3 & 11.2 & 2.1 \\
    YOLOv8l & 1.2 & 18.7 & 2.2 \\
    YOLOv8x & 0.9 & 24.8 & 2.3 \\
    \SetCell[c=4]{l} \textit{YOLOv9} & & & \\
    YOLOv9-t & 1.3 & 14.1 & 2.2 \\
    YOLOv9-s & 1.3 & 11.5 & 2.2 \\
    YOLOv9-m & 1.3 & 14.0 & 2.1 \\
    YOLOv9-c & 1.3 & 17.0 & 2.0 \\
    YOLOv9-b & 1.2 & 17.2 & 2.0 \\
    YOLOv9-e & 1.1 & 33.5 & 1.9 \\
    \SetCell[c=4]{l} \textit{YOLOv10} & & & \\
    YOLOv10n & 1.4 & 5.5 & 1.6 \\
    YOLOv10s & 1.3 & 7.7 & 1.6 \\
    YOLOv10m & 1.3 & 13.0 & 1.6 \\
    YOLOv10b & 1.3 & 16.7 & 1.5 \\
    YOLOv10l & 1.4 & 19.6 & 1.5 \\
    YOLOv10x & 1.2 & 26.5 & 1.5 \\
    \SetCell[c=4]{l} \textit{YOLOv11} & & & \\
    YOLOv11n & \textbf{0.1} & \textbf{2.4} & 2.2 \\
    YOLOv11s & \textbf{0.1} & 5.0 & 2.5 \\
    YOLOv11m & \textbf{0.1} & 11.9 & 0.6 \\
    YOLOv11l & 0.2 & 14.6 & 1.2 \\
    YOLOv11x & \textbf{0.1} & 26.3 & 0.6 \\
    \SetCell[c=4]{l} \textit{YOLOv12} & & & \\
    YOLOv12n & 0.8 & 4.6 & 0.7 \\
    YOLOv12s & 3.8 & 9.6 & 1.0 \\
    YOLOv12m & 0.3 & 15.9 & \textbf{0.5} \\
    YOLOv12l & 0.6 & 19.2 & \textbf{0.5} \\
\end{longtblr}
}
%NOTE: Perlu tambahin studi literatur penggunaan model YOLO biasanya sebagai bagian pertama dari two-stage detector

\subsection{OWLv2}
OWLv2 (\textit{Open-Vocabulary World Object Detection}) merupakan model OVD yang dikembangkan oleh \textcite{mindererScalingOpenVocabularyObject2024} untuk mendeteksi dan mengklasifikasikan objek berdasarkan deskripsi teks bebas. Kemampuan ini memberikan fleksibilitas tinggi dalam mengenali objek baru tanpa memerlukan pelatihan ulang. Pengembangan model ini berakar pada penelitian sebelumnya oleh \textcite{mindererSimpleOpenVocabularyObject2022} yang memperkenalkan OWL-ViT (\textit{Vision Transformer for Open-World Localization}). OWL-ViT mengadaptasi arsitektur \textit{Vision Transformer} (ViT) standar dengan sedikit modifikasi untuk tugas deteksi objek dan dilatih menggunakan pasangan gambar serta teks dalam skala besar. Pendekatan ini bertujuan mentransfer kemampuan pengenalan visual-semantik dari model seperti CLIP ke dalam tugas lokalisasi objek, sehingga menjembatani kesenjangan antara klasifikasi gambar tingkat tinggi dan deteksi objek yang lebih rinci melalui pencocokan fitur dari \textit{image encoder} dengan \textit{text embedding} dari \textit{query} label.

\textcite{mindererScalingOpenVocabularyObject2024} kemudian mengembangkan OWLv2 dengan fokus utama pada skalabilitas data pelatihan. Tantangan utama dalam OVD adalah kelangkaan data anotasi kotak pembatas (\textit{bounding box}) yang terkurasi. Untuk mengatasi hal ini, OWLv2 menerapkan strategi pelatihan mandiri (\textit{self-training}) yang memanfaatkan detektor yang sudah ada untuk menghasilkan label semu (\textit{pseudo-labels}) pada dataset gambar-teks berskala web (WebLI) yang sangat besar. Dengan cara ini, model dapat belajar dari miliaran contoh data yang bising namun kaya informasi, yang secara signifikan meningkatkan kemampuan deteksi \textit{zero-shot}.

Secara arsitektur, \textcite{mindererScalingOpenVocabularyObject2024} mengoptimalkan desain OWL-ViT pada OWLv2 dengan meningkatkan efisiensi pada komponen \textit{detection head} dan strategi pelatihan. Model ini menggunakan \textit{image encoder} untuk mengekstraksi representasi visual dan \textit{text encoder} untuk memproses \textit{query} teks menjadi \textit{embedding}. Interaksi antara kedua representasi ini memungkinkan model untuk memprediksi lokasi objek dan probabilitas kelasnya berdasarkan kesesuaian semantik. Arsitektur umum dari pendekatan ini diilustrasikan pada Gambar \ref{fig:owlv2_architecture}.

\begin{figure}[H]
\centering
\includegraphics[width=\textwidth]{bab2/arsitekturowlv2.png}
\caption[Arsitektur OWLv2]{Arsitektur OWLv2 yang menggunakan \textit{image encoder} dan \textit{text encoder} untuk OVD \parencite{mindererScalingOpenVocabularyObject2024}}
\label{fig:owlv2_architecture}
\end{figure}

Keunggulan performa OWLv2 dibuktikan melalui evaluasi komprehensif oleh \textcite{wuEvaluatingPerformanceOpenVocabulary2026} yang membandingkan berbagai model OVD terkini. Hasil studi tersebut menunjukkan bahwa OWLv2 secara konsisten mengungguli model serupa, seperti OWL-ViT dan Grounding DINO dalam hal akurasi deteksi dan kemampuan generalisasi pada kategori objek yang belum pernah dilihat sebelumnya. \citeauthor{wuEvaluatingPerformanceOpenVocabulary2026} menyoroti bahwa strategi skala besar yang diterapkan pada OWLv2 memberikan keuntungan signifikan dalam menangani variasi visual dan kompleksitas deskripsi teks dibandingkan pendekatan lainnya.


\section{\textit{Retrieval-Augmented Generation}}
\textit{Retrieval-Augmented Generation} (RAG) didefinisikan oleh \textcite{huyenAIEngineeringBuilding2025} sebagai teknik yang meningkatkan kemampuan generasi model melalui pengambilan informasi relevan dari sumber memori eksternal, seperti basis data internal atau internet. Pola pengambilan lalu generasi ini pertama kali diperkenalkan oleh \textcite{chenReadingWikipediaAnswer2017} dalam sistem yang memanfaatkan halaman Wikipedia untuk menjawab pertanyaan domain terbuka. Istilah RAG kemudian dikukuhkan oleh \textcite{lewisRetrievalaugmentedGenerationKnowledgeintensive2020} sebagai solusi untuk tugas-tugas intensif pengetahuan (\textit{knowledge-intensive tasks}) di mana kapasitas model tidak mencukupi untuk menampung seluruh informasi. Melalui pendekatan ini, hanya informasi yang paling relevan yang diambil dan dimasukkan ke dalam model, yang terbukti mampu menghasilkan respons yang lebih detail serta mengurangi risiko halusinasi.

Secara arsitektural, sistem RAG mengintegrasikan tiga komponen utama, yaitu \textit{Retrieval}, \textit{Augmentation}, dan \textit{Generation} \parencite{singhAgenticRetrievalAugmentedGeneration2025}. Komponen \textit{Retrieval} bertugas melakukan \textit{query}, yakni permintaan pencarian informasi, pada sumber data eksternal, seperti basis pengetahuan, API, atau basis data vektor. Komponen ini sering memanfaatkan pencarian vektor padat dan model berbasis \textit{Transformer} untuk meningkatkan presisi pengambilan serta relevansi semantik. Selanjutnya, komponen \textit{Augmentation} memproses data yang diambil dengan mengekstraksi dan meringkas informasi yang paling relevan agar selaras dengan konteks \textit{query}. Terakhir, komponen \textit{Generation} menggabungkan informasi tersebut dengan pengetahuan pra-latih LLM untuk menghasilkan respons yang koheren dan sesuai konteks. Komponen-komponen utama dalam arsitektur RAG ini diilustrasikan pada Gambar \ref{fig:rag_components}.

\begin{figure}[H]
\centering
\includegraphics[width=0.8\textwidth]{bab2/rag_component.png}
\caption[Komponen inti arsitektur RAG]{Komponen inti arsitektur RAG \parencite{singhAgenticRetrievalAugmentedGeneration2025}}
\label{fig:rag_components}
\end{figure}

\textcite{huyenAIEngineeringBuilding2025} mengonseptualisasikan RAG sebagai mekanisme konstruksi konteks yang spesifik untuk setiap \textit{query}. Mekanisme ini membedakan RAG dari penggunaan konteks statis yang bersifat seragam. Walaupun jendela konteks model terus berkembang, RAG tetap relevan karena kemampuannya dalam menyeleksi informasi esensial secara efisien. Hal tersebut berdampak pada pengurangan biaya komputasi dan latensi dibandingkan dengan pemrosesan konteks yang sangat panjang secara keseluruhan.

\textcite{singhAgenticRetrievalAugmentedGeneration2025} dan \textcite{huyenAIEngineeringBuilding2025} mengemukakan bahwa paradigma pengembangan RAG dapat bervariasi bergantung pada kebutuhan spesifik aplikasi yang dibangun. Secara garis besar, \textcite{huyenAIEngineeringBuilding2025} mengklasifikasikan kerangka kerja RAG menjadi dua varian utama, yaitu \textit{Vanilla RAG} (atau \textit{Naive RAG}) dan \textit{Advanced RAG}. Karakteristik \textit{Vanilla RAG} menurut \textcite{singhAgenticRetrievalAugmentedGeneration2025} terletak pada kesederhanaan implementasinya yang sesuai untuk penanganan \textit{query} berbasis fakta dengan konteks yang tidak terlalu kompleks. Sebaliknya, \textit{Advanced RAG} mengintegrasikan teknik \textit{retrieval} yang lebih canggih untuk memenuhi kebutuhan aplikasi yang menuntut presisi tinggi serta pemahaman nuansa yang mendalam. Meskipun demikian, keterbatasan alur kerja statis pada kedua pendekatan tradisional tersebut mendorong munculnya paradigma baru, yakni \textit{Agentic RAG}. \textcite{singhAgenticRetrievalAugmentedGeneration2025} menjelaskan bahwa pendekatan ini mengintegrasikan agen otonom yang mampu melakukan pengambilan keputusan dinamis dan penalaran iteratif untuk mengatasi kendala pada paradigma sebelumnya. 

\subsection{\textit{Vanilla RAG}}

Sistem \textit{Vanilla RAG} terdiri dari tiga modul utama, yakni pipa ingesti (\textit{ingestion pipeline}), pipa pengambilan (\textit{retrieval pipeline}), dan pipa generasi (\textit{generation pipeline}). Modul ingesti berfungsi mengisi basis data vektor secara \textit{batch} atau \textit{streaming}, sementara modul pengambilan melakukan \textit{query} untuk mendapatkan entri relevan terhadap masukan pengguna. Modul generasi kemudian menggunakan data tersebut untuk memperkaya \textit{prompt} bagi LLM guna menghasilkan jawaban. Gambaran mengenai kerangka kerja \textit{Vanilla RAG} ini ditampilkan pada Gambar \ref{fig:naive_rag}.

\begin{figure}[H]
\centering
\includegraphics[width=0.7\textwidth]{bab2/naiverag.png}
\caption[Kerangka kerja \textit{Vanilla RAG}]{Kerangka kerja \textit{Vanilla RAG} yang terdiri dari pipa ingesti, pengambilan, dan generasi \parencite{iusztinLLMEngineersHandbook2024}}
\label{fig:naive_rag}
\end{figure}

\textcite{iusztinLLMEngineersHandbook2024} menjelaskan bahwa pipa ingesti bertugas mengekstraksi dokumen mentah dari berbagai sumber data, seperti \textit{data warehouse}, \textit{data lake}, atau halaman web. Data tersebut kemudian dibersihkan, dipecah menjadi bagian-bagian kecil (\textit{chunking}), dan diubah menjadi \textit{embedding} sebelum dimuat ke dalam basis data vektor. Tahapan pemrosesan data agar dapat ditemukan kembali dengan cepat ini oleh \textcite{huyenAIEngineeringBuilding2025} didefinisikan sebagai proses pengindeksan (\textit{indexing}). Lebih lanjut, \textcite{auffarthGenerativeAILangChain2025} menyatakan bahwa lapisan penyimpanan (\textit{storage layer}) yang menampung informasi eksternal setelah melalui proses pengindeksan ini dapat disebut sebagai basis pengetahuan (\textit{knowledge base}). Dalam konteks ini, \textcite{huyenAIEngineeringBuilding2025} juga mencatat bahwa istilah "dokumen" dan "\textit{chunk}" sering digunakan secara bergantian, karena secara teknis potongan dokumen juga merupakan sebuah dokumen tersendiri.

Pada tahap selanjutnya, komponen pengambilan memproses masukan pengguna, baik berupa teks, citra, maupun audio, dengan mengubahnya menjadi \textit{embedding} dan melakukan \textit{query} ke basis data vektor untuk mencari vektor yang serupa \parencite{iusztinLLMEngineersHandbook2024}. Fungsi utama dari langkah ini adalah memproyeksikan masukan pengguna ke dalam ruang vektor yang sama dengan indeks dokumen guna menemukan entri paling mirip (\textit{top-K}) sebagai konteks tambahan. Konteks yang diperoleh tersebut kemudian digabungkan dengan \textit{query} pengguna dalam sebuah \textit{template prompt} untuk diteruskan ke LLM pada tahap generasi. \textit{Prompt} akhir yang terbentuk merupakan hasil konstruksi dari \textit{template} sistem yang telah diisi dengan \textit{query} dan konteks relevan, di mana rekayasa \textit{prompt} (\textit{prompt engineering}) umumnya diterapkan pada level \textit{template} ini.

\subsection{\textit{Advanced RAG}}

\textcite{iusztinLLMEngineersHandbook2024} menyatakan bahwa kerangka kerja \textit{Vanilla RAG} masih memiliki keterbatasan dalam menangani aspek fundamental seperti relevansi dokumen, redundansi informasi, serta latensi sistem. Oleh sebab itu, kerangka kerja \textit{Advanced RAG} dikembangkan untuk mengoptimalkan proses tersebut melalui tiga tahapan utama, yaitu pra-pengambilan (\textit{pre-retrieval}), pengambilan (\textit{retrieval}), dan pasca-pengambilan (\textit{post-retrieval}). Tahap pra-pengambilan berfokus pada strukturisasi dan pemrosesan data guna memaksimalkan efisiensi pengindeksan serta \textit{query}. Tahap pengambilan bertujuan meningkatkan kualitas model \textit{embedding} dan strategi penyaringan metadata untuk memperbaiki hasil pencarian vektor. Sementara itu, tahap pasca-pengambilan menargetkan eliminasi gangguan (\textit{noise}) dari dokumen yang diambil serta kompresi \textit{prompt} sebelum diteruskan ke LLM. Sistem ini juga mensyaratkan adanya modul evaluasi yang kuat untuk mengukur kualitas data yang diambil dan jawaban yang dihasilkan secara kuantitatif. Gambaran mengenai kerangka kerja \textit{Advanced RAG} diperlihatkan pada Gambar \ref{fig:advanced_rag}.

\begin{figure}[H]
\centering
\includegraphics[width=\textwidth]{bab2/advancedrag.png}
\caption[Kerangka kerja \textit{Advanced RAG}]{Kerangka kerja \textit{Advanced RAG} yang terdiri dari tahapan pra-pengambilan, pengambilan, dan pasca-pengambilan \parencite{iusztinLLMEngineersHandbook2024}}
\label{fig:advanced_rag}
\end{figure}

Tahap pra-pengambilan dijalankan melalui dua pendekatan strategis, yakni pengindeksan data dan optimasi \textit{query} \parencite{iusztinLLMEngineersHandbook2024}. Pengindeksan data umumnya diimplementasikan pada modul pembersihan atau pemecahan data (\textit{chunking}) untuk memperbaiki struktur informasi, misalnya dengan menyematkan metadata berupa tanggal, URL, atau penanda bab. Metadata tersebut berperan krusial dalam memfasilitasi penyaringan data yang efisien saat proses pencarian berlangsung. Di sisi lain, optimasi \textit{query} berfokus pada penyempurnaan masukan pengguna agar lebih efektif saat diproses oleh sistem, salah satunya melalui teknik \textit{self-query}. Teknik ini memetakan \textit{query} tidak terstruktur menjadi bentuk terstruktur dengan bantuan LLM untuk mengidentifikasi entitas dan hubungan kunci sebagai parameter filter. Penggunaan parameter tersebut memungkinkan ruang pencarian vektor direduksi secara signifikan untuk meningkatkan relevansi hasil. \citeauthor{iusztinLLMEngineersHandbook2024} menekankan bahwa persiapan data dan \textit{query} yang matang ini merupakan fondasi bagi efisiensi tahapan selanjutnya.

Pada tahap pengambilan, optimasi dapat dicapai dengan meningkatkan kualitas model \textit{embedding} atau memanfaatkan fitur pencarian pada basis data \parencite{iusztinLLMEngineersHandbook2024}. Peningkatan model \textit{embedding} dilakukan untuk menangkap kemiripan semantik antara \textit{query} dan dokumen dengan lebih akurat pada saat inferensi. Pemanfaatan fitur basis data sering kali menjadi solusi yang lebih praktis, terutama melalui penerapan metode pencarian hibrida (\textit{hybrid search}). Metode ini menggabungkan pencarian berbasis kata kunci yang presisi dalam mencocokkan istilah spesifik dengan pencarian vektor yang unggul dalam menemukan kemiripan semantik. Selain itu, teknik pencarian vektor terfilter (\textit{filtered vector search}) juga dapat diterapkan dengan memanfaatkan indeks metadata untuk membatasi ruang pencarian sebelum atau sesudah pencarian vektor. Sinergi antara pencocokan kata kunci dan pemahaman semantik ini dinilai sangat efektif untuk tugas yang menuntut akurasi pengambilan informasi yang tinggi.

Tahap pasca-pengambilan difokuskan pada pemrosesan ulang data yang telah diambil untuk memastikan kinerja LLM tidak terganggu oleh konteks yang terlalu besar atau informasi yang tidak relevan \parencite{iusztinLLMEngineersHandbook2024}. Salah satu metode populer dalam tahap ini adalah \textit{reranking}, yang mengurutkan ulang dokumen menggunakan model \textit{machine learning} untuk memastikan hanya informasi yang paling relevan yang digunakan sebagai konteks akhir. Teknik-teknik optimasi ini merupakan contoh aspek yang dapat ditingkatkan dalam alur kerja RAG, namun penerapannya sangat bergantung pada jenis data yang dikelola. Teknik yang dirancang untuk teks mungkin tidak sepenuhnya dapat diterapkan pada data \textit{multimodal} seperti citra atau video yang memiliki karakteristik berbeda.

\subsection{\textit{Information Retrieval}}
Di balik layar, RAG menggunakan teknik \textit{Information Retrieval} (IR) untuk mengambil dokumen yang relevan guna menjawab pertanyaan \parencite{huyenAIEngineeringBuilding2025}. IR merupakan bidang yang mencakup pengambilan segala jenis media berdasarkan kebutuhan informasi pengguna \parencite{jm3}. Konsep ini merupakan gagasan yang telah berumur satu abad dan menjadi tulang punggung bagi mesin pencari, sistem rekomendasi, serta analitik \textit{log} \parencite{huyenAIEngineeringBuilding2025}. Banyak algoritma \textit{retrieval} yang dikembangkan untuk sistem pencarian tradisional tersebut juga dapat diterapkan dalam RAG. Sejarahnya dapat ditarik kembali hingga tahun 1920-an, saat Emanuel Goldberg mendeskripsikan "mesin statistik" dalam patennya untuk mencari dokumen yang tersimpan pada film \parencite{al-aminHistoryGenerativeArtificial2024}. 

Teknik IR yang umumnya digunakan dalam RAG adalah \textit{ad hoc retrieval}, di mana pengguna mengajukan \textit{query} kepada sistem pengambilan informasi dan sistem tersebut mengembalikan sekumpulan dokumen yang diurutkan berdasarkan relevansinya \parencite{jm3}. Dalam konteks ini, \textit{query} merepresentasikan kebutuhan informasi pengguna, dokumen merujuk pada satuan teks yang diindeks, dan koleksi merupakan kumpulan dokumen yang digunakan untuk memenuhi permintaan tersebut. Arsitektur umum sistem \textit{ad hoc retrieval} diilustrasikan pada Gambar \ref{fig:adhoc_ir}. Peringkat dokumen ditentukan dengan menghitung skor relevansi setiap dokumen kandidat terhadap \textit{query} yang diberikan, yang mengekspresikan seberapa relevan dokumen tersebut dalam memenuhi kebutuhan informasi pengguna. Secara umum, sistem IR ini terbagi menjadi dua kelas berdasarkan jenis vektor yang digunakan untuk merepresentasikan \textit{query} dan dokumen, yaitu \textit{sparse vector} dan \textit{dense vector}.

\begin{figure}[H]
\centering
\includegraphics[width=0.9\textwidth]{bab2/adhoc-ir.png}
\caption[Arsitektur sistem \textit{ad hoc information retrieval}]{Arsitektur sistem \textit{ad hoc information retrieval} \parencite{jm3}}
\label{fig:adhoc_ir}
\end{figure}

\textit{Sparse retrieval}, yang juga disebut \textit{term-based retrieval} \parencite{huyenAIEngineeringBuilding2025} atau \textit{keyword-based retrieval} \parencite{auffarthGenerativeAILangChain2025}, bekerja pada level leksikal dengan mencocokkan istilah dalam \textit{query} secara langsung terhadap dokumen. Setiap kata dalam kosakata direpresentasikan sebagai \textit{one-hot vector} dengan ukuran sebesar panjang kosakata, sehingga representasi dokumen dan \textit{query} menjadi sangat jarang. Pendekatan ini mampu mencocokkan istilah-istilah spesifik secara tepat, namun cenderung gagal menangkap kesamaan semantik antara kata-kata yang bersinonim atau memiliki makna serupa.

\textit{Dense retrieval} di sisi lain mencocokkan dokumen berdasarkan kesamaan semantik antara \textit{query} dan dokumen, bukan sekadar tumpang tindih kata kunci \parencite{huyenAIEngineeringBuilding2025}. Pendekatan ini dikenal sebagai \textit{embedding-based retrieval}, karena representasi vektor padat dihasilkan oleh model \textit{embedding} yang menangkap makna semantik dari teks. Walaupun pembagian antara \textit{sparse} dan \textit{dense} tampak jelas, terdapat teknik yang berada di antara keduanya. Contohnya, \textcite{formalSPLADESparseLexical2021} memperkenalkan SPLADE (\textit{Sparse Lexical and Expansion}), algoritma pengambilan informasi yang menggunakan \textit{sparse embedding}. SPLADE memanfaatkan representasi dari model BERT dengan regularisasi untuk mendorong sebagian besar nilai \textit{embedding} menuju nol sehingga operasi menjadi lebih efisien, dan karena karakteristik \textit{sparse}-nya tersebut, SPLADE umumnya dikelompokkan bersama algoritma \textit{sparse retrieval}.

\subsubsection{\textit{Sparse Retrieval}}
\textit{Sparse retrieval} yang disebut juga \textit{lexical retrieval} merupakan pendekatan pencarian dokumen yang paling mendasar dengan memanfaatkan kata kunci dari \textit{query} \parencite{huyenAIEngineeringBuilding2025}. Metode ini bekerja dengan asumsi bahwa semakin sering suatu istilah muncul dalam sebuah dokumen, semakin relevan dokumen tersebut terhadap \textit{query} yang diberikan. Namun, ketergantungan semata pada frekuensi kemunculan kata (\textit{Term Frequency} atau TF) memiliki kelemahan signifikan, di mana kata-kata umum yang tidak informatif seperti kata penghubung dapat mendominasi skor relevansi. Untuk mengatasi hal ini, diperkenalkan konsep \textit{Inverse Document Frequency} (IDF) yang memberikan bobot lebih rendah pada kata yang muncul di banyak dokumen, sehingga sistem dapat memprioritaskan istilah yang lebih spesifik dan bermakna.

Algoritma TF-IDF menggabungkan kedua metrik tersebut untuk menghitung skor relevansi dokumen secara lebih akurat \parencite{huyenAIEngineeringBuilding2025}. Dalam mekanisme ini, nilai penting suatu kata berbanding terbalik dengan jumlah dokumen yang memuatnya. Secara matematis, skor total dokumen dihitung dengan menjumlahkan hasil perkalian antara IDF dan TF dari setiap kata yang terdapat dalam \textit{query}. Implementasi teknik ini umumnya memanfaatkan struktur data \textit{inverted index}, sebuah "kamus" yang memetakan istilah ke daftar dokumen yang memuatnya. Struktur ini memungkinkan pencarian teks yang sangat efisien sekaligus menyimpan metadata statistik yang diperlukan untuk perhitungan skor.

Salah satu varian TF-IDF yang paling banyak diadopsi dalam industri adalah Okapi BM25 yang pertama kali diperkenalkan oleh \textcite{harman1995trec3}. Algoritma ini menyempurnakan pendekatan TF-IDF standar dengan menambahkan normalisasi terhadap panjang dokumen dan saturasi frekuensi istilah \parencite{huyenAIEngineeringBuilding2025}. Normalisasi ini penting karena dokumen yang lebih panjang secara alami cenderung memuat lebih banyak kata kunci tanpa tentu lebih relevan. Meskipun metode \textit{retrieval} modern berbasis vektor semakin populer, BM25 dan variannya tetap menjadi standar industri yang kuat dan sering digunakan sebagai \textit{baseline} yang tangguh untuk membandingkan kinerja model pencarian yang lebih canggih.

\subsubsection{\textit{Dense Retrieval}}
Berbeda dengan \textit{sparse retrieval} yang bekerja pada level leksikal, \textit{dense retrieval} mencocokkan dokumen berdasarkan kesamaan semantik antara \textit{query} dan dokumen \parencite{huyenAIEngineeringBuilding2025}. Pendekatan ini, yang juga dikenal sebagai \textit{semantic retrieval} atau \textit{embedding-based retrieval}, menggunakan model \textit{embedding} untuk mengubah teks menjadi vektor padat yang merepresentasikan maknanya. Ketika suatu \textit{query} diajukan, \textit{query} tersebut diubah menjadi \textit{embedding} menggunakan model yang sama dengan yang dipakai saat pengindeksan, kemudian sistem mengambil $k$ potongan data yang vektornya paling dekat dengan vektor \textit{query}. Dengan demikian, \textit{dense retrieval} dapat menangani kasus di mana kata kunci dalam \textit{query} tidak muncul secara harfiah dalam dokumen yang relevan.

Pembentukan representasi vektor padat umumnya dilakukan menggunakan arsitektur \textit{Bi-Encoder} \parencite{jm3}. \textcite{reimersSentenceBERTSentenceEmbeddings2019} memperkenalkan Sentence-BERT (SBERT), sebuah modifikasi dari model BERT yang menggunakan jaringan siam (\textit{siamese network}) untuk menghasilkan \textit{embedding} kalimat yang bermakna secara semantik. Dalam arsitektur ini, dua input diproses secara paralel oleh model transformer dengan bobot yang sama (\textit{shared weights}), kemudian representasi keduanya dibandingkan menggunakan ukuran kesamaan seperti \textit{cosine similarity}, sebagaimana diilustrasikan pada Gambar \ref{fig:sentence_transformer}. Pendekatan ini memungkinkan komputasi kesamaan antar dokumen dalam skala besar secara efisien, berbeda dengan arsitektur \textit{Cross-Encoder} yang harus memproses pasangan teks secara bersamaan dan tidak praktis untuk pengindeksan skala besar.

\begin{figure}[H]
\centering
\includegraphics[width=0.6\textwidth]{bab2/sentencetransformer.png}
\caption[Arsitektur Sentence-BERT]{Arsitektur Sentence-BERT untuk menghasilkan \textit{embedding} kalimat \parencite{reimersSentenceBERTSentenceEmbeddings2019}}
\label{fig:sentence_transformer}
\end{figure}

Komponen krusial lain dalam \textit{dense retrieval} adalah basis data vektor (\textit{vector database}), yaitu basis data khusus yang dirancang untuk menyimpan dan mencari vektor \parencite{huyenAIEngineeringBuilding2025}. Pencarian vektor umumnya dirumuskan sebagai masalah \textit{nearest-neighbor search}: diberikan sebuah vektor \textit{query}, temukan $k$ vektor dalam basis data yang paling mirip berdasarkan metrik seperti \textit{cosine similarity}. Pendekatan naif $k$-NN yang menghitung kemiripan terhadap seluruh vektor dalam basis data memberikan hasil yang presisi, namun tidak efisien secara komputasi untuk dataset besar. Karena itu, pencarian vektor pada skala besar umumnya dilakukan menggunakan algoritma \textit{approximate nearest neighbor} yang menukar sedikit presisi dengan peningkatan kecepatan yang signifikan. Salah satu implementasi basis data vektor yang mendukung pencarian \textit{approximate nearest neighbor} adalah Qdrant \parencite{qdrantQdrant2025}.

Secara umum, basis data vektor mengorganisasi vektor ke dalam struktur seperti \textit{bucket}, pohon, atau graf untuk meningkatkan efisiensi pencarian \parencite{huyenAIEngineeringBuilding2025}. Salah satu algoritma \textit{approximate nearest neighbor} yang banyak digunakan adalah \textit{Hierarchical Navigable Small World} (HNSW), yang diperkenalkan oleh \textcite{malkovEfficientRobustApproximate2020}. HNSW membangun graf berlapis di mana setiap simpul merepresentasikan sebuah vektor, dan tepi-tepinya menghubungkan vektor-vektor yang berdekatan, sehingga pencarian \textit{nearest-neighbor} dilakukan dengan menelusuri tepi-tepi graf tersebut.

Walaupun basis data vektor muncul sebagai kategori tersendiri seiring berkembangnya RAG, pada dasarnya setiap basis data yang mampu menyimpan vektor dapat berfungsi sebagai basis data vektor \parencite{huyenAIEngineeringBuilding2025}. Saat ini, banyak basis data tradisional yang telah memperluas kapabilitasnya untuk mendukung penyimpanan dan pencarian vektor sebagai fitur tambahan. Fenomena ini menyebabkan batas klasifikasi antara basis data konvensional dan basis data vektor murni menjadi semakin kabur.

\subsection{\textit{Hybrid Search}}
\textit{Hybrid search} merupakan pendekatan yang menggabungkan beberapa metode \textit{retrieval} untuk memanfaatkan kelebihan masing-masing algoritma \parencite{huyenAIEngineeringBuilding2025}. Penggabungan ini dapat dilakukan secara berurutan maupun secara paralel sebagai \textit{ensemble}. Pada pendekatan berurutan, \textit{retriever} yang lebih ringan secara komputasi, seperti BM25, terlebih dahulu mengambil kandidat dokumen, kemudian hasil tersebut disaring lebih lanjut melalui tahap \textit{reranking} menggunakan pencarian vektor yang lebih presisi. Pada pendekatan paralel, beberapa \textit{retriever} dijalankan bersamaan, dan hasil peringkat dari masing-masing \textit{retriever} digabungkan menjadi peringkat akhir. Algoritma yang umum digunakan untuk menggabungkan peringkat dari beberapa \textit{retriever} adalah \textit{Reciprocal Rank Fusion} (RRF) dari \textcite{cormackReciprocalRankFusion2009}, di mana setiap dokumen mendapat skor berdasarkan posisi peringkatnya pada tiap \textit{retriever}, dan skor-skor tersebut dijumlahkan untuk menghasilkan peringkat akhir.

\subsection{\textit{Reranking}}
Peringkat dokumen awal yang dihasilkan oleh \textit{retriever} dapat disempurnakan lebih lanjut melalui proses \textit{reranking} \parencite{huyenAIEngineeringBuilding2025}. \textit{Reranking} bertujuan menghasilkan urutan dokumen yang lebih presisi sekaligus memungkinkan pengurangan jumlah dokumen yang diteruskan ke model, baik untuk menyesuaikan batas konteks model maupun mengurangi jumlah token masukan.Pada sistem \textit{dense retrieval}, \textit{retriever} umumnya menggunakan arsitektur \textit{bi-encoder} \parencite{jm3}. Pada arsitektur ini, dokumen di-\textit{encode} terlebih dahulu dan disimpan sebagai vektor, sehingga ketika \textit{query} masuk, hanya \textit{query} yang perlu di-\textit{encode} dan relevansinya dihitung menggunakan \textit{dot product} terhadap vektor dokumen yang telah tersimpan. Pendekatan ini efisien secara komputasi, namun kurang presisi karena tidak dapat menangkap sepenuhnya interaksi makna antar token antara \textit{query} dan dokumen akibat keduanya diproses secara independen.

Alternatif dari \textit{bi-encoder} adalah \textit{cross-encoder}, di mana \textit{query} dan dokumen dimasukkan sekaligus ke dalam \textit{Transformer} untuk menghasilkan skor relevansi \parencite{reimersSentenceBERTSentenceEmbeddings2019}. Karena kedua teks diproses bersamaan, \textit{cross-encoder} dapat memodelkan hubungan antar token secara langsung sehingga presisinya lebih tinggi. Akan tetapi, \textit{cross-encoder} tidak menghasilkan \textit{embedding} tersendiri dan setiap pasangan \textit{query}-dokumen harus diproses bersama, sehingga tidak praktis untuk digunakan langsung pada korpus besar. Perbandingan arsitektur keduanya diilustrasikan pada Gambar \ref{fig:biencoder_vs_crossencoder}.

\begin{figure}[H]
\centering
\includegraphics[width=0.7\textwidth]{bab2/biencoder_vs_crossencoder.png}
\caption[Perbandingan arsitektur \textit{bi-encoder} dan \textit{cross-encoder}]{Perbandingan arsitektur \textit{bi-encoder} dan \textit{cross-encoder} \parencite{reimersSentenceBERTSentenceEmbeddings2019}}
\label{fig:biencoder_vs_crossencoder}
\end{figure}

Dari karakteristik masing-masing arsitektur tersebut, pola umum dalam sistem \textit{reranking} adalah menggunakan \textit{bi-encoder} yang murah, namun kurang presisi untuk mengambil banyak kandidat dokumen pada tahap pertama guna memastikan \textit{recall}, kemudian mengurutkan ulang kandidat tersebut menggunakan \textit{cross-encoder} yang lebih mahal, namun lebih presisi untuk menghasilkan peringkat akhir \parencite{huyenAIEngineeringBuilding2025}. Sebagai alternatif \textit{cross-encoder}, ColBERT \parencite{khattabColBERTEfficientEffective2020} menawarkan pendekatan \textit{late interaction} yang berupaya menyeimbangkan biaya komputasi dan presisi. ColBERT mengenkode \textit{query} dan dokumen secara terpisah hingga representasi per token, menghasilkan sekumpulan vektor untuk setiap token yang dikenal sebagai \textit{multivectors}. Relevansi kemudian dihitung menggunakan operator \textit{MaxSim}, yaitu setiap token \textit{query} dicocokkan dengan token dokumen yang paling mirip, kemudian seluruh skor tersebut dijumlahkan. Karena \textit{multivectors} dokumen dapat disimpan sebelumnya, ColBERT dapat lebih efisien dari \textit{cross-encoder}. Cara kerja ColBERT diilustrasikan pada Gambar \ref{fig:colbert}.

\begin{figure}[H]
\centering
\includegraphics[width=0.4\textwidth]{bab2/colbert.png}
\caption[Ilustrasi mekanisme \textit{late interaction} pada ColBERT]{Ilustrasi mekanisme \textit{late interaction} pada ColBERT menggunakan operator \textit{MaxSim} \parencite{khattabColBERTEfficientEffective2020}}
\label{fig:colbert}
\end{figure}

\textcite{huyenAIEngineeringBuilding2025} juga menekankan bahwa \textit{reranking} dalam konteks RAG sedikit berbeda dari \textit{reranking} pada sistem pencarian konvensional. Pada sistem pencarian, posisi peringkat sangat menentukan karena pengguna cenderung memperhatikan hasil teratas. Dalam RAG, urutan dokumen dalam konteks tetap berpengaruh karena model cenderung lebih baik memproses informasi yang terletak di awal dan akhir konteks, sehingga penempatan dokumen yang paling relevan pada posisi tersebut dapat meningkatkan kualitas respons yang dihasilkan.

\subsection{Qdrant}
\textcite{qdrantQdrant2025} adalah \textit{library} (pustaka) \textit{open-source} (sumber terbuka) yang menyediakan mesin pencari vektor (\textit{vector search engine}) dan basis data vektor yang dirancang untuk menyimpan, mengelola, dan mencari vektor \textit{embedding} dalam skala besar dengan efisien. Qdrant menyediakan mekanisme pencarian kesamaan (\textit{similarity search}) melalui berbagai teknik IR dan algoritma pengambilan yang ada di industri.

Entitas utama yang menjadi unit penyimpanan dasar dalam \textcite{qdrantQdrant2025} dikenal sebagai \textit{Point}. Setiap \textit{Point} merepresentasikan catatan data tunggal yang terdiri dari tiga komponen utama, yaitu pengenal unik (ID), representasi vektor, dan \textit{payload} opsional. Struktur dari sebuah \textit{Point} dalam Qdrant diilustrasikan pada Gambar \ref{fig:qdrant_points}.

\begin{figure}[H]
\centering
\includegraphics[width=0.8\textwidth]{bab2/qdrantpoints.png}
\caption[Struktur Data \textit{Point} pada Qdrant]{Struktur Data \textit{Point} pada Qdrant yang Terdiri dari ID, Vektor, dan \textit{Payload} \parencite{qdrantQdrant2025}}
\label{fig:qdrant_points}
\end{figure}

Komponen ID pada struktur data \textcite{qdrantQdrant2025} berfungsi sebagai identitas unik untuk setiap \textit{Point}, yang dapat berupa bilangan bulat 64-bit (\textit{unsigned integer}) atau UUID. Jika pengguna tidak menyediakan ID saat pembuatan \textit{Point}, klien Qdrant secara otomatis akan menghasilkan UUID acak. Komponen kedua adalah vektor, yang merupakan representasi numerik dari data. Qdrant memiliki fleksibilitas tinggi dalam penyimpanan vektor, dengan dukungan terhadap berbagai jenis vektor seperti vektor padat (\textit{dense vector}), vektor jarang (\textit{sparse vector}), maupun \textit{multivector}. Lebih lanjut, Qdrant memungkinkan penyimpanan beberapa jenis vektor sekaligus dalam satu \textit{Point} melalui fitur \textit{named vectors}. Fitur ini berguna dalam skenario yang kompleks, seperti pada model \textit{multimodal} di mana satu entitas data mungkin memiliki representasi vektor untuk teks dan citra secara bersamaan.

Komponen terakhir pada struktur data \textcite{qdrantQdrant2025} adalah \textit{payload}, yang merupakan metadata tambahan dalam format JSON yang terkait dengan vektor tersebut. \textit{Payload} memungkinkan penyimpanan informasi non-vektor seperti atribut deskriptif, kategori, atau label waktu. Fungsi utama dari \textit{payload} adalah untuk memfasilitasi mekanisme penyaringan (\textit{filtering}) selama proses pencarian. Dengan fitur ini, Qdrant dapat melakukan pencarian hibrida yang menggabungkan pencarian semantik berbasis vektor dengan penyaringan berbasis aturan pada metadata, sehingga hasil yang diperoleh menjadi lebih spesifik dan relevan sesuai dengan kriteria yang diinginkan.

\subsection{Evaluasi Sistem \textit{Information Retrieval}}

Kualitas sebuah \textit{retriever} dapat dievaluasi berdasarkan kualitas data yang berhasil diambilnya \parencite{huyenAIEngineeringBuilding2025}. Dua metrik yang paling umum digunakan dalam evaluasi sistem IR adalah \textit{precision} dan \textit{recall} \parencite{jm3}. \textit{Precision} mengukur proporsi dokumen yang dikembalikan sistem yang benar-benar relevan, sedangkan \textit{recall} mengukur proporsi dokumen relevan yang berhasil ditemukan dari keseluruhan dokumen relevan yang tersedia. Untuk menghitung kedua metrik ini, dibutuhkan sebuah \textit{evaluation set} yang terdiri dari sekumpulan \textit{query} pengujian beserta dokumen yang telah dianotasi relevansinya, baik oleh manusia maupun \textit{AI judge} \parencite{huyenAIEngineeringBuilding2025}. Nilai \textit{precision} dan \textit{recall} kemudian dihitung berdasarkan hasil pengambilan \textit{retriever} pada \textit{evaluation set} tersebut.

Kedua metrik di atas tidak secara memadai mengukur kinerja sistem yang mengembalikan dokumen dalam urutan peringkat tertentu \parencite{jm3}. Penilaian berbasis \textit{precision} dan \textit{recall} tidak memperhitungkan posisi dokumen relevan dalam daftar hasil, padahal dokumen yang lebih relevan idealnya berada di posisi lebih atas. Untuk kasus seperti ini, metrik seperti NDCG (\textit{normalized discounted cumulative gain}) dan MRR (\textit{Mean Reciprocal Rank}) lebih tepat digunakan karena keduanya memperhitungkan urutan peringkat dalam penilaian kualitas pengambilan \parencite{huyenAIEngineeringBuilding2025}.

Seluruh metrik evaluasi IR tersebut dapat dihitung menggunakan ranx \parencite{bassaniRanxBlazingFastPython2022}, sebuah \textit{library} Python yang menyediakan berbagai metrik evaluasi peringkat secara cepat. ranx mendukung evaluasi \textit{offline} yang memanfaatkan \textit{ground truth} berupa pasangan \textit{query} dan dokumen relevan yang lazim disebut \textit{qrels} dalam IR. \textit{Library} ini mendukung seluruh metrik umum yang telah disebutkan, serta memungkinkan evaluasi beberapa \textit{run} dan metrik sekaligus dengan memparalelkan komputasi, sehingga cocok untuk evaluasi berskala besar.

\section{Sistem Agentik}
Menurut \textcite{anthropicBuildingEffectiveAgents2024}, sistem agentik (\textit{agentic system}) adalah pendekatan untuk memperluas kemampuan LLM dengan memberi akses pada alat, memori, atau sumber informasi eksternal sehingga model dapat bertindak lebih dari sekadar menghasilkan teks. Kemampuan \textit{foundation model} yang terus berkembang telah membuka kemungkinan bagi aplikasi-aplikasi agentik yang sebelumnya sulit diwujudkan \parencite{huyenAIEngineeringBuilding2025}. Paradigma ini umumnya dikenal dengan sebutan \textit{AI agent}, walaupun \textcite{anthropicBuildingEffectiveAgents2024} lebih memilih istilah sistem agentik untuk menekankan sifat sistemiknya. Konsep agen otonom sendiri bukan gagasan baru, di mana \textcite{russell2020artificial} secara eksplisit mendefinisikan kecerdasan buatan sebagai studi dan perancangan agen rasional, menjadikan agen cerdas yang otonom sebagai tujuan akhir dari bidang ini. Walaupun potensinya besar, \textcite{huyenAIEngineeringBuilding2025} mencatat bahwa sistem agentik berbasis AI masih merupakan bidang yang berkembang, dengan belum adanya kerangka teoretis yang mapan untuk mendefinisikan, mengembangkan, dan mengevaluasinya.

Secara struktural, sistem agentik terdiri dari empat komponen utama yang bekerja secara terintegrasi, sebagaimana diilustrasikan pada Gambar \ref{fig:komponen_sistem_agentik} \parencite{singhAgenticRetrievalAugmentedGeneration2025}. LLM dengan peran (\textit{role}) dan tugas (\textit{task}) yang terdefinisi berfungsi sebagai mesin penalaran utama yang menginterpretasikan \textit{query} pengguna dan menghasilkan respons. Memori dibagi menjadi memori jangka pendek (\textit{short-term memory}) untuk melacak konteks percakapan terkini, dan memori jangka panjang (\textit{long-term memory}) untuk menyimpan pengetahuan atau pengalaman yang terakumulasi. Komponen perencanaan (\textit{planning}) memandu proses penalaran sistem melalui mekanisme refleksi dan evaluasi mandiri (\textit{self-critique}) agar tugas-tugas kompleks dapat dipecah secara sistematis. Terakhir, alat-alat (\textit{tools}) seperti \textit{vector search}, \textit{web search}, dan API memperluas kemampuan sistem melampaui pembangkitan teks biasa, memungkinkan akses ke sumber data eksternal dan komputasi khusus.

\begin{figure}[H]
    \centering
    \includegraphics[width=0.85\textwidth]{bab2/arsitektur_sistem_agentik.png}
    \caption[Komponen-komponen sistem agentik]{Komponen-komponen sistem agentik yang terdiri dari LLM, memori, perencanaan, dan alat \parencite{singhAgenticRetrievalAugmentedGeneration2025}}
    \label{fig:komponen_sistem_agentik}
\end{figure}


\textcite{anthropicBuildingEffectiveAgents2024} membuat pembedaan arsitektur yang penting dalam sistem agentik antara \textit{workflow} dan \textit{agent}. \textit{Workflow} adalah sistem di mana LLM dan alat diatur melalui jalur kode yang telah ditentukan sebelumnya, sedangkan \textit{agent} adalah sistem di mana LLM secara dinamis mengarahkan proses dan penggunaan alatnya sendiri untuk menyelesaikan tugas. Pembedaan ini bukan bersifat preskriptif, melainkan merupakan pola-pola umum yang dapat dikombinasikan dan disesuaikan dengan kebutuhan kasus penggunaan yang berbeda \parencite{anthropicBuildingEffectiveAgents2024}.

\textcite{anthropicBuildingEffectiveAgents2024} mengenalkan beberapa pola \textit{workflow} yang umum dalam pengembangan sistem agentik. Pola \textit{prompt chaining} menyusun pemanggilan LLM secara berurutan, di mana \textit{output} satu tahap menjadi \textit{input} tahap berikutnya, seperti yang diterapkan \textcite{trautmannLargeLanguageModel2023} untuk memproses dokumen hukum melalui tahap peringkasan, pencarian semantik, hingga klasifikasi berbasis LLM. Pola \textit{routing} mengklasifikasikan \textit{input} lalu mengarahkannya ke alat atau proses yang sesuai, seperti yang digunakan \textcite{xuMultimodalAgriculturalAgent2025} dengan model klasifikasi sebagai \textit{router} untuk menentukan alat yang relevan berdasarkan gambar dan pertanyaan yang diterima. Pola \textit{parallelization} memungkinkan beberapa LLM memproses \textit{subtask} secara bersamaan (\textit{sectioning}) atau menjalankan tugas yang sama berulang kali untuk keragaman hasil (\textit{voting}), sebagaimana diimplementasikan \textcite{xiaParallelismMeetsAdaptiveness2025} yang mengombinasikannya dengan adaptasi multi-agen.

Dua pola yang lebih kompleks adalah \textit{orchestrator-workers} dan \textit{evaluator-optimizer}. Pola \textit{orchestrator-workers} menggunakan satu LLM sebagai orkestrator yang secara dinamis mendelegasikan \textit{subtask} ke beberapa LLM pekerja, sementara pola \textit{evaluator-optimizer} menambahkan evaluasi iteratif di mana satu LLM memberikan umpan balik atas \textit{output} LLM lain dalam sebuah siklus penyempurnaan. Kedua pola ini terlihat pada kerangka multi-agen yang dikembangkan \textcite{keDynamicOrchestrationMultiAgent2025}, di mana orkestrator dinamis mengelola peran \textit{Retriever}, \textit{Reflector}, dan dua \textit{Answerer} yang beroperasi secara paralel, sementara komponen \textit{Improver} secara iteratif menyempurnakan jawaban akhir dengan menyelaraskan informasi dari banyak gambar.

Berbeda dengan \textit{workflow}, arsitektur \textit{agent} menempatkan LLM sebagai pengendali yang secara dinamis menentukan langkah-langkah dan alat yang digunakan selama eksekusi tugas \parencite{anthropicBuildingEffectiveAgents2024}. Agen dimulai dengan instruksi atau diskusi awal bersama pengguna, lalu beroperasi secara mandiri sambil terus mendapatkan umpan balik dari lingkungan di setiap langkah, seperti hasil pemanggilan alat atau keluaran eksekusi kode. Mekanisme ini memungkinkan agen menilai kemajuannya sendiri, memutuskan kapan harus berhenti, atau kembali meminta masukan manusia. Pada dasarnya, sebuah \textit{agent} tidak lebih dari LLM yang menggunakan alat secara berulang berdasarkan umpan balik lingkungan, sehingga desain alat dan dokumentasinya menjadi faktor penentu kualitas sistem \parencite{anthropicBuildingEffectiveAgents2024}. LLM yang berperan sebagai pengendali dalam arsitektur ini oleh \textcite{minaeeLargeLanguageModels2025} disebut \textit{LLM agent} atau oleh \textcite{huyenAIEngineeringBuilding2025} disebut sebagai \textit{foundation model agent}, karena mencerminkan bahwa kemampuan model fondasi seperti pemahaman instruksi kompleks, penalaran, dan penggunaan alat merupakan prasyarat utamanya.

Kemampuan penggunaan alat pada model fondasi mulai dikembangkan secara sistematis melalui Toolformer oleh \textcite{schickToolformerLanguageModels2023}, yang pertama kali menunjukkan bahwa LLM dapat dilatih untuk menggunakan alat eksternal secara mandiri, suatu kemampuan yang kini ummum menjadi fitur pada model fondasi seperti pada keluarga model Gemini. Untuk mengembangkan \textit{LLM agent} yang efektif, model perlu diaugmentasi dengan teknik-teknik tertentu, salah satunya adalah teknik ReAct dari \textcite{yaoReActSynergizingReasoning2023}. Teknik ini kemudian banyak dijadikan fondasi dalam pengembangan sistem yang melibatkan agen dan telah banyak diadopsi oleh berbagai proyek, seperti LangChain, AutoGPT, dan DSPy \parencite{yuSurveyAgentWorkflow2025}.

\subsection{Ekosistem LangChain}
Ekosistem LangChain adalah kumpulan \textit{framework} (kerangka kerja), \textit{tools} (alat), dan layanan yang dikembangkan oleh \textcite{langchainaiLangChain2025} untuk memfasilitasi pembangunan, pengujian, dan penerapan (\textit{deployment}) aplikasi berbasis LLM, termasuk sistem agentik \parencite{narayanMasteringLangChainComprehensive2025}. Pada awalnya, LangChain membangun aplikasi LLM di atas konsep \textit{chain}, yakni rangkaian yang menghubungkan LLM, instruksi \textit{prompt}, dan \textit{output parser} opsional menjadi satu alur pemrosesan \parencite{oshinLearningLangchainBuilding2025}. Untuk mendukung pola ini, LangChain menyediakan komponen-komponen utama, di antaranya \textit{prompts and templates} untuk mendefinisikan interaksi yang terstruktur dengan model, \textit{chains} sebagai mekanisme penghubung antarkompoben dalam suatu urutan pemrosesan, \textit{memory} untuk mempertahankan konteks lintas percakapan, \textit{agents} sebagai pengambil keputusan otonom yang dapat menggunakan berbagai alat, serta \textit{document loaders} dan sistem \textit{retrieval} untuk mengintegrasikan sumber pengetahuan eksternal ke dalam aplikasi \parencite{narayanMasteringLangChainComprehensive2025}.

Sebelum rilis LangChain 1.0, seluruh komponen ini menerapkan antarmuka yang seragam melalui metode \textit{invoke}, \textit{stream}, dan \textit{batch}, yang dapat dikomposisikan secara deklaratif menggunakan \textit{LangChain Expression Language} (LCEL) \parencite{oshinLearningLangchainBuilding2025}. LCEL mengompilasi komposisi tersebut menjadi rencana eksekusi yang dioptimalkan dengan paralelisasi otomatis, \textit{streaming}, \textit{tracing}, dan dukungan operasi asinkron. Ekosistem LangChain pada masa itu juga diperluas dengan tiga \textit{library} pendamping, yaitu LangGraph, LangServe, dan LangSmith \parencite{narayanMasteringLangChainComprehensive2025}. LangGraph \parencite{langchainaiLangGraph2025} merupakan \textit{framework} tingkat rendah yang memungkinkan pembangunan \textit{agent} berbasis \textit{stateful graph}, di mana alur eksekusi dimodelkan sebagai simpul (\textit{node}) dan sisi (\textit{edge}) dalam graf terarah, dengan dukungan \textit{persistence}, \textit{memory}, \textit{human-in-the-loop}, dan \textit{streaming} secara bawaan. LangServe \parencite{langserve} memungkinkan pengembang memaparkan \textit{chain} atau \textit{agent} sebagai API berbasis FastAPI \parencite{fastapi}, sedangkan LangSmith \parencite{langchainaiLangSmithSDK2025} menyediakan \textit{Software Development Kit} (SDK) dan \textit{platform} observabilitas untuk memantau serta mengevaluasi perilaku aplikasi LLM.

Ekosistem ini kemudian mengalami perubahan signifikan ketika LangChain mengumumkan rilis LangChain dan LangGraph 1.0, yang menandai pergeseran arsitektur dari pendekatan \textit{chain} berbasis LCEL menuju LangGraph sebagai \textit{runtime} (lingkungan eksekusi) utama \parencite{langchain2024alpha}. Bersamaan dengan rilis tersebut, LangGraph Platform diumumkan sebagai pengganti resmi LangServe untuk proyek baru, sementara LangServe sendiri beralih ke status \textit{maintenance-only} (hanya pemeliharaan) \parencite{langserve}. LangGraph Platform kemudian diintegrasikan ke dalam LangSmith dan di-\textit{rebrand} (diganti nama) menjadi LangSmith Deployment, memperluas peran LangSmith dari sekadar \textit{platform} observabilitas menjadi pusat pengelolaan siklus hidup aplikasi LLM secara menyeluruh \parencite{langchain2024langgraph}. LangChain 1.0, yang dibangun di atas LangGraph, menawarkan antarmuka tingkat tinggi dengan berbagai komponen \textit{prebuilt} (siap pakai) untuk mempercepat pengembangan aplikasi. Perbandingan keempat komponen utama ekosistem LangChain dirangkum pada Tabel \ref{tab:langchain_ecosystem}.

{
\setstretch{1.0}
\small
\begin{longtblr}[
    caption = {Perbandingan komponen ekosistem LangChain},
    entry = {Perbandingan komponen ekosistem LangChain},
    label = {tab:langchain_ecosystem},
]{
    width = \linewidth,
    colspec = {X[1.5,l] X[2,l] X[2,l] X[2,l] X[2,l]},
    hlines, vlines,
    rows = {m},
    rowsep = 3pt,
    rowhead = 1,
    row{1} = {font=\bfseries, halign=c, valign=m},
    column{1} = {font=\bfseries},
}
    Aspek & LangChain & LangGraph & LangServe & LangSmith \\
    Tujuan & \textit{Framework} tingkat tinggi untuk membangun aplikasi LLM dengan komponen \textit{prebuilt} & \textit{Framework} tingkat rendah untuk membangun \textit{agent} berbasis \textit{stateful graph} & Memaparkan aplikasi LangChain sebagai REST API & \textit{Platform} komersial untuk observabilitas, evaluasi \textit{prompt}, dan \textit{deployment} aplikasi LLM \\
    Fungsi utama & \textit{Chain} siap pakai, \textit{pipeline} RAG, dan abstraksi \textit{agent} di atas LangGraph & \textit{Stateful graph}, memori, \textit{persistence}, \textit{human-in-the-loop}, dan \textit{streaming} & Memaparkan \textit{chain} atau \textit{agent} sebagai \textit{endpoint} FastAPI & \textit{Tracing}, pemantauan, manajemen \textit{prompt}, evaluasi, dan \textit{deployment} melalui LangSmith Deployment \\
    Kapan digunakan & Pengembangan aplikasi LLM secara cepat tanpa perlu mendefinisikan logika graf secara manual & Dibutuhkan kontrol penuh atas alur \textit{agent}, \textit{workflow} multi-agen yang kompleks, atau logika produksi yang sangat spesifik & Hanya untuk pemeliharaan aplikasi lama & Keperluan \textit{debug}, pemantauan, dan \textit{deployment} \\
    \textit{Open-source} & Ya & Ya & Ya & Tidak \\
    Gratis & Ya & Ya & Ya & Sebagian
\end{longtblr}
}

\subsubsection{LangChain 1.0}

LangChain dan LangGraph 1.0 dirilis bersamaan sebagai pembaruan arsitektur yang signifikan dalam ekosistem ini \parencite{langchain2024alpha}. LangGraph 1.0 dipromosikan ke status stabil tanpa perubahan yang memecah kompatibilitas sebelumnya, dengan menyediakan \textit{runtime} agen bawaan yang mencakup eksekusi permanen (\textit{durable execution}), memori jangka pendek, pola \textit{human-in-the-loop}, dan \textit{streaming}. Sementara itu, LangChain 1.0 menyederhanakan antarmukanya dengan berfokus pada satu abstraksi agen inti, yaitu memberikan LLM akses ke alat, memanggil model dengan masukan tertentu, dan berulang kali mengeksekusi alat yang dipilih hingga tidak ada alat yang dipanggil lagi. Abstraksi ini diimplementasikan melalui fungsi \textit{create\_agent} yang dibangun di atas LangGraph, sehingga LangChain 1.0 mewarisi seluruh kapabilitas \textit{runtime}-nya secara langsung. Dengan demikian, LangGraph berperan sebagai \textit{engine} eksekusi yang fleksibel untuk kustomisasi tinggi, sementara LangChain berfungsi sebagai lapisan abstraksi tingkat tinggi yang mempercepat pengembangan. Rantai dan agen lama dari versi sebelumnya tetap tersedia melalui paket \textit{langchain-legacy} agar kompatibilitas mundur terjaga.

\subsubsection{LangSmith Deployment}
LangSmith Deployment, yang sebelumnya dikenal sebagai LangGraph Platform, merupakan infrastruktur khusus yang dirancang untuk \textit{deployment} agen AI pada skala besar \parencite{langchain2024langgraph}. Platform ini menyediakan solusi terpadu untuk mengelola siklus hidup aplikasi agen, mencakup LangGraph Server untuk eksekusi, LangGraph Studio untuk visualisasi dan penelusuran kesalahan (\textit{debugging}), serta SDK terkait. Perubahan nama ini mencerminkan integrasi yang lebih erat dengan ekosistem LangSmith, menawarkan opsi \textit{deployment} yang fleksibel mulai dari Self-Hosted Lite untuk penggunaan lokal hingga Cloud SaaS yang dikelola penuh. Selain itu, tersedia pula opsi Bring Your Own Cloud (BYOC) dan Self-Hosted Enterprise untuk memenuhi kebutuhan keamanan data dan infrastruktur yang spesifik. Infrastruktur ini dirancang untuk menangani beban kerja agen yang berstatus (\textit{stateful}) dan berjalan dalam waktu lama (\textit{long-running}), yang sering kali sulit ditangani oleh infrastruktur web konvensional.

Fitur utama dari LangSmith Deployment mencakup kemampuan skalabilitas horizontal dan antrean tugas untuk menangani lonjakan permintaan secara efisien \parencite{langchain2024langgraph}. Platform ini mendukung persistensi data lintas percakapan, yang memungkinkan pelacakan keadaan (\textit{state tracking}) interaktif serta mekanisme \textit{human-in-the-loop} untuk kontrol agen yang lebih baik. Dukungan terhadap \textit{streaming} respons secara \textit{real-time} dan pemrosesan latar belakang (\textit{background runs}) memberikan fleksibilitas dalam merancang pengalaman pengguna yang responsif maupun tugas pemrosesan \textit{batch}. Lebih lanjut, fitur kontrol konkurensi memastikan integritas proses ketika menangani banyak pesan pengguna secara bersamaan sebelum agen merespons. Keseluruhan fitur ini bertujuan untuk menyederhanakan kompleksitas infrastruktur yang biasanya diperlukan dalam membangun aplikasi agen yang andal dan siap produksi. Gambaran mengenai arsitektur LangSmith Deployment ini diperlihatkan pada Gambar \ref{fig:langgraph_platform}.

\begin{figure}[H]
\centering
\includegraphics[width=0.8\textwidth]{bab2/langgraph_platform.png}
\caption[Arsitektur LangSmith Deployment]{Arsitektur LangSmith Deployment (sebelumnya LangGraph Platform) \parencite{langchainaiLangSmithSDK2025}}
\label{fig:langgraph_platform}
\end{figure}



\section{\textit{LLM-as-a-judge}}

\textit{LLM-as-a-judge} adalah pendekatan evaluasi yang memanfaatkan LLM untuk menilai keluaran LLM lainnya, sebagai respons atas tantangan dalam mengevaluasi respons yang bersifat \textit{open-ended} \parencite{huyenAIEngineeringBuilding2025}. Tantangan ini sebelumnya banyak diselesaikan dengan evaluasi manusia, yang tidak skalabel untuk kebutuhan produksi. Model yang berperan sebagai penilai disebut oleh \textcite{huyenAIEngineeringBuilding2025} sebagai \textit{AI judge}, sementara \textcite{oshinLearningLangchainBuilding2025} menyebutnya \textit{LLM-as-a-judge evaluators} atau sederhananya evaluator LLM. Secara umum, evaluator ini bekerja dengan mengintegrasikan aturan penilaian ke dalam \textit{prompt} LLM, yang kemudian menilai keluaran berdasarkan jawaban referensi yang tersedia \parencite{oshinLearningLangchainBuilding2025}. Dalam penerapannya, skor yang dihasilkan perlu diaudit dan evaluator perlu disetel secara iteratif agar penilaian yang dihasilkan dapat diandalkan.

Keunggulan utama dari metode ini terletak pada efisiensi biaya dan waktu dibandingkan dengan evaluasi manusia, serta kemampuannya untuk memberikan penilaian yang lebih bernuansa dibandingkan metrik heuristik. Sejumlah penelitian menunjukkan bahwa evaluator LLM memiliki korelasi yang tinggi dengan penilaian manusia. \textcite{zhengJudgingLLMasaJudgeMTBench2023} menemukan bahwa pada \textit{benchmark} MT-Bench, tingkat kesepakatan antara GPT-4 \parencite{openaiGPT4TechnicalReport2024} dan manusia mencapai 85\%, lebih tinggi dari kesepakatan antarmanusia yang sebesar 81\%. \textcite{duboisLengthControlledAlpacaEvalSimple2025} juga melaporkan bahwa evaluator LLM pada AlpacaEval memiliki korelasi hampir sempurna (0,98) dengan papan peringkat LMSYS Chat Arena yang dievaluasi secara manusia. 

Selain itu, evaluator LLM dapat menjelaskan alasan di balik penilaiannya, yang memudahkan proses audit hasil evaluasi \parencite{huyenAIEngineeringBuilding2025}. Fleksibilitas ini menjadikan \textit{LLM-as-a-judge} dapat diterapkan pada berbagai jenis aplikasi, bahkan menjadi satu-satunya opsi evaluasi otomatis untuk beberapa kasus penggunaan tertentu. Akan tetapi, \textcite{wangLargeLanguageModels2024} mengidentifikasi sejumlah tantangan, seperti bias posisi (\textit{positional bias}) di mana model cenderung memfavoritkan jawaban yang muncul pertama, serta bias kebertele-telean (\textit{verbosity bias}) yang memberikan skor lebih tinggi pada jawaban yang lebih panjang meskipun tidak lebih informatif.

Sebagai upaya untuk meningkatkan keselarasan dengan penilaian manusia, \textcite{liuGEvalNLGEvaluation2023} memperkenalkan kerangka kerja G-Eval. Pendekatan ini memanfaatkan teknik \textit{Chain-of-Thought} untuk mengarahkan LLM agar menghasilkan langkah-langkah penalaran evaluasi secara eksplisit sebelum memberikan skor akhir. Melalui proses penalaran bertahap ini, model dapat melakukan penilaian yang lebih terstruktur dan transparan, yang pada gilirannya meningkatkan akurasi evaluasi. Selain itu, adanya jejak penalaran memungkinkan verifikasi terhadap logika di balik setiap skor yang diberikan. Mekanisme kerja dari G-Eval ini diilustrasikan pada Gambar \ref{fig:geval_framework}.

\begin{figure}[H]
\centering
\includegraphics[width=0.65\textwidth]{bab2/kerangkakerjageval.png}
\caption[Kerangka kerja G-Eval]{Kerangka kerja G-Eval yang memanfaatkan \textit{Chain-of-Thought} untuk evaluasi \parencite{liuGEvalNLGEvaluation2023}}
\label{fig:geval_framework}
\end{figure}

Berdasarkan eksperimen yang dilakukan, G-Eval menunjukkan korelasi Spearman yang lebih tinggi dengan penilaian manusia dibandingkan dengan metode evaluasi otomatis sebelumnya, membuktikan bahwa penambahan langkah penalaran dapat meningkatkan reliabilitas LLM sebagai juri yang objektif. Hasil eksperimen pada \textit{benchmark} SummEval memperlihatkan bahwa G-Eval-4 mencatatkan kinerja unggul dengan rata-rata korelasi Spearman ($\rho$) sebesar 0,514 dan Kendall-Tau ($\tau$) sebesar 0,418. Capaian ini melampaui metrik konvensional seperti ROUGE-L yang hanya memperoleh rata-rata korelasi Spearman 0,165 dan Kendall-Tau 0,128, serta metrik berbasis model lainnya seperti BERTScore ($\rho=0,225$; $\tau=0,175$) dan UniEval ($\rho=0,474$; $\tau=0,377$). Dominasi G-Eval dalam metrik korelasi ini menegaskan bahwa pendekatan evaluasi berbasis LLM dengan penalaran mampu mendekati penilaian manusia secara lebih akurat dibandingkan metode-metode sebelumnya.

Walaupun gagasan penggunaan AI untuk mengotomasi evaluasi sudah ada sejak lama, pendekatan ini baru menjadi praktis seiring kemampuan model AI yang semakin memadai, yaitu sekitar saat GPT-3 dirilis oleh \textcite{brownLanguageModelsAre2020} \parencite{huyenAIEngineeringBuilding2025}. Sejak saat itu, \textit{LLM-as-a-judge} berkembang menjadi salah satu metode evaluasi yang paling umum digunakan dalam sistem AI produksi. \textcite{langchain_state_of_ai_2023} mencatat bahwa 58\% evaluasi pada platform LangChain dilakukan oleh \textit{AI judge}, mencerminkan adopsi yang luas di industri. Di samping itu, metode ini juga terus berkembang sebagai area penelitian yang aktif.

\subsection{Ragas (\textit{Retrieval Augmented Generation Assessment})}
Ragas (\textit{Retrieval Augmented Generation Assessment}) adalah kerangka kerja evaluasi otomatis yang dikembangkan oleh \textcite{esRagasAutomatedEvaluation2025} untuk mengukur kinerja sistem RAG, dan dapat diakses melalui \textit{library} Python dengan nama yang sama \parencite{ragas2024}. Kerangka kerja ini mengadopsi paradigma \textit{LLM-as-a-judge} dalam mengukur metrik untuk evaluasinya. Evaluasi dalam Ragas dilakukan secara \textit{reference-free} (tanpa referensi jawaban manusia) untuk sebagian metriknya, yang memungkinkan skalabilitas evaluasi pada dataset yang dinamis. Penilaian dilakukan terhadap tiga komponen utama dalam ekosistem RAG, yaitu pertanyaan (\textit{query}), konteks yang diambil (\textit{retrieved context}), dan jawaban yang dihasilkan (\textit{generated answer}). Ragas juga dirancang agar dapat diintegrasikan dengan berbagai \textit{framework} dan \textit{toolchain} yang berbeda seperti LangChain \parencite{langchainaiLangChain2025}, sehingga evaluasi dapat dilakukan tanpa harus mengganti tumpukan teknologi yang sudah digunakan.

\textcite{ragas2024} menyediakan sejumlah metrik evaluasi yang dikelompokkan berdasarkan tipe aplikasi, mencakup sistem RAG, \textit{workflow} agentik, dan perbandingan bahasa alami \parencite{ragas2024}. Untuk sistem RAG, dua metrik utama yang tersedia adalah \textit{Context Precision} dan \textit{Context Recall}, yang keduanya berfokus pada kualitas komponen \textit{retrieval}. \textit{Context Precision} mengukur sejauh mana potongan konteks yang relevan berada di peringkat atas dari konteks yang diambil, dihitung sebagai rata-rata \textit{precision@k} untuk setiap potongan. \textit{Context Recall} melengkapinya dengan mengukur proporsi klaim dalam jawaban referensi yang dapat diatribusikan ke konteks yang diambil, memastikan tidak ada informasi penting yang terlewat. Bersama-sama, kedua metrik ini mencerminkan kualitas \textit{retriever} dari dua sisi, yaitu ketepatan dan kelengkapan.

Untuk evaluasi \textit{workflow} agentik, \textcite{ragas2024} menyediakan beberapa metrik, di antaranya \textit{Tool Call Accuracy} dan \textit{Tool Call F1}. \textit{Tool Call Accuracy} mengukur keakuratan urutan dan argumen pemanggilan alat oleh agen terhadap referensi yang diharapkan, cocok digunakan saat urutan pemanggilan alat bersifat kritis. \textit{Tool Call F1} menawarkan pendekatan yang lebih longgar berbasis \textit{precision} dan \textit{recall} tanpa mempertimbangkan urutan, sehingga berguna untuk mengukur kedekatan perilaku agen dengan ekspektasi meskipun urutan panggilannya berbeda. Selain itu, tersedia pula metrik perbandingan bahasa alami seperti \textit{Factual Correctness} yang mengukur keakuratan faktual respons terhadap referensi melalui dekomposisi klaim dan \textit{natural language inference}, serta \textit{Semantic Similarity} yang mengukur kemiripan semantik antara respons dan referensi menggunakan \textit{embedding}.

\subsection{DeepEval}
DeepEval merupakan kerangka kerja evaluasi \textit{open-source} yang dirancang khusus untuk menguji dan mengevaluasi sistem berbasis LLM, seperti aplikasi RAG dan sistem agentik \parencite{confidentaiDeepEval2025}. Kerangka kerja ini menawarkan pendekatan yang modular dan mudah diintegrasikan dalam siklus pengembangan perangkat lunak serta mengadopsi prinsip pengujian unit (\textit{unit testing}) untuk memvalidasi luaran LLM secara sistematis. Hampir seluruh metrik bawaan DeepEval menggunakan paradigma \textit{LLM-as-a-judge} untuk menilai \textit{test case}, yakni satuan interaksi atomik pengguna dengan aplikasi LLM. DeepEval menyediakan berbagai metrik evaluasi yang sebagian besar memiliki dukungan riset dan seluruhnya bersifat \textit{multimodal} \parencite{confidentaiDeepEval2025}. Hal ini berarti metrik-metrik tersebut mendukung evaluasi sistem berbasis LLM yang memproses lebih dari sekadar teks, di mana \textit{test case} dapat memuat teks maupun citra sekaligus.

Salah satu metrik yang didukung DeepEval adalah implementasi G-Eval dari \textcite{liuGEvalNLGEvaluation2023}, yang menggunakan \textit{CoT prompting} untuk menstabilkan komputasi skor dan membuat penilaian LLM judge lebih andal dan konsisten. Walaupun efektif untuk tugas yang memerlukan penilaian subjektif seperti koherensi dan relevansi jawaban, menurut \textcite{confidentaiDeepEval2025}, G-Eval masih bersifat non-deterministik, sehingga hasil evaluasi dapat bervariasi antar eksekusi pada masukan yang sama. Keterbatasan ini menjadi kendala saat kriteria evaluasi bersifat eksplisit dan terstruktur, seperti pada kebutuhan verifikasi format luaran. 

Untuk mengatasi hal tersebut, DeepEval juga mengembangkan metrik berbasis \textit{Directed Acyclic Graph} (DAG) sebagai alternatif yang lebih deterministik \parencite{confidentaiDeepEval2025}. Dalam pendekatan ini, setiap simpul merepresentasikan LLM judge yang menangani satu keputusan spesifik, sedangkan sisi-sisinya mendefinisikan alur logis antar keputusan. Dengan memecah proses evaluasi menjadi unit-unit atomik yang lebih kecil, ambiguitas diminimalkan sehingga keselarasan dengan ekspektasi dapat ditegakkan secara lebih ketat, seperti yang diilustrasikan pada Gambar~\ref{fig:dag_contoh}.

\begin{figure}[H]
\centering
\includegraphics[width=0.75\textwidth]{bab2/contoh_dag.png}
\caption[Contoh arsitektur metrik DAG pada DeepEval]{Contoh arsitektur metrik DAG pada DeepEval yang mengevaluasi \textit{actual output} melalui serangkaian keputusan atomik yang terstruktur \parencite{confidentaiDeepEval2025}}
\label{fig:dag_contoh}
\end{figure}

\section{Penelitian Terkait}
Pada penelitian terkait, berbagai pendekatan sistem agentik dengan beberapa ragam augmentasi LLM dan integrasi model \textit{deep learning} hingga \textit{computer vision} telah digunakan untuk tugas yang berkaitan dengan identifikasi penyakit tanaman. Pada implementasi paling sederhana, \textcite{sbDesignImplementationAIPowered2024} mengembangkan sistem agentik untuk deteksi penyakit tanaman sekaligus memberikan rekomendasi penanganan. Sistem ini memanfaatkan YOLOv5 untuk mengidentifikasi daun tanaman yang terinfeksi pada berbagai spesies dan jenis penyakit. Selanjutnya, sistem diintegrasikan dengan LLM dari Gemini API guna menyediakan rekomendasi penanganan yang disesuaikan dengan jenis penyakit yang terdeteksi. Evaluasi sistem hanya dilakukan pada model YOLOv5 yakni dengan \textit{mean Average Precision} (mAP) sebesar 0,85 dan \textit{precision} sebesar 0,88 dalam diagnosis penyakit. Pendekatan dari sistem agentik ini dapat dikatakan sebagai \textit{workflow} tanpa ada teknik augmentasi dari LLM, sehingga beresiko terjadi halusinasi.

Mengembangkan dari pendekatan sebelumnya, \textcite{mondalVisionMeetsLanguage2025} memperkenalkan sistem yang lebih maju dengan menggabungkan YOLO dan RAG untuk augmentasi LLM. Sistem ini memanfaatkan YOLOv8 untuk terlebih dahulu mendeteksi penyakit pada daun kopi, kemudian mengintegrasikan mekanisme RAG untuk memberikan hasil diagnosis yang relevan. Pengujian dan evaluasi hanya dilakukan pada model YOLOv8 yang mendapatkan mAP@0.5 sebesar 0.681 dalam deteksi. Sistem agentik ini sendiri merupakan pendekatan \textit{workflow} dengan teknik augmentasi tambahan untuk LLM melalui RAG.

Alternatif pendekatan lainnya terdapat pada penelitian \textcite{roumeliotisPlantDiseaseDetection2025} yang mengeksplorasi penggunaan GPT-4o sebagai LLM \textit{multimodal} dan digabungkan dengan ResNet-50 untuk klasifikasi penyakit tanaman dari citra daun. \citeauthor{roumeliotisPlantDiseaseDetection2025} melakukan evaluasi dalam \textit{setting zero-shot}, \textit{few-shot}, dan \textit{progressive fine-tuning} menggunakan \textit{dataset} PlantVillage, pada dua spesies tanaman, yaitu apel dan jagung, dengan gambar resolusi 100, 150, dan 256 piksel. Hasilnya menunjukkan bahwa model GPT-4o yang di-\textit{fine-tune} mencapai akurasi tertinggi sekitar 98,12\% pada citra daun apel 256 px, sedikit lebih baik dibanding ResNet-50 di $\sim$96,88\% pada kondisi yang sama. Namun, performa \textit{zero-shot} GPT-4o, menunjukkan bahwa meski model \textit{multimodal} LLM kuat, pelatihan tetap diperlukan. Selain itu, meskipun model menunjukkan adaptasi yang baik terhadap berbagai resolusi gambar, generalisasi ke spesies tanaman yang berbeda tetap memiliki batasan, khususnya ketika struktur daun cukup berbeda. Secara keseluruhan, penelitian ini menunjukkan penerapan \textit{workflow} dengan augmentasi \textit{fine-tuning} pada LLM.

Pendekatan yang menarik terdapat pada penelitian \textcite{xuMultimodalAgriculturalAgent2025} yang mengusulkan MA3 sebagai pendekatan \textit{workflow} yang dapat mengambil keputusan dalam identifikasi penyakit tanaman secara adaptif. Arsitektur ini dirancang agar lebih adaptif melalui modul \textit{router} berisi model BERT yang dilatih secara \textit{supervised learning} untuk mengklasifikasikan \textit{query} (\textit{input} teks) pengguna. Berdasarkan hasil klasifikasi tersebut, sistem dapat mengarahkan proses ke modul klasifikasi penyakit, deteksi penyakit, atau tanya jawab. Modul klasifikasi dan deteksi dibangun menggunakan \textit{backbone} CLIP-ViT yang di-\textit{fine-tune} untuk data penyakit tebu, sementara komponen terakhir memanfaatkan Qwen2.5-32B sebagai LLM guna menggabungkan keluaran seluruh alat dan menyusun jawaban akhir yang lebih informatif. Integrasi komponen-komponen ini menjadikan MA3 sebagai sistem yang \textit{multimodal}, karena sistem mampu memproses masukan berupa teks sekaligus citra, di mana teks diproses oleh \textit{router} dan LLM, sementara citra diproses oleh modul berbasis CLIP-ViT untuk klasifikasi maupun deteksi penyakit.

Dengan desain MA3, \textcite{xuMultimodalAgriculturalAgent2025} menghadirkan pendekatan \textit{workflow} adaptif yang membedakannya dari pendekatan agen, karena pemilihan alat dikendalikan oleh \textit{router}, bukan oleh LLM itu sendiri. Pilihan penggunaan \textit{router} dalam penggunaan alat yang adaptif dimaksudkan untuk menghindari halusinasi dari LLM. Akan tetapi \citeauthor{xuMultimodalAgriculturalAgent2025} luput menyebutkan pendekatan seperti agen dengan pola ReAct yang sebenarnya sudah menawarkan solusi atas masalah tersebut. Di mana, pola ReAct dapat diterapkan langsung ke LLM untuk memberikan kemampuan pengambilan keputusan yang adaptif dalam memilih alat, tanpa perlu menambahkan komponen \textit{router} terpisah. 

Hasil evaluasi pada sistem yang dikembangkan penelitian \textcite{xuMultimodalAgriculturalAgent2025} menunjukkan bahwa \textit{router} BERT mencapai akurasi sangat tinggi di atas 99\%, modul klasifikasi penyakit memperoleh \textit{precision} hingga 96,2\%, sedangkan modul deteksi penyakit masih menghadapi tantangan dengan mAP sekitar 0,325. Evaluasi dilakukan dengan kombinasi metrik tradisional serta metrik kualitas tanya jawab seperti \textit{semantic consistency}, \textit{information completeness}, \textit{information leakage}, dan \textit{redundancy} oleh \textit{evaluator} pakar di bidang pertanian secara \textit{offline}. Pada sistem agentik yang dikembangkan, \textcite{xuMultimodalAgriculturalAgent2025} juga memperkenalkan metode evaluasi \textit{online} menggunakan LLM berupa DeepSeek-V3 sebagai evaluator LLM untuk memeriksa kualitas dan relevansi dari penggunaan berbagai komponen MA3, baik dalam tugas klasifikasi, deteksi, maupun tanya jawab.

Sejauh ini belum ada penelitian yang menerapkan pendekatan agen dalam sistem agentik untuk identifikasi penyakit tanaman melalui pola ReAct, maupun mengintegrasikannya dengan VLM dan dapat mengevaluasi sistemnya secara komprehensif. Pendekatan yang paling mendekati ditunjukkan oleh ShizishanGPT yang diperkenalkan oleh \textcite{yangShizishanGPTAgriculturalLarge2024}, yakni sebuah sistem agentik yang dirancang untuk mendukung berbagai tugas agrikultur, termasuk seputar penyakit tanaman, meski penggunaannya masih terbatas pada skenario tanya jawab tanpa analisis citra. Dalam sistem, terdapat agen LLM dengan pola ReAct yang menggunakan GPT-4 dengan akses modul yang berisi alat eksternal seperti pencarian internet, prediksi fenotipe tanaman, analisis gen, dan model cuaca. ShizishanGPT juga memiliki modul pengambilan informasi berbasis RAG, tetapi modul ini tidak terintegrasi pada agen sebagai alat, sehingga penggunaannya tidak adaptif. 

Untuk eksperimen, \textcite{yangShizishanGPTAgriculturalLarge2024} menyiapkan \textit{dataset} sintetis berisi 100 pasangan pertanyaan dan jawaban pertanian yang digenerasi oleh LLM. Evaluasi dilakukan melalui \textit{machine scoring} pada modul pengambilan informasi berbasis RAG menggunakan beberapa metrik, di mana ShizishanGPT meraih skor komposit tertinggi di 0,923. Selain itu, \textit{manual scoring} oleh \textit{evaluator} untuk akurasi, profesionalisme, dan kemampuan bahasa dari jawaban menempatkan ShizishanGPT pada skor 9,65 dari 10, lebih unggul dibanding GPT-4 maupun model lainnya. Evaluasi juga dilakukan melalui eksperimen yang menghapus beberapa modul dari sistem untuk menilai kontribusi masing-masing komponen. Hasilnya menunjukkan konfigurasi penuh memberikan hasil terbaik, sementara penghapusan modul pengambilan informasi berbasis RAG memberikan skor terendah di 9,21, menunjukkan peran penting RAG dalam menjaga kualitas jawaban. 

Secara keseluruhan, penelitian-penelitian sebelumnya menunjukkan berbagai pendekatan dalam penerapan sistem agentik untuk identifikasi penyakit tanaman. Penerapan ini dimulai dari integrasi \textit{workflow} sederhana dengan YOLO dan LLM, pengambilan informasi melalui RAG, \textit{fine-tuning} pada LLM, hingga arsitektur \textit{workflow} yang kompleks dengan \textit{router} dan model-model khusus. Peningkatan signifikan terlihat pada \textcite{xuMultimodalAgriculturalAgent2025} dengan munculnya sistem agentik yang lebih otonom dan reliabel. Akan tetapi, belum ada yang secara komprehensif mengintegrasikan pola ReAct dengan VLM untuk identifikasi penyakit tanaman secara \textit{end-to-end} dengan evaluasi yang menyeluruh. Rangkuman dari penelitian-penelitian tersebut ditunjukkan pada Tabel \ref{tab:summary_related_research}. Rangkuman ini disusun dengan menyoroti aspek pendekatan sistem agentik yang digunakan, augmentasi LLM, integrasi model atau teknologi tambahan, serta hasil evaluasi utama yang dilaporkan dalam setiap penelitian.

{
\setstretch{1.0}
\begin{longtblr}[
    caption = {Rangkuman penelitian terkait sistem agentik untuk identifikasi penyakit tanaman},
    entry = {Rangkuman penelitian terkait},
    label = {tab:summary_related_research},
]{
    width = \linewidth,
    colspec = {X[1.5,l] X[2,l] X[2,l] X[2,l]},
    rows = {m},
    rowhead = 1,
    row{1} = {font=\bfseries, halign=c, valign=m},
    hlines, vlines,
    rowsep = 3pt,
}
    Peneliti dan Tahun & Pendekatan Sistem Agentik & Integrasi Tambahan & Evaluasi \\
    \textcite{sbDesignImplementationAIPowered2024} & \textit{Workflow} sederhana & YOLOv5 untuk deteksi penyakit tanaman & mAP 0,85 dan \textit{precision} 0,88 pada model YOLOv5 \\
    \textcite{roumeliotisPlantDiseaseDetection2025} & \textit{Workflow} sederhana dengan \textit{fine-tuning} LLM & ResNet-50 untuk klasifikasi tanaman & Akurasi hingga 98,12\% pada GPT-4o (\textit{fine-tuned}) dan $\sim$96,88\% pada ResNet-50 \\
    \textcite{xuMultimodalAgriculturalAgent2025} & \textit{Workflow} adaptif dengan \textit{fine-tuning} LLM & BERT sebagai \textit{router} penggunaan alat adaptif dan CLIP-ViT untuk deteksi lesi serta klasifikasi penyakit tanaman & \textit{Router} akurasi $>$99\%; klasifikasi \textit{precision} 96,2\%; deteksi mAP $\sim$0,325; Evaluasi komponen menggunakan \textit{LLM-as-a-judge} \\
    \textcite{yangShizishanGPTAgriculturalLarge2024} & Kombinasi \textit{workflow} dengan ReAct & RAG, model--model untuk berbagai tugas, dan pencarian internet & Skor komposit RAG 0,923 (\textit{machine scoring}) dan skor manual 9,65/10 (\textit{human scoring}) \\
\end{longtblr}
}

